% This paper has been transcribed in Plain TeX by
% David R. Wilkins
% School of Mathematics, Trinity College, Dublin 2, Ireland
% (dwilkins@maths.tcd.ie)
%
% Trinity College, 2000.

\magnification=\magstep1
\vsize=227 true mm \hsize=170 true mm
   \voffset=-0.4 true mm \hoffset=-5.4 true mm

\def\folio{\ifnum\pageno>0 \number\pageno \else
   \ifnum\pageno<0 \romannumeral-\pageno \else\fi\fi}

\font\Largebf=cmbx10  scaled \magstep2
\font\largerm=cmr12
\font\largeit=cmti12
\font\largesc=cmcsc10  scaled \magstep1
\font\sc=cmcsc10

\def\multieqalign#1{\null\,\vcenter{\openup1\jot \mathsurround=0pt
   \ialign{\strut\hfil$\displaystyle{##}$&$\displaystyle{{}##}$\hfil
   &&\quad \strut\hfil$\displaystyle{##}$&$\displaystyle{{}##}$\hfil
    \crcr#1\crcr}}\,}

\def\sum{\mathop{\Sigma}\nolimits}

\def\sectiontitle{\it\noindent}

\pageno=0

\null\vskip72pt

\centerline{\Largebf SECOND ESSAY ON A}

\vskip12pt

\centerline{\Largebf GENERAL METHOD IN DYNAMICS}

\vskip24pt

\centerline{\Largebf By}

\vskip24pt

\centerline{\Largebf William Rowan Hamilton}

\vskip24pt

\centerline{\largerm (Philosophical Transactions of the Royal Society,
   part~I for 1835, pp.\ 95--144.)}

\vskip36pt

\vfill

\centerline{\largerm Edited by David R. Wilkins}

\vskip 12pt

\centerline{\largerm 2000}

\vskip36pt\eject

\pageno=-1

\null\vskip36pt

\centerline{\Largebf NOTE ON THE TEXT}

\bigskip

This edition is based on the original publication in the {\it
Philosophical Transactions of the Royal Society}, part~I for
1835.

The following errors in the original published text have been
corrected:

\smallskip

\item{}
a term $dt$, missing from the integral (R${}^1$.) has been inserted;

\smallskip

\item{}
a factor $t^2$ has been inserted in the
integrand in the second integral of equation (160.).

\smallskip

The {\it Second Essay on a General Method in Dynamics} has also
been republished in {\it The Mathematical Papers of Sir
William Rowan Hamilton, Volume II: Dynamics}, edited for
the Royal Irish Academy by A.~W. Conway and A.~J. McConnell, and
published by Cambridge University Press in 1940.

\bigbreak\bigskip

\line{\hfil David R. Wilkins}

\vskip3pt

\line{\hfil Dublin, February 2000}


\vskip6pt

\line{\hfil Corrected March 2005}

\vfill\eject

\pageno=1

\null\vskip36pt

\noindent
{\largeit
Second Essay on a General Method in Dynamics.
By {\largesc William Rowan Hamilton}, Member of several Scientific
Societies in Great Britain and in Foreign Countries,
Andrews' Professor of Astronomy in the University of Dublin,
and Royal Astronomer of Ireland.  Communicated by
Captain {\largesc Beaufort}, R.N. F.R.S.}

\vskip12pt

\centerline{Received October 29, 1834,---Read January 15, 1835.}
\vskip12pt
\centerline{[{\it Philosophical Transactions of the Royal Society},
   part~I for 1835, pp.\ 95--144.]}

\bigbreak

{\sectiontitle
Introductory Remarks.\par}

\nobreak\bigskip

The former Essay\footnote*{Philosophical Transactions for the
year 1834, Second Part.}
contained a general method for reducing all the
most important problems of dynamics to the study of one
characteristic function, one central or radical relation.  It was
remarked at the close of that Essay, that many eliminations
required by this method in its first conception, might be avoided
by a general transformation, introducing the time explicitly into
a part~$S$ of the whole characteristic function~$V$; and it is
now proposed to fix the attention chiefly on this part~$S$,
and to call it the {\it Principal Function}.  The properties
of this part or function~$S$, which were noticed briefly in the
former Essay, are now more fully set forth; and especially its
uses in questions of perturbation, in which it dispenses with
many laborious and circuitous processes, and enables us to
express accurately the disturbed configuration of a system by the
rules of undisturbed motion, if only the initial components of
velocities be changed in a suitable manner.  Another manner of
extending rigorously to disturbed motion the rules of
undisturbed, by the gradual variation of elements, in number
double the number of coordinates or other marks of position of
the system, which was first invented by {\sc Lagrange}, and was
afterwards improved by {\sc Poisson}, is considered in this Second
Essay under a form perhaps a little more general; and the general
method of calculation which has already been applied to other
analogous questions in optics and in dynamics by the author of the
present Essay, is now applied to the integration of the equations
which determine these elements.  This general method is founded
chiefly on a combination of the principles of variations with
those of partial differentials, and may furnish, when it shall be
matured by the labours of other analysts, a separate branch of
algebra, which may be called perhaps the {\it Calculus of
Principal Functions\/}; because, in all the chief applications of
algebra to physics, and in a very extensive class of purely
mathematical questions, it reduces the determination of many
mutually connected functions to the search and study of one
principal or central relation.  When applied to the integration
of the equations of varying elements, it suggests, as is now
shown, the consideration of a certain {\it Function of Elements},
which may be variously chosen, and may either be rigorously
determined, or at least approached to, with an indefinite
accuracy, by a corollary of the general method.  And to
illustrate all these new general processes, but especially those
which are connected with problems of perturbation, they are
applied in this Essay to a very simple example, suggested by the
motions of projectiles, the parabolic path being treated as the
undisturbed.  As a more important example, the problem of
determining the motions of a ternary or multiple system, with any
laws of attraction or repulsion, and with one predominant mass,
which was touched on in the former Essay, is here resumed in a
new way, by forming and integrating the differential equations of
a new set of varying elements, entirely distinct in theory
(though little differing in practice) from the elements conceived
by {\sc Lagrange}, and having this advantage, that the differentials of
all the new elements for {\it both\/} the disturbed and
disturbing masses may be expressed by the coefficients of
{\it one\/} disturbing function.

\vfill\eject

{\sectiontitle
Transformations of the Differential Equations of Motion of an
Attracting or Repelling System.\par}

\nobreak\bigskip

1.
It is well known to mathematicians, that the differential
equations of motion of any system of free points, attracting or
repelling one another according to any functions of their distances,
and not disturbed by any foreign force, may be comprised in the
following formula:
$$\sum \mathbin{.} m (x'' \delta x + y'' \delta y + z'' \delta z)
   = \delta U:
   \eqno {\rm (1.)}$$
the sign of summation $\sum$ extending to all the points of the
system; $m$ being, for any one such point, the constant called its
mass, and $x$~$y$~$z$ being its rectangular coordinates; while
$x''$~$y''$~$z''$ are the accelerations, or second differential
coefficients taken with respect to the time, and $\delta x$,
$\delta y$, $\delta z$ are any arbitrary infinitesimal variations
of those coordinates, and $U$ is a certain {\it force-function},
introduced into dynamics by {\sc Lagrange}, and involving the masses and
mutual distances of the several points of the system.  If the number
of those points be $n$, the formula (1.) may be decomposed into $3n$
ordinary differential equations of the second order, between the
coordinates and the time,
$$m_i x''_i = {\delta U \over \delta x_i};\quad
  m_i y''_i = {\delta U \over \delta y_i};\quad
  m_i z''_i = {\delta U \over \delta z_i}:
   \eqno {\rm (2.)}$$
and to integrate these differential equations of motion of an
attracting or repelling system, or some transformations of these, is
the chief and perhaps ultimately the only problem of mathematical
dynamics.

\medskip

2.
To facilitate and generalize the solution of this problem, it is
useful to express previously the $3n$ rectangular coordinates
$x$~$y$~$z$ as functions of $3n$ other and more general marks of
position $\eta_1 \, \eta_2 \,\ldots\, \eta_{3n}$; and then the
differential equations of motion take this more general form,
discovered by {\sc Lagrange},
$${d \over dt} {\delta T \over \delta \eta_i'}
   - {\delta T \over \delta \eta_i}
   = {\delta U \over \delta \eta_i},
   \eqno {\rm (3.)}$$
in which
$$T = {\textstyle {1 \over 2}} \sum \mathbin{.} m
         (x'^2 + y'^2 + z'^2).
   \eqno {\rm (4.)}$$

For, from the equations (2.) or (1.),
$$\left. \eqalign{
{\delta U \over \delta \eta_i}
   &= \sum \mathbin{.} m \left(
           x'' {\delta x \over \delta \eta_i}
         + y'' {\delta y \over \delta \eta_i}
         + z'' {\delta z \over \delta \eta_i}
         \right) \cr
   &= {d \over dt} \sum \mathbin{.} m \left(
           x' {\delta x \over \delta \eta_i}
         + y' {\delta y \over \delta \eta_i}
         + z' {\delta z \over \delta \eta_i}
         \right) \cr
   &\mathrel{\phantom{=}} \mathord{}
       - \sum \mathbin{.} m \left(
           x' {d \over dt} {\delta x \over \delta \eta_i}
         + y' {d \over dt} {\delta y \over \delta \eta_i}
         + z' {d \over dt} {\delta z \over \delta \eta_i}
         \right);\cr}
   \right\}
   \eqno {\rm (5.)}$$
in which
$$\left. \eqalign{
   &\mathrel{\phantom{=}} \mathord{}
      \sum \mathbin{.} m \left(
           x' {\delta x \over \delta \eta_i}
         + y' {\delta y \over \delta \eta_i}
         + z' {\delta z \over \delta \eta_i}
         \right) \cr
   &= \sum \mathbin{.} m \left(
           x' {\delta x' \over \delta \eta_i'}
         + y' {\delta y' \over \delta \eta_i'}
         + z' {\delta z' \over \delta \eta_i'}
         \right)
      = {\delta T \over \delta \eta_i'},\cr}
   \right\}
   \eqno {\rm (6.)}$$
and
$$\left. \eqalign{
   &\mathrel{\phantom{=}} \mathord{}
      \sum \mathbin{.} m \left(
           x' {d \over dt} {\delta x \over \delta \eta_i}
         + y' {d \over dt} {\delta y \over \delta \eta_i}
         + z' {d \over dt} {\delta z \over \delta \eta_i}
         \right) \cr
   &= \sum \mathbin{.} m \left(
           x' {\delta x' \over \delta \eta_i}
         + y' {\delta y' \over \delta \eta_i}
         + z' {\delta z' \over \delta \eta_i}
         \right)
      = {\delta T \over \delta \eta_i},\cr}
   \right\}
   \eqno {\rm (7.)}$$
$T$ being here considered as a function of the $6n$ quantities
of the forms $\eta'$ and $\eta$, obtained by introducing into its
definition (4.), the values
$$x' =   \eta_1' {\delta x \over \delta \eta_1}
       + \eta_2' {\delta x \over \delta \eta_2}
       + \cdots
       + \eta_{3n}' {\delta x \over \delta \eta_{3n}},
   \hbox{ \&c.}
   \eqno {\rm (8.)}$$

A different proof of this important transformation (3.) is
given in the M\'{e}canique Analytique.

\medskip

3.
The function $T$, being homogeneous of the second dimension
with respect to the quantities $\eta'$, must satisfy the condition
$$2T = \sum \mathbin{.} \eta' {\delta T \over \delta \eta'};
   \eqno {\rm (9.)}$$
and since the variation of the same function~$T$ may evidently
be expressed as follows,
$$\delta T = \sum \left( {\delta T \over \delta \eta'} \delta \eta'
                         + {\delta T \over \delta \eta} \delta \eta
                         \right),
   \eqno {\rm (10.)}$$
we see that this variation may be expressed in this other way,
$$\delta T = \sum \left( \eta' \delta {\delta T \over \delta \eta'}
                         - {\delta T \over \delta \eta} \delta \eta
                         \right).
   \eqno {\rm (11.)}$$

If then we put, for abridgement,
$${\delta T \over \delta \eta_1'}    = \varpi_1,\quad \ldots \quad
  {\delta T \over \delta \eta_{3n}'} = \varpi_{3n},
   \eqno {\rm (12.)}$$
and consider $T$ (as we may) as a function of the following form,
$$T = F(\varpi_1, \varpi_2,\ldots\, \varpi_{3n},
        \eta_1, \eta_2,\ldots\, \eta_{3n}),
   \eqno {\rm (13.)}$$
we see that
$${\delta F \over \delta \varpi_1} = \eta_1',\quad \ldots \quad
  {\delta F \over \delta \varpi_{3n}} = \eta_{3n}',
   \eqno {\rm (14.)}$$
and
$${\delta F \over \delta \eta_1}
  = - {\delta T \over \delta \eta_1},\quad \ldots \quad
  {\delta F \over \delta \eta_{3n}}
  = - {\delta T \over \delta \eta_{3n}};
   \eqno {\rm (15.)}$$
and therefore that the general equation (3.) may receive this
new transformation,
$${d \varpi_i \over dt} = {\delta (U - F) \over \delta \eta_i}.
   \eqno {\rm (16.)}$$

If then we introduce, for abridgement, the following expression~$H$,
$$H = F - U = F(\varpi_1, \varpi_2,\ldots\, \varpi_{3n},
                \eta_1, \eta_2,\ldots\, \eta_{3n})
              - U(\eta_1, \eta_2,\ldots\, \eta_{3n}),
   \eqno {\rm (17.)}$$
we are conducted to this new manner of presenting the differential
equations of motion of a system of $n$ points, attracting or
repelling one another:
$$\left. \multieqalign{
      {d \eta_1 \over dt} &= {\delta H \over \delta \varpi_1}; &
      {d \varpi_1 \over dt} & = - {\delta H \over \delta \eta_1};\cr
      {d \eta_2 \over dt} &= {\delta H \over \delta \varpi_2}; &
      {d \varpi_2 \over dt} &= - {\delta H \over \delta \eta_2};\cr
   \noalign{\vskip6pt\hbox{.........}\vskip6pt}
      {d \eta_{3n} \over dt} &= {\delta H \over \delta \varpi_{3n}}; &
      {d \varpi_{3n} \over dt} &= - {\delta H \over \delta \eta_{3n}}.\cr}
   \right\}
   \eqno {\rm (A.)}$$
In this view, the problem of mathematical dynamics, for a system of
$n$ points, is to integrate a system (A.) of $6n$ ordinary differential
equations of the first order, between the $6n$ variables
$\eta_i$~$\varpi_i$ and the time~$t$; and the solution of the problem
must consist in assigning these $6n$ variables as functions of the time,
and of their own initial values, which we may call $e_i$~$p_i$.
And all these $6n$ functions, or $6n$ relations to determine them,
may be expressed, with perfect generality and rigour, by the method
of the former Essay, or by the following simplified process.

\bigbreak

{\sectiontitle
Integration of the Equations of Motion, by means of one
Principal Function.\par}

\nobreak\bigskip

4.
If we take the variation of the definite integral
$$S = \int_0^t \left( \sum \mathbin{.}
         \varpi {\delta H \over \delta \varpi} - H \right) dt
   \eqno {\rm (18.)}$$
without varying $t$ or $dt$, we find, by the Calculus of Variations,
$$\delta S = \int_0^t \delta S' \mathbin{.} dt,
   \eqno {\rm (19.)}$$
in which
$$S' = \sum \mathbin{.} \varpi {\delta H \over \delta \varpi} - H,
   \eqno {\rm (20.)}$$
and therefore
$$\delta S'
   = \sum \left( \varpi \, \delta {\delta H \over \delta \varpi}
                 - {\delta H \over \delta \eta} \delta \eta \right),
   \eqno {\rm (21.)}$$
that is, by the equations of motion (A.),
$$\delta S'
   = \sum \left( \varpi \, \delta {d \eta \over dt}
                 + {d \varpi \over dt} \delta \eta \right)
   = {d \over dt} \sum \mathbin{.} \varpi \, \delta \eta;
   \eqno {\rm (22.)}$$
the variation of the integral~$S$ is therefore
$$\delta S = \sum ( \varpi \, \delta \eta - p \, \delta e),
   \eqno {\rm (23.)}$$
($p$ and $e$ being still initial values,) and it decomposes itself
into the following $6n$ expressions, when $S$ is considered as
a function of the $6n$ quantities $\eta_i$~$e_i$, (involving also
the time,)
$$\left. \multieqalign{
      \varpi_1 &= {\delta S \over \delta \eta_1}; &
      p_1 &= - {\delta S \over \delta e_1};\cr
      \varpi_2 &= {\delta S \over \delta \eta_2}; &
      p_2 &= - {\delta S \over \delta e_2};\cr
   \noalign{\vskip6pt\hbox{.........}\vskip6pt}
      \varpi_{3n} &= {\delta S \over \delta \eta_{3n}}; &
      p_{3n} &= - {\delta S \over \delta e_{3n}};\cr}
   \right\}
   \eqno {\rm (B.)}$$
which are evidently forms for the sought integrals of the $6n$
differential equations of motion (A.), containing only one
unknown function~$S$.  The difficulty of mathematical dynamics
is therefore reduced to the search and study of this one function~$S$,
which may for that reason be called the {\sc Principal Function}
of motion of a system.

This function~$S$ was introduced in the first Essay under the form
$$S = \int_0^t (T + U) \, dt,$$
the symbols $T$ and $U$ having in this form their recent meanings;
and it is worth observing, that when $S$ is expressed by this definite
integral, the conditions for its variation vanishing (if the final
and initial coordinates and the time be given) are precisely the
differential equations of motion (3.), under the forms assigned
by {\sc Lagrange}.  The variation of this definite integral~$S$ has
therefore the double property, of giving the differential equations
of motion for any transformed coordinates when the extreme positions
are regarded as fixed, and of giving the integrals of those
differential equations when the extreme positions are treated
as varying.

\medskip

5.
Although the function~$S$ seems to deserve the name here given
it of {\it Principal Function}, as serving to express, in what
appears the simplest way, the integrals of the equations of motion,
and the differential equations themselves; yet the same analysis
conducts to other functions, which also may be used to express the
integrals of the same equations.  Thus if we put
$$Q = \int_0^t \left( - \sum \mathbin{.} \eta
                      {\delta H \over \delta \eta} + H \right) dt,
   \eqno {\rm (24.)}$$
and take the variation of this integral~$Q$ without varying $t$
or $dt$, we find, by a similar process,
$$\delta Q = \sum ( \eta \, \delta \varpi - e \, \delta p);
   \eqno {\rm (25.)}$$
so that if we consider $Q$ as a function of the $6n$ quantities
$\varpi_i$~$p_i$ and of the time, we shall have $6n$ expressions
$$\eta_i = + {\delta Q \over \delta \varpi_i},\quad
  e_i    = - {\delta Q \over \delta p_i},
   \eqno {\rm (26.)}$$
which are other forms of the integrals of the equations of
motion (A.), involving the function~$Q$ instead of $S$.  We might
also employ the integral
$$V = \int_0^t \sum \mathbin{.} \varpi {\delta H \over \delta \varpi} \,dt
    = \sum \int_e^\eta \varpi \,d\eta,
   \eqno {\rm (27.)}$$
which was called the {\it Characteristic Function\/} in the former
Essay, and of which, when considered as a function of the $6n + 1$
quantities $\eta_i$~$e_i$~$H$, the variation is
$$\delta V = \sum (\varpi \, \delta \eta - p \, \delta e)
      + t \, \delta H.
   \eqno {\rm (28.)}$$
And all these functions $S$, $Q$, $V$, are connected in such a way,
that the forms and properties of any one may be deduced from those
of any other.

\bigbreak

{\sectiontitle
Investigation of a Pair of Partial Differential Equations of the
first Order, which the Principal Function must satisfy.\par}

\nobreak\bigskip

6.
In forming the variation (23.), or the partial differential
coefficients (B.) of the Principal Function~$S$, the variation
of the time was omitted; but it is easy to calculate the coefficient
$\displaystyle {\delta S \over \delta t}$ corresponding to this
variation, since the evident equation
$${d S \over dt}
   = {\delta S \over \delta t}
      + \sum {\delta S \over \delta \eta} {d\eta \over dt}
   \eqno {\rm (29.)}$$
gives, by (20.), and by (A.), (B.),
$${\delta S \over \delta t}
   = S' - \sum \mathbin{.} \varpi {\delta H \over \delta \varpi}
   = -H.
   \eqno {\rm (30.)}$$

It is evident also that this coefficient, or the quantity $-H$,
is constant, so as not to alter during the motion of the system;
because the differential equations of motion (A.) give
$${dH \over dt} = \sum \left(
        {\delta H \over \delta \eta}{d\eta \over dt}
      + {\delta H \over \delta \varpi}{d\varpi \over dt} \right)
   = 0.
   \eqno {\rm (31.)}$$
If, therefore, we attend to the equation (17.), and observe that
the function~$F$ is necessarily rational and integer and homogeneous
of the second dimension with respect to the quantities $\varpi_i$,
we shall perceive that the principal function~$S$ must satisfy the
two following equations between its partial differential coefficients
of the first order, which offer the chief means of discovering
its form:
$$\left. \eqalign{
{\delta S \over \delta t}
   + F \left( {\delta S \over \delta \eta_1},
              {\delta S \over \delta \eta_2},\ldots\,
              {\delta S \over \delta \eta_{3n}},
              \eta_1, \eta_2,\ldots\, \eta_{3n}
              \right)
   &= U(\eta_1, \eta_2,\ldots\, \eta_{3n}),\cr
{\delta S \over \delta t}
   + F \left( {\delta S \over \delta e_1},
              {\delta S \over \delta e_2},\ldots\,
              {\delta S \over \delta e_{3n}},
              e_1, e_2,\ldots\, e_{3n}
              \right)
   &= U(e_1, e_2,\ldots\, e_{3n}).\cr}
   \right\}
   \eqno {\rm (C.)}$$

Reciprocally, if the form of $S$ be known, the forms of these
equations (C.) can be deduced from it, by elimination of the quantities
$e$ or $\eta$ between the expressions of its partial differential
coefficients; and thus we can return from the principal function~$S$
to the functions $F$ and $U$, and consequently to the expression~$H$,
and the equations of motion (A.)

Analogous remarks apply to the functions $Q$ and $V$, which must
satisfy the partial differential equations,
$$\left. \eqalign{
-{\delta Q \over \delta t}
    + F \left( \varpi_1, \varpi_2,\ldots, \varpi_{3n},
               {\delta Q \over \delta \varpi_1},
               {\delta Q \over \delta \varpi_2},\ldots\,
               {\delta Q \over \delta \varpi_{3n}}
               \right)
   &= U \left( {\delta Q \over \delta \varpi_1},
               {\delta Q \over \delta \varpi_2},\ldots\,
               {\delta Q \over \delta \varpi_{3n}}
               \right),\cr
-{\delta Q \over \delta t}
    + F \left( p_1, p_2,\ldots, p_{3n},
               -{\delta Q \over \delta p_1},
               -{\delta Q \over \delta p_2},\ldots\,
               -{\delta Q \over \delta p_{3n}}
               \right)
   &= U \left( -{\delta Q \over \delta p_1},
               -{\delta Q \over \delta p_2},\ldots\,
               -{\delta Q \over \delta p_{3n}}
               \right),\cr}
   \right\}
   \eqno {\rm (32.)}$$
and
$$\left. \eqalign{
F \left( {\delta V \over \delta \eta_1},
         {\delta V \over \delta \eta_2},\ldots\,
         {\delta V \over \delta \eta_{3n}},
         \eta_1, \eta_2,\ldots\, \eta_{3n}
              \right)
   &= H + U(\eta_1, \eta_2,\ldots\, \eta_{3n}),\cr
F \left( {\delta V \over \delta e_1},
         {\delta V \over \delta e_2},\ldots\,
         {\delta V \over \delta e_{3n}},
         e_1, e_2,\ldots\, e_{3n}
              \right)
   &= H + U(e_1, e_2,\ldots\, e_{3n}).\cr}
   \right\}
   \eqno {\rm (33.)}$$

\bigbreak

{\sectiontitle
General Method of improving an approximate Expression for the
Principal Function in any Problem of Dynamics.\par}

\nobreak\bigskip

7.
If we separate the principal function~$S$ into any two parts,
$$S_1 +S_2 = S,
   \eqno {\rm (34.)}$$
and substitute their sum for $S$ in the first equation (C.), the
function~$F$, from its rational and integer and homogeneous form
and dimension, may be expressed in this new way,
$$\left. \eqalign{
   & F \left( {\delta S \over \delta \eta_1},\ldots\,
              {\delta S \over \delta \eta_{3n}},
              \eta_1,\ldots\, \eta_{3n} \right)
   = F \left( {\delta S_1 \over \delta \eta_1},\ldots\,
              {\delta S_1 \over \delta \eta_{3n}},
              \eta_1,\ldots\, \eta_{3n} \right) \cr
   &\qquad\mathrel{\phantom{=}} \mathord{}
   + F' \left( {\delta S_1 \over \delta \eta_1} \right)
               {\delta S_2 \over \delta \eta_1} + \cdots
   + F' \left( {\delta S_1 \over \delta \eta_{3n}} \right)
               {\delta S_2 \over \delta \eta_{3n}}
   + F \left( {\delta S_2 \over \delta \eta_1},\ldots\,
              {\delta S_2 \over \delta \eta_{3n}},
              \eta_1,\ldots\, \eta_{3n} \right) \cr
   &\qquad
   = F \left( {\delta S_1 \over \delta \eta_1},\ldots\,
              {\delta S_1 \over \delta \eta_{3n}},
              \eta_1,\ldots\, \eta_{3n} \right)
   - F \left( {\delta S_2 \over \delta \eta_1},\ldots\,
              {\delta S_2 \over \delta \eta_{3n}},
              \eta_1,\ldots\, \eta_{3n} \right) \cr
   &\qquad\mathrel{\phantom{=}} \mathord{}
   + F' \left( {\delta S   \over \delta \eta_1} \right)
               {\delta S_2 \over \delta \eta_1} + \cdots
   + F' \left( {\delta S   \over \delta \eta_{3n}} \right)
               {\delta S_2 \over \delta \eta_{3n}},\cr}
   \right\}
   \eqno {\rm (35.)}$$
because
$$F' \left( {\delta S_1 \over \delta \eta_i} \right)
   =   F' \left( {\delta S   \over \delta \eta_i} \right)
     - F' \left( {\delta S_2 \over \delta \eta_i} \right),
   \eqno {\rm (36.)}$$
and
$$\sum \mathbin{.}  F' \left( {\delta S_2 \over \delta \eta} \right)
                    {\delta S_2 \over \delta \eta}
  = 2F \left( {\delta S_2 \over \delta \eta_1},\ldots\,
              {\delta S_2 \over \delta \eta_{3n}},
              \eta_1,\ldots\, \eta_{3n} \right);
   \eqno {\rm (37.)}$$
and since, by (A.) and (B.),
$$F' \left( {\delta S   \over \delta \eta_i} \right)
   = F'(\varpi_i) = {\delta H \over \delta \varpi_i}
   = {d\eta_i \over dt},
   \eqno {\rm (38.)}$$
we easily transform the first equation (C.) to the following,
$$\left. \eqalign{
{d S_2 \over dt}
  &= - {\delta S_1 \over \delta t}
     + U(\eta_1,\ldots\, \eta_{3n})
     - F \left( {\delta S_1 \over \delta \eta_1},\ldots\,
                {\delta S_1 \over \delta \eta_{3n}},
                \eta_1,\ldots\, \eta_{3n} \right) \cr
  &\mathrel{\phantom{=}} \mathord{}
     + F \left( {\delta S_2 \over \delta \eta_1},\ldots\,
                {\delta S_2 \over \delta \eta_{3n}},
                \eta_1,\ldots\, \eta_{3n} \right),\cr}
   \right\}
   \eqno {\rm (D.)}$$
which gives rigorously
$$\left. \eqalign{
S_2 &= \int_0^t \left\{
        - {\delta S_1 \over \delta t}
        + U(\eta_1,\ldots\, \eta_{3n})
        - F \left( {\delta S_1 \over \delta \eta_1},\ldots\,
                   {\delta S_1 \over \delta \eta_{3n}},
                   \eta_1,\ldots\, \eta_{3n} \right)
     \right\} dt \cr
   &\mathrel{\phantom{=}} \mathord{}
      + \int_0^t
          F \left( {\delta S_2 \over \delta \eta_1},\ldots\,
                   {\delta S_2 \over \delta \eta_{3n}},
                   \eta_1,\ldots\, \eta_{3n} \right) dt,\cr}
   \right\}
   \eqno {\rm (E.)}$$
supposing only that the two parts $S_1$, $S_2$, like the whole
principal function~$S$, are chosen so as to vanish with the time.

This general and rigorous transformation offers a general method
of improving an approximate expression for the principal function~$S$,
in any problem of dynamics.  For if the part~$S_1$ be such an approximate
expression, then the remaining part~$S_2$ will be small; and the
homogeneous function~$F$ involving the squares and products of the
coefficients of this small part, in the second definite integral
(E.), will be in general also small, and of a higher order of
smallness; we may therefore in general neglect this second definite
integral, in passing to a second approximation, and may in general
improve a first approximate expresssion~$S_1$ by adding to it
the following correction,
$$\Delta S_1 = \int_0^t \left\{
        - {\delta S_1 \over \delta t}
        + U(\eta_1,\ldots\, \eta_{3n})
        - F \left( {\delta S_1 \over \delta \eta_1},\ldots\,
                   {\delta S_1 \over \delta \eta_{3n}},
                   \eta_1,\ldots\, \eta_{3n} \right)
     \right\} dt;
   \eqno {\rm (F.)}$$
in calculating which definite integral we may employ the following
approximate forms for the integrals of the equations of motion,
$$p_1    = - {\delta S_1 \over \delta e_1},\quad
  p_2    = - {\delta S_1 \over \delta e_2},\quad \ldots \quad
  p_{3n} = - {\delta S_1 \over \delta e_{3n}},
   \eqno {\rm (39.)}$$
expressing first, by these, the variables $\eta_i$ as functions
of the time and of the $6n$ constants $e_i$~$p_i$, and then
eliminating, after the integration, the $3n$ quantities $p_i$,
by the same approximate forms.  And when an improved expression,
or second approximate value $S_1 + \Delta S_1$, for the principal
function $S$, has been thus obtained, it may be substituted in
like manner for the first approximate value~$S_1$, so as to obtain
a still closer approximation, and the process may be repeated
indefinitely.

An analogous process applies to the indefinite improvement of a
first approximate expression for the function $Q$ or $V$.

\bigbreak

{\sectiontitle
Rigorous Theory of Perturbations, founded on the Properties
of the Disturbing Part of the whole Principal Function.\par}

\nobreak\bigskip

8.
If we separate the expression $H$ (17.) into any two
parts of the same kind,
$$H_1 + H_2 = H,
   \eqno {\rm (40.)}$$
in which
$$H_1 = F_1(\varpi_1, \varpi_2,\ldots\, \varpi_{3n},
            \eta_1, \eta_2,\ldots\, \eta_{3n})
        - U_1(\eta_1, \eta_2,\ldots\, \eta_{3n}),
   \eqno {\rm (41.)}$$
and
$$H_2 = F_2(\varpi_1, \varpi_2,\ldots\, \varpi_{3n},
            \eta_1, \eta_2,\ldots\, \eta_{3n})
        - U_2(\eta_1, \eta_2,\ldots\, \eta_{3n}),
   \eqno {\rm (42.)}$$
the functions $F_1$~$F_2$~$U_1$~$U_2$ being such that
$$F_1 + F_2 = F,\quad U_1 + U_2 = U;
   \eqno {\rm (43.)}$$
the differential equations of motion (A.) will take this form,
$${d \eta_i \over dt}
   =   {\delta H_1 \over \delta \varpi_i}
     + {\delta H_2 \over \delta \varpi_i},\quad
{d \varpi_i \over dt}
   = - {\delta H_1 \over \delta \eta_i}
     - {\delta H_2 \over \delta \eta_i},
   \eqno {\rm (G.)}$$
and if the part~$H_2$ and its coefficients be small, they will not
differ much from these other differential equations,
$${d \eta_i \over dt}
   =   {\delta H_1 \over \delta \varpi_i},\quad
{d \varpi_i \over dt}
   = - {\delta H_1 \over \delta \eta_i};
   \eqno {\rm (H.)}$$
so that the rigorous integrals of the latter system will be
approximate integrals of the former.  Whenever then, by a
proper choice of the predominant term $H_1$, a system of
$6n$ equations such as (H.) has been formed and rigorously
integrated, giving expressions for the $6n$ variables
$\eta_i$~$\varpi_i$ as functions of the time~$t$, and of their
own initial values $e_i$~$p_i$, which may be thus denoted:
$$\eta_i   = \phi_i(t, e_1, e_2,\ldots\, e_{3n},
                    p_1, p_2,\ldots\, p_{3n}),
   \eqno {\rm (44.)}$$
and
$$\varpi_i = \psi_i(t, e_1, e_2,\ldots\, e_{3n},
                    p_1, p_2,\ldots\, p_{3n});
   \eqno {\rm (45.)}$$
the simpler motion thus defined by the rigorous integrals of (H.)
may be called the {\it undisturbed motion\/} of the proposed system
of $n$ points, and the more complex motion expressed by the rigorous
integrals of (G.) may be called by contrast the {\it disturbed motion\/}
of that system; and to pass from one to the other, may be called a
{\it Problem of Perturbation}.

\medskip

9.
To accomplish this passage, let us observe that the differential
equations of undisturbed motion (H.), being of the same form as the
original equations (A.), may have their integrals similarly expressed,
that is, as follows:
$$\varpi_i =   {\delta S_1 \over \delta \eta_i},\quad
  p_i      = - {\delta S_1 \over \delta e_i},
   \eqno {\rm (I.)}$$
$S_1$ being here the {\it principal function of undisturbed motion},
or the definite integral
$$S_1 = \int_0^t \left(
            \sum \mathbin{.} \varpi {\delta H_1 \over \delta \varpi} - H_1
            \right) dt,
   \eqno {\rm (46.)}$$
considered as a function of the time and of the
quantities $\eta_i$~$e_i$.  In like manner if we represent
by $S_1 +S_2$ the whole principal function of disturbed motion,
the rigorous integrals of (G.) may be expressed by (B.), as follows:
$$\varpi_i =   {\delta S_1 \over \delta \eta_i}
             + {\delta S_2 \over \delta \eta_i},\quad
  p_i      = - {\delta S_1 \over \delta e_i}
             - {\delta S_2 \over \delta e_i}.
   \eqno {\rm (K.)}$$
Comparing the forms (44.) with the second set of equations (I.)
for the integrals of undisturbed motion, we find that the following
relations between the functions $\phi_i$~$S_1$ must be rigorously
and {\it identically\/} true:
$$\eta_i   = \phi_i \left( t, e_1, e_2,\ldots\, e_{3n},
                           - {\delta S_1 \over \delta e_1},
                           - {\delta S_1 \over \delta e_2},\ldots\,
                           - {\delta S_1 \over \delta e_{3n}}
                           \right);
   \eqno {\rm (47.)}$$
and therefore, by (K.), that the integrals of disturbed motion
may be put under the following forms,
$$\eta_i   = \phi_i \left( t, e_1, e_2,\ldots\, e_{3n},
                           p_1    + {\delta S_2 \over \delta e_1},
                           p_2    + {\delta S_2 \over \delta e_2},\ldots\,
                           p_{3n} + {\delta S_2 \over \delta e_{3n}}
                           \right).
   \eqno {\rm (L.)}$$
We may therefore calculate rigorously the disturbed variables $\eta_i$
by the rules of undisturbed motion (44.), if without altering the
time~$t$, or the initial values $e_i$ of those variables, which
determine the initial configuration, we alter (in general) the
initial velocities and directions, by adding to the elements~$p_i$
the following perturbational terms,
$$\Delta p_1 = {\delta S_2 \over \delta e_1},\quad
  \Delta p_2 = {\delta S_2 \over \delta e_2},\quad \ldots \quad
  \Delta p_{3n} = {\delta S_2 \over \delta e_{3n}}:
   \eqno {\rm (M.)}$$
a remarkable result, which includes the whole theory of perturbation.
We might deduce from it the differential coefficients $\eta_i'$, or
the connected quantities $\varpi_i$, which determine the disturbed
directions and velocities of motion at any time~$t$; but a similar
reasoning gives at once the general expression,
$$\varpi_i = {\delta S_2 \over \delta \eta_i}
      + \psi_i \left( t, e_1, e_2,\ldots\,e_{3n},
                      p_1    + {\delta S_2 \over \delta e_1},
                      p_2    + {\delta S_2 \over \delta e_2},\ldots\,
                      p_{3n} + {\delta S_2 \over \delta e_{3n}}
                      \right),
   \eqno {\rm (N.)}$$
implying, that after altering the initial velocities and directions
or the elements $p_i$ as before, by the perturbational terms (M.),
we may then employ the rules of undisturbed motion (45.) to
calulate the velocities and directions at the time~$t$, or the
varying quantities $\varpi_i$, if we finally apply to these
quantities thus calculated the following new corrections for
perturbation:
$$\Delta \varpi_1 = {\delta S_2\over \delta \eta_1},\quad
  \Delta \varpi_2 = {\delta S_2\over \delta \eta_2},\quad \ldots \quad
  \Delta \varpi_{3n} = {\delta S_2\over \delta \eta_{3n}}.
   \eqno {\rm (O.)}$$

\bigbreak

{\sectiontitle
Approximate expressions deduced from the foregoing rigorous Theory.\par}

\nobreak\bigskip

10.
The foregoing theory gives indeed rigorous expressions for the
perturbations, in passing from the simpler motion (H.) or (I.)
to the more complex motion (G.) or (K.): but it may seem that these
expressions are of little use, because they involve an unknown
{\it disturbing function\/} $S_2$, (namely, the perturbational part
of the whole principal function $S$,) and also unknown or disturbed
coordinates or marks of position~$\eta_i$.  However, it was lately
shown that whenever a first approximate form for the principal
function~$S$, such as here the principal function $S_1$ of
undisturbed motion, has been found, the correction $S_2$ can in
general be assigned, with an indefinitely increasing accuracy;
and since the perturbations (M.) and (O.) involve the disturbed
coordinates $\eta_i$ only as they enter into the coefficients
of this small disturbing function $S_2$, it is evidently
permitted to substitute for these coordinates, at first, their
undisturbed values, and then to correct the results by
substituting more accurate expressions.

\medskip

11.
The function $S_1$ of undisturbed motion must satisfy rigorously
two partial differential equations of the form (C.), namely,
$$\left. \eqalign{
{\delta S_1 \over \delta t}
   + F_1 \left( {\delta S_1 \over \delta \eta_1},\ldots\,
                {\delta S_1 \over \delta \eta_{3n}},
                \eta_1,\ldots\, \eta_{3n}
                \right)
   &= U_1(\eta_1,\ldots\, \eta_{3n}),\cr
{\delta S_1 \over \delta t}
   + F_1 \left( {\delta S_1 \over \delta e_1},\ldots\,
                {\delta S_1 \over \delta e_{3n}},
                e_1,\ldots\, e_{3n}
                \right)
   &= U_1(e_1,\ldots\, e_{3n});\cr}
   \right\}
   \eqno {\rm (P.)}$$
and therefore, by (D.), the disturbing function $S_2$ must
satisfy rigorously the following other condition:
$$\left. \eqalign{
{d S_2 \over dt}
  &= U_2(\eta_1,\ldots\, \eta_{3n})
     - F_2 \left( {\delta S_1 \over \delta \eta_1},\ldots\,
                  {\delta S_1 \over \delta \eta_{3n}},
                  \eta_1,\ldots\, \eta_{3n} \right) \cr
  &\mathrel{\phantom{=}} \mathord{}
     + F \left( {\delta S_2 \over \delta \eta_1},\ldots\,
                {\delta S_2 \over \delta \eta_{3n}},
                \eta_1,\ldots\, \eta_{3n} \right),\cr}
   \right\}
   \eqno {\rm (Q.)}$$
and may, on account of the homogeneity and dimension of $F$,
be approximately expressed as follows:
$$S_2 = \int_0^t \left\{
      U_2(\eta_1,\ldots\, \eta_{3n})
      - F_2 \left( {\delta S_1 \over \delta \eta_1},\ldots\,
                   {\delta S_1 \over \delta \eta_{3n}},
                   \eta_1,\ldots\, \eta_{3n} \right)
      \right\} dt,
   \eqno {\rm (R.)}$$
or thus, by (I.),
$$S_2 = \int_0^t \left\{
      U_2(\eta_1,\ldots\, \eta_{3n})
      - F_2 \left( \varpi_1,\ldots\, \varpi_{3n},
                   \eta_1,\ldots\, \eta_{3n} \right)
      \right\} dt,
   \eqno {\rm (S.)}$$
that is, by (42.),
$$S_2 = - \int_0^t H_2 \,dt.
   \eqno {\rm (T.)}$$
In this expression, $H_2$ is given immediately as a function
of the varying quantities $\eta_i$~$\varpi_i$, but it may
be considered in the same order of approximation as a known
function of their initial values $e_i$~$p_i$ and of the
time~$t$, obtained by substituting for $\eta_i$~$\varpi_i$
their undisturbed values (44.), (45.) as functions of those
quantities; its variation may therefore be expressed in either
of the two following ways:
$$\delta H_2 = \sum \left(
           {\delta H_2 \over \delta \eta} \delta \eta
         + {\delta H_2 \over \delta \varpi} \delta \varpi
           \right),
   \eqno {\rm (48.)}$$
or
$$\delta H_2 = \sum \left(
           {\delta H_2 \over \delta e} \delta e
         + {\delta H_2 \over \delta p} \delta p
           \right)
         + {\delta H_2 \over \delta t} \delta t.
   \eqno {\rm (49.)}$$

Adopting the latter view, and effecting the integration (T.)
with respect to the time, by treating the elements $e_i$~$p_i$
as constant, we are afterwards to substitute for the quantities
$p_i$ their undisturbed expressions (39.) or (I.), and then
we find for the variation of the disturbing function $S_2$
the expression
$$\delta S_2 = - H_2 \, \delta t + \sum \left(
         - \delta e \mathbin{.} \int_0^t {\delta H_2 \over \delta e} dt
         + \delta {\delta S_1 \over \delta e}
            \mathbin{.} \int_0^t {\delta H_2 \over \delta p} dt
         \right),
   \eqno {\rm (50.)}$$
which enables us to transform the perturbational terms (M.), (O.)
into the following approximate forms:
$$\Delta p_i = - \int_0^t {\delta H_2 \over \delta e_i} dt
      + \sum \mathbin{.} {\delta^2 S_1 \over \delta e \, \delta e_i}
            \int_0^t {\delta H_2 \over \delta p} dt,
   \eqno {\rm (U.)}$$
and
$$\Delta \varpi_i
   = \sum \mathbin{.} {\delta^2 S_1 \over \delta e \, \delta \eta_i}
            \int_0^t {\delta H_2 \over \delta p} dt,
   \eqno {\rm (V.)}$$
containing only functions and quantities which may be regarded
as given, by the theory of undisturbed motion.

\medskip

12.
In the same order of approximation, if the variation of
the expression (44.) for an undisturbed coordinate $\eta_i$
be thus denoted,
$$\delta \eta_i
   = {\delta \eta_i \over \delta t} \delta t + \sum \left(
         {\delta \eta_i \over \delta e} \delta e
       + {\delta \eta_i \over \delta p} \delta p \right),
   \eqno {\rm (51.)}$$
the perturbation of that coordinate may be expressed as follows:
$$\Delta \eta_i
   = \sum \mathbin{.} {\delta \eta_i \over \delta p} \Delta p;
   \eqno {\rm (W.)}$$
that is, by (U.),
$$\left. \eqalign{\Delta \eta_i
   &=  - {\delta \eta_i \over \delta p_1}
         \int_0^t {\delta H_2 \over \delta e_1} dt
       - {\delta \eta_i \over \delta p_2}
         \int_0^t {\delta H_2 \over \delta e_2} dt
       - \cdots
       - {\delta \eta_i \over \delta p_{3n}}
         \int_0^t {\delta H_2 \over \delta e_{3n}} dt \cr
   &\mathrel{\phantom{=}} \mathord{}
    + \left(
         {\delta \eta_i \over \delta p_1}
         {\delta^2 S_1 \over \delta e_1^2}
       + {\delta \eta_i \over \delta p_2}
         {\delta^2 S_1 \over \delta e_1 \, \delta e_2}
       + \cdots
       + {\delta \eta_i \over \delta p_{3n}}
         {\delta^2 S_1 \over \delta e_1 \, \delta e_{3n}}
         \right) \int_0^t {\delta H_2 \over \delta p_1} dt \cr
   &\mathrel{\phantom{=}} \mathord{}
    + \ldots\ldots \cr
   &\mathrel{\phantom{=}} \mathord{}
    + \left(
         {\delta \eta_i \over \delta p_1}
         {\delta^2 S_1 \over \delta e_{3n} \, \delta e_1}
       + {\delta \eta_i \over \delta p_2}
         {\delta^2 S_1 \over \delta e_{3n} \, \delta e_2}
       + \cdots
       + {\delta \eta_i \over \delta p_{3n}}
         {\delta^2 S_1 \over \delta e_{3n}^2}
         \right) \int_0^t {\delta H_2 \over \delta p_{3n}} dt.\cr}
   \right\}
   \eqno {\rm (52.)}$$
Besides, the identical equation (47.) gives
$${\delta \eta_i \over \delta e_k}
   =   {\delta \eta_i \over \delta p_1}
       {\delta^2 S_1 \over \delta e_k \, \delta e_1}
     + {\delta \eta_i \over \delta p_2}
       {\delta^2 S_1 \over \delta e_k \, \delta e_2}
     + \cdots
     + {\delta \eta_i \over \delta p_{3n}}
       {\delta^2 S_1 \over \delta e_k \, \delta e_{3n}};
   \eqno {\rm (53.)}$$
the expression (52.) may therefore be thus abridged,
$$\left. \eqalign{
\Delta \eta_i
   &=  - {\delta \eta_i \over \delta p_1}
         \int_0^t {\delta H_2 \over \delta e_1} dt
       - \cdots
       - {\delta \eta_i \over \delta p_{3n}}
         \int_0^t {\delta H_2 \over \delta e_{3n}} dt \cr
   &\mathrel{\phantom{=}} \mathord{}
       + {\delta \eta_i \over \delta e_1}
         \int_0^t {\delta H_2 \over \delta p_1} dt
       + \cdots
       + {\delta \eta_i \over \delta e_{3n}}
         \int_0^t {\delta H_2 \over \delta p_{3n}} dt,\cr}
   \right\}
   \eqno {\rm (X.)}$$
and shows that instead of the rigorous perturbational terms (M.)
we may approximately employ the following,
$$\Delta p_i
   = - \int_0^t {\delta H_2 \over \delta e_i} dt,
   \eqno {\rm (Y.)}$$
in order to calculate the disturbed configuration at any time~$t$
by the rules of undisturbed motion, provided that besides thus
altering the initial velocities and directions we alter also the
initial configuration, by the formula
$$\Delta e_i
   =   \int_0^t {\delta H_2 \over \delta p_i} dt.
   \eqno {\rm (Z.)}$$
It would not be difficult to calculate, in like manner, approximate
expressions for the disturbed directions and velocities at any
time~$t$; but it is better to resume, in another way, the
rigorous problem of perturbation.

\bigbreak

{\sectiontitle
Other Rigorous Theory of Perturbation, founded on the properties
of the disturbing part of the constant of living force, and giving
formul{\ae} for the Variation of Elements more analogous to
those already known.\par}

\nobreak\bigskip

13.
Suppose that the theory of undisturbed motion has given the
$6n$ constants $e_i$~$p_i$ or any combinations of these,
$\kappa_1, \kappa_2,\ldots, \kappa_{6n}$, as functions of the
$6n$ variables $\eta_i$~$\varpi_i$, and of the time~$t$, which
may be thus denoted:
$$\kappa_i = \chi_i(t,\eta_1, \eta_2,\ldots\, \eta_{3n},
                    \varpi_1, \varpi_2,\ldots\, \varpi_{3n}),
   \eqno {\rm (54.)}$$
and which give reciprocally expressions for the variables
$\eta_i$~$\varpi_i$ in terms of these elements and of the time,
analogous to (44.) and (45.), and capable of being denoted
similarly,
$$\eta_i   = \phi_i(t,\kappa_1, \kappa_2,\ldots\, \kappa_{6n}),\quad
  \varpi_i = \psi_i(t,\kappa_1, \kappa_2,\ldots\, \kappa_{6n});
   \eqno {\rm (55.)}$$
then, the total differential coefficient of every such {\it element}
or function $\kappa_i$, taken with respect to the time, (both as
it enters explicitly and implicitly into the expressions (54.),)
must vanish in the undisturbed motion; so that, by the differential
equations of such motion (H.), the following general relation must
be rigorously and {\it identically\/} true:
$$0  = {\delta \kappa_i \over \delta t} + \sum \left(
            {\delta \kappa_i \over \delta \eta}
            {\delta H_1 \over \delta \varpi}
          - {\delta \kappa_i \over \delta \varpi}
            {\delta H_1 \over \delta \eta}
   \right).
   \eqno {\rm (56.)}$$

In passing to disturbed motion, if we retain the equation (54.)
as a {\it definition\/} of the quantity $\kappa_i$, that quantity
will no longer be constant, but it will continue to satisfy the
inverse relations (55.), and may be called, by analogy, a
{\it varying element\/} of the motion; and its total differential
coefficient, taken with respect to the time, may, by the identical
equation (56.), and by the differential equations of disturbed motion
(G.), be rigorously expressed as follows:
$${d\kappa_i \over dt}
   = \sum \left( {\delta \kappa_i \over \delta \eta}
                 {\delta H_2 \over \delta \varpi}
               - {\delta \kappa_i \over \delta \varpi}
                 {\delta H_2 \over \delta \eta}
   \right).
   \eqno {\rm (A^1.)}$$

\medskip

14.
This result (A${}^1$.) contains the whole theory of the gradual
variation of the elements of disturbed motion of a system; but it
may receive an advantageous transformation, by the substitution
of the expressions (55.) for the variables $\eta_i$~$\varpi_i$
as functions of the time and of the elements; since it will thus
conduct to a system of $6n$ rigorous and ordinary differential
equations of the first order between those varying elements
and the time.  Expressing, therefore, the quantity $H_2$ as
a function of these latter variables, its variation $\delta H_2$
takes this new form,
$$\delta H_2
   = \sum \mathbin{.} {\delta H_2 \over \delta \kappa} \delta \kappa
     + {\delta H_2 \over \delta t} \delta t,
   \eqno {\rm (57.)}$$
and gives, by comparison with the form (48.), and by (54.),
$${\delta H_2 \over \delta \eta_r}
   = \sum \mathbin{.} {\delta H_2 \over \delta \kappa}
            {\delta \kappa \over \delta \eta_r};\quad
  {\delta H_2 \over \delta \varpi_r}
   = \sum \mathbin{.} {\delta H_2 \over \delta \kappa}
            {\delta \kappa \over \delta \varpi_r};
   \eqno {\rm (58.)}$$
and thus the general equation (A${}^1$.) is transformed to
the following,
$${d\kappa_i \over dt}
   =   a_{i,1}  {\delta H_2 \over \delta \kappa_1}
     + a_{i,2}  {\delta H_2 \over \delta \kappa_2}
     + \cdots
     + a_{i,6n} {\delta H_2 \over \delta \kappa_{6n}},
   \eqno {\rm (B^1.)}$$
in which
$$a_{i,s} = \sum \left(
         {\delta \kappa_i \over \delta \eta}
         {\delta \kappa_s \over \delta \varpi}
       - {\delta \kappa_i \over \delta \varpi}
         {\delta \kappa_s \over \delta \eta}
         \right):
   \eqno {\rm (C^1.)}$$
so that it only remains to eliminate the variables $\eta$~$\varpi$
from the expressions of these latter coefficients.  Now it is
remarkable that this elimination removes the symbol~$t$ also,
and leaves the coefficients $a_{i,s}$ expressed as functions
of the elements $\kappa$ alone, not explicitly involving
the time.  This general theorem of dynamics, which is, perhaps,
a little more extensive than the analogous results discovered by
{\sc Lagrange} and by {\sc Poisson}, since it does not limit the
disturbing terms in the differential equations of motion to depend
on the configuration only, may be investigated in the following way.

\medskip

\def\hamiltonsum#1#2#3{\sum
      \lower.5ex\hbox{${}_{#1\,#2}$}\raise.5ex\llap{${}^{#3}$}}

15.
The sign of summation $\sum$ in (C${}^1$.), like the same
sign in those other analogous equations in which it has already
occurred without an index in this Essay, refers not to the
expressed indices, such as here $i$, $s$, in the quantity to
be summed, but to an index which is not expressed and which
may be here called~$r$; so that if we introduce for greater
clearness this variable index and its limits, the expression
(C${}^1$.) becomes
$$a_{i,s} = \hamiltonsum{(r)}{1}{3n} \left(
         {\delta \kappa_i \over \delta \eta_r}
         {\delta \kappa_s \over \delta \varpi_r}
       - {\delta \kappa_i \over \delta \varpi_r}
         {\delta \kappa_s \over \delta \eta_r}
         \right):
   \eqno {\rm (59.)}$$
and its total differential coefficient, taken with respect
to time, may be separated into the two following parts,
$$\left. \eqalign{
{d \over dt} a_{i,s}
   &= \hamiltonsum{(r)}{1}{3n} \left(
         {\delta \kappa_i \over \delta \eta_r}
         {d \over dt}
         {\delta \kappa_s \over \delta \varpi_r}
       - {\delta \kappa_s \over \delta \eta_r}
         {d \over dt}
         {\delta \kappa_i \over \delta \varpi_r}
         \right) \cr
   &\mathrel{+} \hamiltonsum{(r)}{1}{3n} \left(
         {\delta \kappa_s \over \delta \varpi_r}
         {d \over dt}
         {\delta \kappa_i \over \delta \eta_r}
       - {\delta \kappa_i \over \delta \varpi_r}
         {d \over dt}
         {\delta \kappa_s \over \delta \eta_r}
         \right),\cr}
   \right\}
   \eqno {\rm (60.)}$$
which we shall proceed to calculate separately, and then to add
them together.  By the definition (54.), and the differential
equations of disturbed motion (G.),
$${d \over dt} {\delta \kappa_i \over \delta \varpi_r}
   = {\delta^2 \kappa_i \over \delta t \, \delta \varpi_r}
      + \hamiltonsum{(u)}{1}{3n} \left\{
         {\delta^2 \kappa_i \over \delta \eta_u \, \delta \varpi_r}
         \left(   {\delta H_1 \over \delta \varpi_u}
                + {\delta H_2 \over \delta \varpi_u} \right)
      - {\delta^2 \kappa_i \over \delta \varpi_u \, \delta \varpi_r}
         \left(   {\delta H_1 \over \delta \eta_u}
                + {\delta H_2 \over \delta \eta_u} \right)
      \right\},
   \eqno {\rm (61.)}$$
in which, by the identical equation (56.),
$${\delta^2 \kappa_i \over \delta t \, \delta \varpi_r}
   = - {\delta \over \delta \varpi_r}
      \hamiltonsum{(u)}{1}{3n} \left(
         {\delta \kappa_i \over \delta \eta_u}
         {\delta H_1 \over \delta \varpi_u}
       - {\delta \kappa_i \over \delta \varpi_u}
         {\delta H_1 \over \delta \eta_u}
         \right);
   \eqno {\rm (62.)}$$
we have therefore
$${d \over dt} {\delta \kappa_i \over \delta \varpi_r}
   = \hamiltonsum{(u)}{1}{3n} \left(
         {\delta^2 \kappa_i \over \delta \eta_u \, \delta \varpi_r}
         {\delta H_2 \over \delta \varpi_u}
       - {\delta^2 \kappa_i \over \delta \varpi_u \, \delta \varpi_r}
         {\delta H_2 \over \delta \eta_u}
       + {\delta \kappa_i \over \delta \varpi_u}
         {\delta^2 H_1 \over \delta \eta_u \delta \varpi_r}
       - {\delta \kappa_i \over \delta \eta_u}
         {\delta^2 H_1 \over \delta \varpi_u \delta \varpi_r}
         \right),
   \eqno {\rm (63.)}$$
and
$\displaystyle {d \over dt} {\delta \kappa_s \over \delta \varpi_r}$
may be found from this, by merely changing $i$ to $s$: so that
$$\left. \eqalign{
\hamiltonsum{(r)}{1}{3n} \left(
         {\delta \kappa_i \over \delta \eta_r}
         {d \over dt}
         {\delta \kappa_s \over \delta \varpi_r}
       - {\delta \kappa_s \over \delta \eta_r}
         {d \over dt}
         {\delta \kappa_i \over \delta \varpi_r}
         \right)
   \hskip -14em \cr
   &= \hamiltonsum{(r,u)}{1,1}{3n,3n} \biggl\{ \left(
         {\delta \kappa_s \over \delta \eta_r}
         {\delta^2 \kappa_i \over \delta \varpi_u \, \delta \varpi_r}
       - {\delta \kappa_i \over \delta \eta_r}
         {\delta^2 \kappa_s \over \delta \varpi_u \, \delta \varpi_r}
         \right) {\delta H_2 \over \delta \eta_u} \cr
   &\mathrel{\phantom{=}} \quad \mathord{}
       + \left(
         {\delta \kappa_i \over \delta \eta_r}
         {\delta^2 \kappa_s \over \delta \eta_u \, \delta \varpi_r}
       - {\delta \kappa_s \over \delta \eta_r}
         {\delta^2 \kappa_i \over \delta \eta_u \, \delta \varpi_r}
         \right) {\delta H_2 \over \delta \varpi_u} \cr
   &\mathrel{\phantom{=}} \quad \mathord{}
       + \left(
         {\delta \kappa_i \over \delta \eta_r}
         {\delta \kappa_s \over \delta \varpi_u}
       - {\delta \kappa_s \over \delta \eta_r}
         {\delta \kappa_i \over \delta \varpi_u}
         \right) {\delta^2 H_1 \over \delta \eta_u \, \delta \varpi_r}
       + \left(
         {\delta \kappa_s \over \delta \eta_r}
         {\delta \kappa_i \over \delta \eta_u}
       - {\delta \kappa_i \over \delta \eta_r}
         {\delta \kappa_s \over \delta \eta_u}
         \right) {\delta^2 H_1 \over \delta \varpi_u \, \delta \varpi_r}
   \biggr\},\cr}
   \right\}
   \eqno {\rm (64.)}$$
and similarly,
$$\left. \eqalign{
\hamiltonsum{(r)}{1}{3n} \left(
         {\delta \kappa_s \over \delta \varpi_r}
         {d \over dt}
         {\delta \kappa_i \over \delta \eta_r}
       - {\delta \kappa_i \over \delta \varpi_r}
         {d \over dt}
         {\delta \kappa_s \over \delta \eta_r}
         \right)
   \hskip -14em \cr
   &= \hamiltonsum{(r,u)}{1,1}{3n,3n} \biggl\{ \left(
         {\delta \kappa_s \over \delta \varpi_r}
         {\delta^2 \kappa_i \over \delta \eta_u \, \delta \eta_r}
       - {\delta \kappa_i \over \delta \varpi_r}
         {\delta^2 \kappa_s \over \delta \eta_u \, \delta \eta_r}
         \right) {\delta H_2 \over \delta \varpi_u} \cr
   &\mathrel{\phantom{=}} \quad \mathord{}
       + \left(
         {\delta \kappa_i \over \delta \varpi_r}
         {\delta^2 \kappa_s \over \delta \varpi_u \, \delta \eta_r}
       - {\delta \kappa_s \over \delta \varpi_r}
         {\delta^2 \kappa_i \over \delta \varpi_u \, \delta \eta_r}
         \right) {\delta H_2 \over \delta \eta_u} \cr
   &\mathrel{\phantom{=}} \quad \mathord{}
       + \left(
         {\delta \kappa_i \over \delta \varpi_r}
         {\delta \kappa_s \over \delta \eta_u}
       - {\delta \kappa_s \over \delta \varpi_r}
         {\delta \kappa_i \over \delta \eta_u}
         \right) {\delta^2 H_1 \over \delta \varpi_u \, \delta \eta_r}
       + \left(
         {\delta \kappa_s \over \delta \varpi_r}
         {\delta \kappa_i \over \delta \varpi_u}
       - {\delta \kappa_i \over \delta \varpi_r}
         {\delta \kappa_s \over \delta \varpi_u}
         \right) {\delta^2 H_1 \over \delta \eta_u \, \delta \eta_r}
   \biggr\}.\cr}
   \right\}
   \eqno {\rm (65.)}$$
Adding, therefore, the two last expressions, and making the
reductions which present themselves, we find, by (60.),
$${d \over dt} a_{i,s} = \hamiltonsum{(u)}{1}{3n} \left(
         A_{i,s}^{(u)} {\delta H_2 \over \delta \eta_u}
       + B_{i,s}^{(u)} {\delta H_2 \over \delta \varpi_u}
         \right),
   \eqno {\rm (D^1.)}$$
in which
$$\left. \eqalign{
A_{i,s}^{(u)}
   &= \hamiltonsum{(r)}{1}{3n} \left(
         {\delta \kappa_s \over \delta \eta_r}
         {\delta^2 \kappa_i \over \delta \varpi_u \, \delta \varpi_r}
       - {\delta \kappa_i \over \delta \eta_r}
         {\delta^2 \kappa_s \over \delta \varpi_u \, \delta \varpi_r}
       + {\delta \kappa_i \over \delta \varpi_r}
         {\delta^2 \kappa_s \over \delta \varpi_u \, \delta \eta_r}
       - {\delta \kappa_s \over \delta \varpi_r}
         {\delta^2 \kappa_i \over \delta \varpi_u \, \delta \eta_r}
         \right),\cr
B_{i,s}^{(u)}
   &= \hamiltonsum{(r)}{1}{3n} \left(
         {\delta \kappa_s \over \delta \varpi_r}
         {\delta^2 \kappa_i \over \delta \eta_u \, \delta \eta_r}
       - {\delta \kappa_i \over \delta \varpi_r}
         {\delta^2 \kappa_s \over \delta \eta_u \, \delta \eta_r}
       + {\delta \kappa_i \over \delta \eta_r}
         {\delta^2 \kappa_s \over \delta \eta_u \, \delta \varpi_r}
       - {\delta \kappa_s \over \delta \eta_r}
         {\delta^2 \kappa_i \over \delta \eta_u \, \delta \varpi_r}
         \right);\cr}
   \right\}
   \eqno {\rm (66.)}$$
and since this general form (D${}^1$.) for
$\displaystyle {d \over dt} a_{i,s}$
contains no term independent of the disturbing quantities
$\displaystyle {\delta H_2 \over \delta \eta}$,
$\displaystyle {\delta H_2 \over \delta \varpi}$,
it is easy to infer from it the important consequence already
mentioned, namely, that the coefficients $a_{i,s}$, in the
differentials (B${}^1$.) of the elements, may be expressed
as functions of those elements alone, not explicitly involving
the time.

It is evident also, that these coefficients $a_{i,s}$
have the property
$$a_{s,i} = - a_{i,s}
   \eqno {\rm (67.)}$$
and
$$a_{i,i} = 0;
   \eqno {\rm (68.)}$$
the term proportional to
$\displaystyle {\delta H_2 \over \delta \kappa_i}$
disappears therefore from the expression (B${}^1$.) for
$\displaystyle {d\kappa_i \over dt}$;
and the term
$${\delta H_2 \over \delta \kappa_i} \mathbin{.} a_{i,s} \mathbin{.}
  {\delta H_2 \over \delta \kappa_s}
\quad\hbox{in}\quad
  {\delta H_2 \over \delta \kappa_i} {d\kappa_i \over dt},$$
destroys the term
$${\delta H_2 \over \delta \kappa_s} \mathbin{.} a_{s,i} \mathbin{.}
  {\delta H_2 \over \delta \kappa_i}
\quad\hbox{in}\quad
  {\delta H_2 \over \delta \kappa_s} {d\kappa_s \over dt},$$
when these terms are added together; we have, therefore,
$$\sum \mathbin{.} {\delta H_2 \over \delta \kappa} {d\kappa \over dt}
   = 0,
   \eqno {\rm (E^1.)}$$
or
$${d H_2 \over dt} = {\delta H_2 \over \delta t};
   \eqno {\rm (F^1.)}$$
that is, in taking the first total differential coefficient
of the disturbing expression $H_2$ with respect to the time,
the elements may be treated as constant.

\bigbreak

{\sectiontitle
Simplification of the differential equations which determine these
gradually varying elements, in any problem of Perturbation; and
Integration of the simplified equations by means of certain
Functions of Elements.\par}

\nobreak\bigskip

16.
The most natural choice of these elements is that which
makes them correspond, in undisturbed motion, to the initial
quantities $e_i$~$p_i$.  These quantities, by the differential
equations (H.), may be expressed in undisturbed motion as follows,
$$e_i = \eta_i   - \int_0^t {\delta H_1 \over \delta \varpi_i} dt,\quad
  p_i = \varpi_i + \int_0^t {\delta H_1 \over \delta \eta_i}   dt;
   \eqno {\rm (69.)}$$
and if we suppose them found, by elimination, under the forms
$$\left. \eqalign{
e_i &= \eta_i
      + \Phi_i(t, \eta_1, \eta_2,\ldots\, \eta_{3n},
                  \varpi_1, \varpi_2,\ldots\, \varpi_{3n} ),\cr
p_i &= \varpi_i
      + \Psi_i(t, \eta_1, \eta_2,\ldots\, \eta_{3n},
                  \varpi_1, \varpi_2,\ldots\, \varpi_{3n} ),\cr}
   \right\}
   \eqno {\rm (70.)}$$
it is easy to see that the following equations must be rigorously
and identically true, for all values of $\eta_i$~$\varpi_i$,
$$\left. \eqalign{
0 &= \Phi_i(0, \eta_1, \eta_2,\ldots\, \eta_{3n},
                  \varpi_1, \varpi_2,\ldots\, \varpi_{3n} ),\cr
0 &= \Psi_i(0, \eta_1, \eta_2,\ldots\, \eta_{3n},
                  \varpi_1, \varpi_2,\ldots\, \varpi_{3n} ).\cr}
   \right\}
   \eqno {\rm (71.)}$$
When, therefore, in passing to disturbed motion, we establish the
equations of definition,
$$\left. \eqalign{
\kappa_i &= \eta_i
      + \Phi_i(t, \eta_1, \eta_2,\ldots\, \eta_{3n},
                  \varpi_1, \varpi_2,\ldots\, \varpi_{3n} ),\cr
\lambda_i &= \varpi_i
      + \Psi_i(t, \eta_1, \eta_2,\ldots\, \eta_{3n},
                  \varpi_1, \varpi_2,\ldots\, \varpi_{3n} ),\cr}
   \right\}
   \eqno {\rm (72.)}$$
introducing $6n$ varying elements $\kappa_i$~$\lambda_i$, of
which the set $\lambda_i$ would have been represented in our
recent notation as follows:
$$\lambda_i = \kappa_{3n+i};
   \eqno {\rm (73.)}$$
we see that all the partial differential coefficients of the forms
$\displaystyle {\delta \kappa_i  \over \delta \eta_r}$,
$\displaystyle {\delta \kappa_i  \over \delta \varpi_r}$,
$\displaystyle {\delta \lambda_i \over \delta \eta_r}$,
$\displaystyle {\delta \lambda_i \over \delta \varpi_r}$,
vanish when $t = 0$, except the following:
$${\delta \kappa_i  \over \delta \eta_i}   = 1,\quad
  {\delta \lambda_i \over \delta \varpi_i} = 1;
   \eqno {\rm (74.)}$$
and, therefore, that when $t$ is made $= 0$, in the coefficients
$a_{i,s}$, (59.), all those coefficients vanish, except the following:
$$a_{r, 3n+r} = 1,\quad a_{3n+r,r} = -1.
   \eqno {\rm (75.)}$$

   But it has been proved that these coefficients $a_{i,s}$, when
expressed as functions of the elements, do not contain the time
explicitly; and the supposition $t = 0$ introduces no relation
between those $6n$ elements $\kappa_i$~$\lambda_i$, which still
remain independent: the coefficients $a_{i,s}$, therefore, could
not acquire the values $1$, $0$, $-1$, by the supposition $t = 0$,
unless they had those values constantly, and independently of that
supposition.  The differential equations of the forms (B${}^1$.)
may therefore be expressed, for the present system of varying
elements, in the following simpler way:
$${d \kappa_i \over dt}  =   {\delta H_2 \over \delta \lambda_i};\quad
  {d \lambda_i \over dt} = - {\delta H_2 \over \delta \kappa_i};
   \eqno {\rm (G^1.)}$$
and an easy verification of these expressions is offered by the
formula (E${}^1$.), which takes now this form,
$$\sum \left( {\delta H_2 \over \delta \kappa}  {d\kappa  \over dt}
            + {\delta H_2 \over \delta \lambda} {d\lambda \over dt}
              \right) = 0.
   \eqno {\rm (H^1.)}$$

\medskip

17.
The initial values of the varying elements $\kappa_i$~$\lambda_i$
are evidently $e_i$~$p_i$, by the definitions (72.), and by the
identical equations (71.); the problem of integrating rigorously
the equations of disturbed motion (G.), between the variables
$\eta_i$~$\varpi_i$ and the time, or of determining these
variables as functions of the time and of their own initial values
$e_i$~$p_i$, is therefore rigorously transformed into the problem
of integrating the equations (G${}^1$.), or of determining the $6n$
elements $\kappa_i$~$\lambda_i$ as functions of the time and of the
same initial values.  The chief advantage of this transformation is,
that if the perturbations be small, the new variables (namely
the elements,) alter but little: and that, since the new
differential equations are of the same form as the old,
they may be integrated by a similar method.  Considering,
therefore, the definite integral
$$E = \int_0^t \left(
      \sum \mathbin{.} \lambda {\delta H_2 \over \delta \lambda} - H_2
      \right) dt,
   \eqno {\rm (76.)}$$
as a function of the time and of the $6n$ quantities
$\kappa_1, \kappa_2,\ldots\, \kappa_{3n}, e_1, e_2,\ldots\, e_n$,
and observing that its variation, taken with respect to the latter
quantities, may be shown by a process similar to that of the
fourth number of this Essay to be
$$\delta E = \sum ( \lambda \, \delta \kappa - p \, \delta e),
   \eqno {\rm (I^1.)}$$
we find that the rigorous integrals of the differential
equations (G${}^1$.) may be expressed in the following manner:
$$\lambda_i = {\delta E \over \delta \kappa_i},\quad
  p_i = - {\delta E \over \delta e_i},
   \eqno {\rm (K^1.)}$$
in which there enters only one unknown {\it function of elements}~$E$,
to the search and study of which single function the problem of
perturbation is reduced by this new method.

We might also have put
$$C = \int_0^t \left(
      - \sum \mathbin{.} \kappa {\delta H_2 \over \delta \kappa} + H_2
        \right) dt,
   \eqno {\rm (77.)}$$
and considered this definite integral~$C$ as a function of the
time and of the $6n$ quantities $\lambda_i$~$p_i$; and then
we should have found the following other forms for the integrals
of the differential equations of varying elements,
$$\kappa_i = + {\delta C \over \delta \lambda_i},\quad
  e_i = - {\delta C \over \delta p_i},
   \eqno {\rm (L^1.)}$$
And each of these {\it functions of elements}, $C$ and $E$, must
satisfy a certain partial differential equation, analogous to the
first equation of each pair mentioned in the sixth number of this
Essay, and deduced on similar principles.

\medskip

18.
Thus, it is evident, by the form of the function $E$, and by
the equations (K${}^1$.), (G${}^1$.), and (76.), that the
partial differential coefficient of this function, taken with
respect to the time, is
$${\delta E \over \delta t}
   = {dE \over dt}
      - \sum \mathbin{.} {\delta E \over \delta \kappa} {d\kappa \over dt}
   = - H_2;
   \eqno {\rm (M^1.)}$$
and therefore that if we separate this function~$E$ into any two parts
$$E_1 + E_2 = E,
   \eqno {\rm (N^1.)}$$
and if, for greater clearness, we put the expression $H_2$
under the form
$$H_2 = H_2(t,\kappa_1, \kappa_2,\ldots\, \kappa_{3n},
              \lambda_1, \lambda_2,\ldots\, \lambda_{3n}),
   \eqno {\rm (O^1.)}$$
we shall have rigorously the partial differential equation
$$0 =   {\delta E_1 \over \delta t}
      + {\delta E_2 \over \delta t}
      + H_2 \left( t,\kappa_1,\ldots\, \kappa_{3n},
                     {\delta E_1 \over \delta \kappa_1}
                   + {\delta E_2 \over \delta \kappa_1},\ldots\,
                     {\delta E_1 \over \delta \kappa_{3n}}
                   + {\delta E_2 \over \delta \kappa_{3n}}
                   \right):
   \eqno {\rm (P^1.)}$$
which gives, approximately, by (G${}^1$.) and (K${}^1$),
when the part $E_2$ is small, and when we neglect the squares
and products of its partial differential coefficients,
$$0 =   {d E_2 \over dt}
      + {\delta E_1 \over \delta t}
      + H_2 \left( t,\kappa_1,\ldots\, \kappa_{3n},
                     {\delta E_1 \over \delta \kappa_1},\ldots\,
                     {\delta E_1 \over \delta \kappa_{3n}}
                   \right).
   \eqno {\rm (Q^1.)}$$
Hence, in the same order of approximation, if the part $E_1$,
like the whole function $E$, be chosen so as to vanish with
the time, we shall have
$$E_2 = - \int_0^t \left\{
        {\delta E_1 \over \delta t}
      + H_2 \left( t,\kappa_1,\ldots\, \kappa_{3n},
                     {\delta E_1 \over \delta \kappa_1},\ldots\,
                     {\delta E_1 \over \delta \kappa_{3n}}
                   \right) \right\} dt:
   \eqno {\rm (R^1.)}$$
and thus a first approximate expression $E_1$ can be successively
and indefinitely corrected.

Again, by (L${}^1$.) and (G${}^1$.), and by the definition (77.),
$${\delta C \over \delta t}
   = {dC \over dt}
      - \sum \mathbin{.}
         {\delta C \over \delta \lambda} {d\lambda \over dt}
   = H_2;
   \eqno {\rm (S^1.)}$$
the function~$C$ must therefore satisfy rigorously the partial
differential equation,
$${\delta C \over \delta t}
   = H_2 \left( t, {\delta C \over \delta \lambda_1},\ldots\,
                   {\delta C \over \delta \lambda_{3n}},
                   \lambda_1,\ldots\, \lambda_{3n} \right):
   \eqno {\rm (T^1.)}$$
and if we put
$$C = C_1 + C_2,
   \eqno {\rm (U^1.)}$$
and suppose that the part $C_2$ is small, then the rigorous equation
$${\delta C_1 \over \delta t} + {\delta C_2 \over \delta t}
   = H_2 \left( t, {\delta C_1 \over \delta \lambda_1}
                 + {\delta C_2 \over \delta \lambda_1},\ldots\,
                   {\delta C_1 \over \delta \lambda_{3n}}
                 + {\delta C_2 \over \delta \lambda_{3n}},
                   \lambda_1,\ldots\, \lambda_{3n} \right)
   \eqno {\rm (V^1.)}$$
becomes approximately, by (G${}^1$.) and (L${}^1$.),
$${d C_2 \over dt}
   = - {\delta C_1 \over \delta t}
     + H_2 \left( t, {\delta C_1 \over \delta \lambda_1},\ldots\,
                     {\delta C_1 \over \delta \lambda_{3n}},
                     \lambda_1,\ldots\, \lambda_{3n} \right),
   \eqno {\rm (W^1.)}$$
and gives by integration,
$$C_2 = \int_0^t \left\{
     - {\delta C_1 \over \delta t}
     + H_2 \left( t, {\delta C_1 \over \delta \lambda_1},\ldots\,
                     {\delta C_1 \over \delta \lambda_{3n}},
                     \lambda_1,\ldots\, \lambda_{3n} \right)
   \right\} dt,
   \eqno {\rm (X^1.)}$$
the parts $C_1$ and $C_2$ being supposed to vanish separately when
$t = 0$, like the whole function of elements $C$.

And to obtain such a first approximation, $E_1$ or $C_1$, to either
of these two functions of elements $E$, $C$, we may change, in the
definitions (76.), (77.), the varying elements $\kappa$~$\lambda$,
to their initial values $e$~$p$, and then eliminate one set of
these initial values by the corresponding set of the following
approximate equations, deduced from the formul{\ae} (G${}^1$.):
$$\kappa_i = e_i + \int_0^t {\delta H_2 \over \delta p_i} dt;
   \eqno {\rm (Y^1.)}$$
and
$$\lambda_i = p_i - \int_0^t {\delta H_2 \over \delta e_i} dt.
   \eqno {\rm (Z^1.)}$$

It is easy also to see that these two functions of elements
$C$ and $E$ are connected with each other, and with the
disturbing function $S_2$, so that the form of any one
may be deduced from that of any other, when the function $S_1$
of undisturbed motion is known.


\bigbreak

{\sectiontitle
Analogous formul{\ae} for the motion of a Single Point.\par}

\nobreak\bigskip

19.
Our general method in dynamics, though intended chiefly for the
study of attracting and repelling systems, is not confined to such,
but may be used in all questions to which the law of living forces
applies.  And all the analysis of this Essay, but especially the
theory of perturbations, may usefully be illustrated by the following
analogous reasonings and results respecting the motion of a single
point.

Imagine then such a point, having for its three rectangular
coordinates $x$~$y$~$z$, and moving in an orbit determined
by three ordinary differential equations of the second order
of forms analogous to the equations (2.), namely,
$$x'' = {\delta U \over \delta x};\quad
  y'' = {\delta U \over \delta y};\quad
  z'' = {\delta U \over \delta z};
   \eqno {\rm (78.)}$$
$U$ being any given function of the coordinates not expressly
involving the time: and let us establish the following definition,
analogous to (4.),
$$T = {\textstyle {1 \over 2}} (x'^2 + y'^2 + z'^2),
   \eqno {\rm (79.)}$$
$x'$~$y'$~$z'$ being the first, and $x''$~$y''$~$z''$ being the
second differential coefficients of the coordinates, considered
as functions of the time~$t$.  If we express, for greater generality
or facility, the rectangular coordinates $x$~$y$~$z$ as functions
of three other marks of position $\eta_1$~$\eta_2$~$\eta_3$,
$T$ will become a homogeneous function of the second dimension
of their first differential coefficients $\eta_1'$~$\eta_2'$
$\eta_3'$ taken with respect to time; and if we put, for abridgement,
$$\varpi_1 = {\delta T \over \delta \eta_1'},\quad
  \varpi_2 = {\delta T \over \delta \eta_2'},\quad
  \varpi_3 = {\delta T \over \delta \eta_3'},
   \eqno {\rm (80.)}$$
$T$ may considered also as a function of the form
$$T = F(\varpi_1, \varpi_2, \varpi_3, \eta_1, \eta_2, \eta_3),
   \eqno {\rm (81.)}$$
which will be homogeneous of the second dimension with respect
to $\varpi_1$~$\varpi_2$~$\varpi_3$.  We may also put,
for abridgement,
$$T = F(\varpi_1, \varpi_2, \varpi_3, \eta_1, \eta_2, \eta_3)
      - U(\eta_1, \eta_2, \eta_3) = H;
   \eqno {\rm (82.)}$$
and then, instead of the three differential equations of the
second order (78.), we may employ the six following of the
first order, analogous to the equations (A.), and obtained
by a similar reasoning,
$$\left. \multieqalign{
{d \eta_1 \over dt}   &= + {\delta H \over \delta \varpi_1}, &
{d \eta_2 \over dt}   &= + {\delta H \over \delta \varpi_2}, &
{d \eta_3 \over dt}   &= + {\delta H \over \delta \varpi_3}, \cr
{d \varpi_1 \over dt} &= - {\delta H \over \delta \eta_1},   &
{d \varpi_2 \over dt} &= - {\delta H \over \delta \eta_2},   &
{d \varpi_3 \over dt} &= - {\delta H \over \delta \eta_3},   \cr}
   \right\}
   \eqno {\rm (83.)}$$

\medskip

20.
The rigorous integrals of these six differential equations
may be expressed under the following forms, analogous to (B.),
$$\left. \multieqalign{
\varpi_1 &=   {\delta S \over \delta \eta_1}, &
\varpi_2 &=   {\delta S \over \delta \eta_2}, &
\varpi_3 &=   {\delta S \over \delta \eta_3}, \cr
p_1      &= - {\delta S \over \delta e_1},    &
p_2      &= - {\delta S \over \delta e_2},    &
p_3      &= - {\delta S \over \delta e_3},    \cr}
   \right\}
   \eqno {\rm (84.)}$$
in which $e_1$~$e_2$~$e_3$ $p_1$~$p_2$~$p_3$ are the
initial values, or values at the time $0$, of
$\eta_1$~$\eta_2$~$\eta_3$ $\varpi_1$~$\varpi_2$~$\varpi_3$;
and $S$ is the definite integral
$$S = \int_0^t \left(
        \varpi_1 {\delta H \over \delta \varpi_1}
      + \varpi_2 {\delta H \over \delta \varpi_2}
      + \varpi_3 {\delta H \over \delta \varpi_3}
      - H \right) dt,
   \eqno {\rm (85.)}$$
considered as a function of $\eta_1$~$\eta_2$~$\eta_3$
$e_1$~$e_2$~$e_3$ and $t$.  The quantity~$H$ does not change in
the course of the motion, and the function~$S$ must satisfy the following
pair of partial differential equations of the first order, analogous
to the pair (C.),
$$\left. \eqalign{
{\delta S \over \delta t} + F \left(
      {\delta S \over \delta \eta_1},
      {\delta S \over \delta \eta_2},
      {\delta S \over \delta \eta_3},
      \eta_1, \eta_2, \eta_3 \right)
   &= U(\eta_1, \eta_2, \eta_3);\cr
{\delta S \over \delta t} + F \left(
      {\delta S \over \delta e_1},
      {\delta S \over \delta e_2},
      {\delta S \over \delta e_3},
      e_1, e_2, e_3 \right)
   &= U(e_1, e_2, e_3).\cr}
   \right\}
   \eqno {\rm (86.)}$$
This important function~$S$, which may be called the
{\it principal function\/} of the motion, may hence be
rigorously expressed under the following form, obtained by
reasonings analogous to those of the seventh number of this Essay:
$$\left. \eqalign{
S &= S_1 + \int_0^t \left\{
      - {\delta S_1 \over \delta t}
      + U( \eta_1, \eta_2, \eta_3 )
      - F \left(
         {\delta S_1 \over \delta \eta_1},
         {\delta S_1 \over \delta \eta_2},
         {\delta S_1 \over \delta \eta_3},
         \eta_1, \eta_2, \eta_3 \right)
      \right\} dt \cr
  &\phantom{= S_1} \mathord{}
    + \int_0^t
        F \left(
         {\delta S \over \delta \eta_1}
          - {\delta S_1 \over \delta \eta_1},
         {\delta S \over \delta \eta_2}
          - {\delta S_1 \over \delta \eta_2},
         {\delta S \over \delta \eta_3}
          - {\delta S_1 \over \delta \eta_3},
         \eta_1, \eta_2, \eta_3 \right) dt;\cr}
   \right\}
   \eqno {\rm (87.)}$$
$S_1$ being any arbitrary function of the same quantities
$\eta_1$~$\eta_2$~$\eta_3$ $e_1$~$e_2$~$e_3$~$t$, so chosen
as to vanish with the time.  And if this arbitrary function $S_1$
be chosen so as to be a first approximate value of the principal
function~$S$, we may neglect, in a second approximation, the
second definite integral in (87.).

\bigbreak

21.
A first approximation of this kind can be obtained, whenever, by
separating the expression~$H$, (82.) into a predominant and a
smaller part, $H_1$ and $H_2$, and by neglecting the part $H_2$,
we have changed the differential equations (83.) to others, namely,
$$\left. \multieqalign{
{d \eta_1 \over dt}   &=   {\delta H_1 \over \delta \varpi_1}, &
{d \eta_2 \over dt}   &=   {\delta H_1 \over \delta \varpi_2}, &
{d \eta_3 \over dt}   &=   {\delta H_1 \over \delta \varpi_3}, \cr
{d \varpi_1 \over dt} &= - {\delta H_1 \over \delta \eta_1},   &
{d \varpi_2 \over dt} &= - {\delta H_1 \over \delta \eta_2},   &
{d \varpi_3 \over dt} &= - {\delta H_1 \over \delta \eta_3},   \cr}
   \right\}
   \eqno {\rm (88.)}$$
and have succeeded in integrating rigorously these simplified
equations, belonging to a simpler motion, which may be called the
{\it undisturbed motion\/} of the point.  For the principal function
of such undisturbed motion, namely, the definite integral
$$S_1 = \int_0^t \left(
        \varpi_1 {\delta H_1 \over \delta \varpi_1}
      + \varpi_2 {\delta H_1 \over \delta \varpi_2}
      + \varpi_3 {\delta H_1 \over \delta \varpi_3}
      - H_1 \right) dt,
   \eqno {\rm (89.)}$$
considered as a function of
$\eta_1$~$\eta_2$~$\eta_3$ $e_1$~$e_2$~$e_3$~$t$,
will then be an approximate value for the original function
of disturbed motion $S$, which original function corresponds
to the more complex differential equations,
$$\left. \multieqalign{
{d \eta_1 \over dt}   &=   {\delta H_1 \over \delta \varpi_1}
                         + {\delta H_2 \over \delta \varpi_1}, &
{d \eta_2 \over dt}   &=   {\delta H_1 \over \delta \varpi_2}
                         + {\delta H_2 \over \delta \varpi_2}, &
{d \eta_3 \over dt}   &=   {\delta H_1 \over \delta \varpi_3}
                         + {\delta H_2 \over \delta \varpi_3}, \cr
{d \varpi_1 \over dt} &= - {\delta H_1 \over \delta \eta_1}
                         - {\delta H_2 \over \delta \eta_1},   &
{d \varpi_2 \over dt} &= - {\delta H_1 \over \delta \eta_2}
                         - {\delta H_2 \over \delta \eta_2},   &
{d \varpi_3 \over dt} &= - {\delta H_1 \over \delta \eta_3}
                         - {\delta H_2 \over \delta \eta_3}.   \cr}
   \right\}
   \eqno {\rm (90.)}$$
The function $S_1$ of undisturbed motion must satisfy a pair
of partial differential equations of the first order, analogous
to the pair (86.); and the integrals of undisturbed motion
may be represented thus,
$$\left. \multieqalign{
\varpi_1 &=   {\delta S_1 \over \delta \eta_1}, &
\varpi_2 &=   {\delta S_1 \over \delta \eta_2}, &
\varpi_3 &=   {\delta S_1 \over \delta \eta_3}, \cr
p_1      &= - {\delta S_1 \over \delta e_1},    &
p_2      &= - {\delta S_1 \over \delta e_2},    &
p_3      &= - {\delta S_1 \over \delta e_3}:    \cr}
   \right\}
   \eqno {\rm (91.)}$$
while the integrals of disturbed motion may be expressed with
equal rigour under the following analogous forms,
$$\left. \multieqalign{
\varpi_1 &=   {\delta S_1 \over \delta \eta_1}
            + {\delta S_2 \over \delta \eta_1}, &
\varpi_2 &=   {\delta S_1 \over \delta \eta_2}
            + {\delta S_2 \over \delta \eta_2}, &
\varpi_3 &=   {\delta S_1 \over \delta \eta_3}
            + {\delta S_2 \over \delta \eta_3}, \cr
p_1      &= - {\delta S_1 \over \delta e_1}
            - {\delta S_2 \over \delta e_1},    &
p_2      &= - {\delta S_1 \over \delta e_2}
            - {\delta S_2 \over \delta e_2},    &
p_3      &= - {\delta S_1 \over \delta e_3}
            - {\delta S_2 \over \delta e_3},    \cr}
   \right\}
   \eqno {\rm (92.)}$$
if $S_1$ denote the rigorous correction of $S_1$, or the
disturbing part of the whole principal function~$S$.  And by
the foregoing general theory of approximation, this disturbing
part or function $S_2$ may be approximately represented by the
definite integral (T.),
$$S_2 = - \int_0^t H_2 \,dt;
   \eqno {\rm (93.)}$$
in calculating which definite integral the equations (91.) may
be employed.

\bigbreak

22.
If the integrals of undisturbed motion (91.) have given
$$\left. \eqalign{
\eta_1   &= \phi_1(t, e_1, e_2, e_3, p_1, p_2, p_3),\cr
\eta_2   &= \phi_2(t, e_1, e_2, e_3, p_1, p_2, p_3),\cr
\eta_3   &= \phi_3(t, e_1, e_2, e_3, p_1, p_2, p_3),\cr}
   \right\}
   \eqno {\rm (94.)}$$
and
$$\left. \eqalign{
\varpi_1 &= \psi_1(t, e_1, e_2, e_3, p_1, p_2, p_3),\cr
\varpi_2 &= \psi_2(t, e_1, e_2, e_3, p_1, p_2, p_3),\cr
\varpi_3 &= \psi_3(t, e_1, e_2, e_3, p_1, p_2, p_3),\cr}
   \right\}
   \eqno {\rm (95.)}$$
then the integrals of disturbed motion (92.) may be
rigorously transformed as follows,
$$\left. \eqalign{
\eta_1
   &= \phi_1 \left( t, e_1, e_2, e_3,
         p_1 + {\delta S_2 \over \delta e_1},
         p_2 + {\delta S_2 \over \delta e_2},
         p_3 + {\delta S_2 \over \delta e_3}
         \right),\cr
\eta_2
   &= \phi_2 \left( t, e_1, e_2, e_3,
         p_1 + {\delta S_2 \over \delta e_1},
         p_2 + {\delta S_2 \over \delta e_2},
         p_3 + {\delta S_2 \over \delta e_3}
         \right),\cr
\eta_3
   &= \phi_3 \left( t, e_1, e_2, e_3,
         p_1 + {\delta S_2 \over \delta e_1},
         p_2 + {\delta S_2 \over \delta e_2},
         p_3 + {\delta S_2 \over \delta e_3}
         \right),\cr}
   \right\}
   \eqno {\rm (96.)}$$
and
$$\left. \eqalign{
\varpi_1
   &= {\delta S_2 \over \delta \eta_1}
      + \psi_1 \left( t, e_1, e_2, e_3,
         p_1 + {\delta S_2 \over \delta e_1},
         p_2 + {\delta S_2 \over \delta e_2},
         p_3 + {\delta S_2 \over \delta e_3}
         \right),\cr
\varpi_2
   &= {\delta S_2 \over \delta \eta_2}
      + \psi_2 \left( t, e_1, e_2, e_3,
         p_1 + {\delta S_2 \over \delta e_1},
         p_2 + {\delta S_2 \over \delta e_2},
         p_3 + {\delta S_2 \over \delta e_3}
         \right),\cr
\varpi_3
   &= {\delta S_2 \over \delta \eta_3}
      + \psi_3 \left( t, e_1, e_2, e_3,
         p_1 + {\delta S_2 \over \delta e_1},
         p_2 + {\delta S_2 \over \delta e_2},
         p_3 + {\delta S_2 \over \delta e_3}
         \right),\cr}
   \right\}
   \eqno {\rm (97.)}$$
$S_2$ being here the rigorous disturbing function.  And the
perturbations of position, at any time~$t$, may be approximately
expressed by the following formula,
$$\left. \eqalign{\Delta \eta_1
   &=    {\delta \eta_1 \over \delta e_1}
            \int_0^t {\delta H_2 \over \delta p_1} dt
       + {\delta \eta_1 \over \delta e_2}
            \int_0^t {\delta H_2 \over \delta p_2} dt
       + {\delta \eta_1 \over \delta e_3}
            \int_0^t {\delta H_2 \over \delta p_3} dt \cr
   &\mathrel{-}
         {\delta \eta_1 \over \delta p_1}
            \int_0^t {\delta H_2 \over \delta e_1} dt
       - {\delta \eta_1 \over \delta p_2}
            \int_0^t {\delta H_2 \over \delta e_2} dt
       - {\delta \eta_1 \over \delta p_3}
            \int_0^t {\delta H_2 \over \delta e_3} dt,\cr}
   \right\}
   \eqno {\rm (98.)}$$
together with two similar formul{\ae} for the perturbations of the
two other coordinates, or marks of position $\eta_2$, $\eta_3$.
In these formul{\ae}, the coordinates and $H_2$ are supposed
to be expressed, by the theory of undisturbed motion, as functions
of the time~$t$, and of the constants
$e_1$~$e_2$~$e_3$ $p_1$~$p_2$~$p_3$.

\bigbreak

23.
Again, if the integrals of undisturbed motion have given, by
elimination, expressions for these constants, of the forms
$$\left. \eqalign{
e_1 &= \eta_1
      + \Phi_1(t, \eta_1, \eta_2, \eta_3,
                  \varpi_1, \varpi_2, \varpi_3),\cr
e_2 &= \eta_2
      + \Phi_2(t, \eta_1, \eta_2, \eta_3,
                  \varpi_1, \varpi_2, \varpi_3),\cr
e_3 &= \eta_3
      + \Phi_3(t, \eta_1, \eta_2, \eta_3,
                  \varpi_1, \varpi_2, \varpi_3),\cr}
   \right\}
   \eqno {\rm (99.)}$$
and
$$\left. \eqalign{
p_1 &= \varpi_1
      + \Psi_1(t, \eta_1, \eta_2, \eta_3,
                  \varpi_1, \varpi_2, \varpi_3),\cr
p_2 &= \varpi_2
      + \Psi_2(t, \eta_1, \eta_2, \eta_3,
                  \varpi_1, \varpi_2, \varpi_3),\cr
p_3 &= \varpi_3
      + \Psi_3(t, \eta_1, \eta_2, \eta_3,
                  \varpi_1, \varpi_2, \varpi_3);\cr}
   \right\}
   \eqno {\rm (100.)}$$
and if, for disturbed motion, we establish the definitions
$$\left. \eqalign{
\kappa_1 &= \eta_1
      + \Phi_1(t, \eta_1, \eta_2, \eta_3,
                  \varpi_1, \varpi_2, \varpi_3),\cr
\kappa_2 &= \eta_2
      + \Phi_2(t, \eta_1, \eta_2, \eta_3,
                  \varpi_1, \varpi_2, \varpi_3),\cr
\kappa_3 &= \eta_3
      + \Phi_3(t, \eta_1, \eta_2, \eta_3,
                  \varpi_1, \varpi_2, \varpi_3),\cr}
   \right\}
   \eqno {\rm (101.)}$$
and
$$\left. \eqalign{
\lambda_1 &= \varpi_1
      + \Psi_1(t, \eta_1, \eta_2, \eta_3,
                  \varpi_1, \varpi_2, \varpi_3),\cr
\lambda_2 &= \varpi_2
      + \Psi_2(t, \eta_1, \eta_2, \eta_3,
                  \varpi_1, \varpi_2, \varpi_3),\cr
\lambda_3 &= \varpi_3
      + \Psi_3(t, \eta_1, \eta_2, \eta_3,
                  \varpi_1, \varpi_2, \varpi_3);\cr}
   \right\}
   \eqno {\rm (102.)}$$
we shall have, for such disturbed motion, the following
rigorous equations, of the forms (94.) and (95.),
$$\left. \eqalign{
\eta_1   &= \phi_1(t, \kappa_1, \kappa_2, \kappa_3,
                      \lambda_1, \lambda_2, \lambda_3),\cr
\eta_2   &= \phi_2(t, \kappa_1, \kappa_2, \kappa_3,
                      \lambda_1, \lambda_2, \lambda_3),\cr
\eta_3   &= \phi_3(t, \kappa_1, \kappa_2, \kappa_3,
                      \lambda_1, \lambda_2, \lambda_3),\cr}
   \right\}
   \eqno {\rm (103.)}$$
and
$$\left. \eqalign{
\varpi_1 &= \psi_1(t, \kappa_1, \kappa_2, \kappa_3,
                      \lambda_1, \lambda_2, \lambda_3),\cr
\varpi_2 &= \psi_2(t, \kappa_1, \kappa_2, \kappa_3,
                      \lambda_1, \lambda_2, \lambda_3),\cr
\varpi_3 &= \psi_3(t, \kappa_1, \kappa_2, \kappa_3,
                      \lambda_1, \lambda_2, \lambda_3);\cr}
   \right\}
   \eqno {\rm (104.)}$$
and may call the quantities
$\kappa_1$~$\kappa_2$~$\kappa_3$
$\lambda_1$~$\lambda_2$~$\lambda_3$
the 6 {\it varying elements\/} of the motion.  To determine these six
varying elements, we may employ the six following rigorous equations
in ordinary differentials of the first order, in which $H_2$ is
supposed to have been expressed by (103.) and (104.) as a function
of the elements and of the time:
$$\left. \multieqalign{
{d\kappa_1 \over dt}
  &=   {\delta H_2 \over \delta \lambda_1}, &
{d\kappa_2 \over dt}
  &=   {\delta H_2 \over \delta \lambda_2}, &
{d\kappa_3 \over dt}
  &=   {\delta H_2 \over \delta \lambda_3}, \cr
{d\lambda_1 \over dt}
  &= - {\delta H_2 \over \delta \kappa_1},  &
{d\lambda_2 \over dt}
  &= - {\delta H_2 \over \delta \kappa_2},  &
{d\lambda_3 \over dt}
  &= - {\delta H_2 \over \delta \kappa_3};  \cr}
   \right\}
   \eqno {\rm (105.)}$$
and the rigorous integrals of these 6 equations may be expressed
in the following manner,
$$\left. \multieqalign{
\lambda_1 &= {\delta E \over \delta \kappa_1}, &
\lambda_2 &= {\delta E \over \delta \kappa_2}, &
\lambda_3 &= {\delta E \over \delta \kappa_3}, \cr
p_1 &= - {\delta E \over \delta e_1}, &
p_2 &= - {\delta E \over \delta e_2}, &
p_3 &= - {\delta E \over \delta e_3}, \cr}
   \right\}
   \eqno {\rm (106.)}$$
the constants $e_1$~$e_2$~$e_3$ $p_1$~$p_2$~$p_3$ retaining
their recent meanings, and being therefore the initial values
of the elements
$\kappa_1$~$\kappa_2$~$\kappa_3$
$\lambda_1$~$\lambda_2$~$\lambda_3$;
while the function~$E$, which may be called the
{\it function of elements}, because its form determines the
laws of their variations, is the definite integral
$$E = \int_0^t \left(
        \lambda_1 {\delta H_2 \over \delta \lambda_1}
      + \lambda_2 {\delta H_2 \over \delta \lambda_2}
      + \lambda_3 {\delta H_2 \over \delta \lambda_3}
      - H_2 \right) dt,
   \eqno {\rm (107.)}$$
considered as depending on $\kappa_1$~$\kappa_2$~$\kappa_3$
$e_1$~$e_2$~$e_3$ and $t$.  The integrals of the equations
(105.) may also be expressed in this other way,
$$\left. \multieqalign{
\kappa_1 &= + {\delta C \over \delta \lambda_1}, &
\kappa_2 &= + {\delta C \over \delta \lambda_2}, &
\kappa_3 &= + {\delta C \over \delta \lambda_3}, \cr
e_1 &= - {\delta C \over \delta p_1}, &
e_2 &= - {\delta C \over \delta p_2}, &
e_3 &= - {\delta C \over \delta p_3}, \cr}
   \right\}
   \eqno {\rm (108.)}$$
$C$ being the definite integral
$$C = - \int_0^t \left(
        \kappa_1 {\delta H_2 \over \delta \kappa_1}
      + \kappa_2 {\delta H_2 \over \delta \kappa_2}
      + \kappa_3 {\delta H_2 \over \delta \kappa_3}
      - H_2 \right) dt,
   \eqno {\rm (109.)}$$
regarded as a function of $\lambda_1$~$\lambda_2$~$\lambda_3$
$p_1$~$p_2$~$p_3$ and $t$: and it is easy to prove that each
of these two {\it functions of elements}, $C$ and $E$, must
satisfy a partial differential equation of the first order,
which can be previously assigned, and which may assist in
discovering the forms of these two functions, and especially
in improving an approximate expression for either.  All these
results for the motion of a single point, are analogous to the
results already deduced in this Essay, for an attracting or
repelling system.

\bigbreak

{\sectiontitle
Mathematical Example, suggested by the motion of Projectiles.\par}

\nobreak\bigskip

24.
If the three marks of position $\eta_1$~$\eta_2$~$\eta_3$ of the
moving point are the rectangular coordinates themselves, and if
the function~$U$ has the form
$$U = - g\eta_3 - {\textstyle {1 \over 2}} \{
      \mu^2 (\eta_1^2 + \eta_2^2) + \nu^2 \eta_3^2 \},
   \eqno {\rm (110.)}$$
$g$, $\mu$, $\nu$ being constants; then the expression
$$H = {\textstyle {1 \over 2}}(\varpi_1^2 + \varpi_2^2 + \varpi_3^2)
      + g\eta_3 + {\textstyle {1 \over 2}} \{
      \mu^2 (\eta_1^2 + \eta_2^2) + \nu^2 \eta_3^2 \}
   \eqno {\rm (111.)}$$
is that which must be substituted in the general forms (83.),
in order to form the 6 differential equations of motion
of the first order, namely,
$$\left. \multieqalign{
{d\eta_1 \over dt} &= \varpi_1, &
{d\eta_2 \over dt} &= \varpi_2, &
{d\eta_3 \over dt} &= \varpi_3, \cr
{d\varpi_1 \over dt} &= - \mu^2 \eta_1, &
{d\varpi_2 \over dt} &= - \mu^2 \eta_2, &
{d\varpi_3 \over dt} &= -g - \nu^2 \eta_3. \cr}
   \right\}
   \eqno {\rm (112.)}$$
These differential equations have for their rigorous integrals
the six following,
$$\left. \eqalign{
\eta_1 &= e_1 \cos \mu t + {p_1 \over \mu} \sin \mu t,\cr
\eta_2 &= e_2 \cos \mu t + {p_2 \over \mu} \sin \mu t,\cr
\eta_3 &= e_3 \cos \nu t + {p_3 \over \nu} \sin \nu t
            - {g \over \nu^2} \mathop{\rm vers} \nu t,\cr}
   \right\}
   \eqno {\rm (113.)}$$
and
$$\left. \eqalign{
\varpi_1 &= p_1 \cos \mu t - \mu e_1 \sin \mu t,\cr
\varpi_2 &= p_2 \cos \mu t - \mu e_2 \sin \mu t,\cr
\varpi_3 &= p_3 \cos \nu t
      - \left( \nu e_3 + {g \over \nu} \right) \sin \nu t;\cr}
   \right\}
   \eqno {\rm (114.)}$$
$e_1$~$e_2$~$e_3$ $p_1$~$p_2$~$p_3$ being still
the initial values of
$\eta_1$~$\eta_2$~$\eta_3$ $\varpi_1$~$\varpi_2$~$\varpi_3$.

Employing these rigorous integral equations to calculate the
function~$S$, that is, by (85.) and (110.) (111.), the
definite integral
$$S = \int_0^t \left(
      {\varpi_1^2 + \varpi_2^2 + \varpi_3^2 \over 2}
      + U \right) dt,
   \eqno {\rm (115.)}$$
we find
$$\left. \eqalign{
{\textstyle {1 \over 2}} (\varpi_1^2 + \varpi_2^2 + \varpi_3^2)
   &= {\textstyle {1 \over 4}} \left\{
         p_1^2 + p_2^2 + p_3^2 + \mu^2 (e_1^2 + e_2^2)
         + \left( \nu e_3 + {g \over \nu} \right)^2 \right\} \cr
      &\mathrel{\phantom{=}} \mathord{}
       + {\textstyle {1 \over 4}} \{
         p_1^2 + p_2^2 - \mu^2(e_1^2 + e_2^2) \} \cos 2\mu t
       - {\textstyle {1 \over 2}} \mu (e_1 p_1 + e_2 p_2) \sin 2\mu t \cr
      &\mathrel{\phantom{=}} \mathord{}
       + {\textstyle {1 \over 4}} \left\{
         p_3^2 - \left(\nu e_3 + {g \over \nu} \right)^2
         \right\} \cos 2\nu t
       - {\textstyle {1 \over 2}} \left( \nu e_3 + {g \over \nu} \right)
            p_3 \sin 2\nu t,\cr}
   \right\}
   \eqno {\rm (116.)}$$
and
$$\left. \eqalign{
U  &= {g^2 \over 2 \nu^2}
       - {\textstyle {1 \over 4}} \left\{
         p_1^2 + p_2^2 + p_3^2 + \mu^2 (e_1^2 + e_2^2)
         + \left( \nu e_3 + {g \over \nu} \right)^2 \right\} \cr
      &\mathrel{\phantom{=}} \mathord{}
       + {\textstyle {1 \over 4}} \{
         p_1^2 + p_2^2 - \mu^2(e_1^2 + e_2^2) \} \cos 2\mu t
       - {\textstyle {1 \over 2}} \mu (e_1 p_1 + e_2 p_2) \sin 2\mu t \cr
      &\mathrel{\phantom{=}} \mathord{}
       + {\textstyle {1 \over 4}} \left\{
         p_3^2 - \left(\nu e_3 + {g \over \nu} \right)^2
         \right\} \cos 2\nu t
       - {\textstyle {1 \over 2}} \left( \nu e_3 + {g \over \nu} \right)
            p_3 \sin 2\nu t,\cr}
   \right\}
   \eqno {\rm (117.)}$$
and therefore,
$$\left. \eqalign{
S  &= {g^2 t \over 2 \nu^2}
       + \{ p_1^2 + p_2^2 - \mu^2(e_1^2 + e_2^2) \}
            {\sin 2\mu t \over 4 \mu}
       - {\textstyle {1 \over 2}} (e_1 p_1 + e_2 p_2)
            \mathop{\rm vers} 2\mu t \cr
      &\mathrel{\phantom{=}} \mathord{}
       + \left\{ p_3^2 - \left(\nu e_3 + {g \over \nu} \right)^2 \right\}
            {\sin 2\nu t \over 4 \nu}
       - {\textstyle {1 \over 2}} p_3 \left( e_3 + {g \over \nu^2} \right)
            \mathop{\rm vers} 2\nu t.\cr}
   \right\}
   \eqno {\rm (118.)}$$

In order, however, to express this function~$S$, as supposed by our
general method, in terms of the final and initial coordinates and
of the time, we must employ the analogous expressions for the
constants $p_1$~$p_2$~$p_3$, deduced from the integrals (113.),
namely, the following:
$$\left. \eqalign{
p_1 &= {\mu \eta_1 - \mu e_1 \cos \mu t \over \sin \mu t},\cr
p_2 &= {\mu \eta_2 - \mu e_2 \cos \mu t \over \sin \mu t},\cr
p_3 &= {\displaystyle \nu \eta_3 + {g \over \nu}
            - \left( \nu e_3 + {g \over \nu} \right) \cos \nu t
         \over \sin \nu t};\cr}
   \right\}
   \eqno {\rm (119.)}$$
and then we find
$$\left. \eqalign{
S  &= {g^2 t \over 2 \nu^2}
    + {\mu \over 2} \mathbin{.}
      {(\eta_1 - e_1)^2 + (\eta_2 - e_2)^2 \over \tan \mu t}
    + {\nu \over 2} \mathbin{.}
      {(\eta_3 - e_3)^2 \over \tan \nu t} \cr
   &\mathrel{\phantom{=}} \mathord{}
    - \mu (\eta_1 e_1 + \eta_2 e_2) \tan {\mu t \over 2}
    - \nu \left( \eta_3 + {g \over \nu^2} \right)
          \left( e_3    + {g \over \nu^2} \right)
          \tan {\nu t \over 2}.\cr}
   \right\}
   \eqno {\rm (120.)}$$

This {\it principal function\/}~$S$ satisfies the following pair
of partial differential equations of the first order, of the
kind (86.),
$$\left. \eqalign{
{\delta S \over \delta t} + {1 \over 2} \left\{
         \left( {\delta S \over \delta \eta_1} \right)^2
       + \left( {\delta S \over \delta \eta_2} \right)^2
       + \left( {\delta S \over \delta \eta_3} \right)^2 \right\}
   &= - g\eta_3 - {\mu^2 \over 2} (\eta_1^2 + \eta_2^2)
       - {\nu^2 \over 2} \eta_3^2,\cr
{\delta S \over \delta t} + {1 \over 2} \left\{
         \left( {\delta S \over \delta e_1} \right)^2
       + \left( {\delta S \over \delta e_2} \right)^2
       + \left( {\delta S \over \delta e_3} \right)^2 \right\}
   &= - ge_3 - {\mu^2 \over 2} (e_1^2 + e_2^2)
       - {\nu^2 \over 2} e_3^2,\cr}
   \right\}
   \eqno {\rm (121.)}$$
and {\it if its form had been previously found}, by the help
of this pair, or in any other way, {\it the integrals of the
equations of motion might (by our general method) have been
deduced from it}, under the forms,
$$\left. \eqalign{
\varpi_1 &= {\delta S \over \delta \eta_1}
          = \mu (\eta_1 - e_1) \mathop{ \rm cotan} \mu t
               - \mu e_1 \tan {\mu t \over 2},\cr
\varpi_2 &= {\delta S \over \delta \eta_2}
          = \mu (\eta_2 - e_2) \mathop{ \rm cotan} \mu t
               - \mu e_2 \tan {\mu t \over 2},\cr
\varpi_3 &= {\delta S \over \delta \eta_3}
          = \nu (\eta_3 - e_3) \mathop{ \rm cotan} \nu t
               - \left( \nu e_3 + {g \over \nu} \right)
               \tan {\nu t \over 2},\cr}
   \right\}
   \eqno {\rm (122.)}$$
and
$$\left. \eqalign{
p_1 &= - {\delta S \over \delta e_1}
          = \mu (\eta_1 - e_1) \mathop{ \rm cotan} \mu t
               + \mu \eta_1 \tan {\mu t \over 2},\cr
p_2 &= - {\delta S \over \delta e_2}
          = \mu (\eta_2 - e_2) \mathop{ \rm cotan} \mu t
               + \mu \eta_2 \tan {\mu t \over 2},\cr
p_3 &= - {\delta S \over \delta e_3}
          = \nu (\eta_3 - e_3) \mathop{ \rm cotan} \nu t
               + \left( \nu \eta_3 + {g \over \nu} \right)
               \tan {\nu t \over 2}:\cr}
   \right\}
   \eqno {\rm (123.)}$$
the last of these two sets of equations coinciding with the set
(119.), or (113.), and conducting, when combined with the first
set, (122.), to the other former set of integrals, (114.).

\bigbreak

25.
Suppose now, to illustrate the theory of perturbation, that the
constants $\mu$, $\nu$ are small, and that, after separating
the expression (111.) for $H$ into two parts,
$$H_1 = {\textstyle {1 \over 2}}
            (\varpi_1^2 + \varpi_2^2 + \varpi_3^2)
         + g \eta_3,
   \eqno {\rm (124.)}$$
and
$$H_2 = {\textstyle {1 \over 2}} \{
      \mu^2 (\eta_1^2 + \eta_2^2) + \nu^2 \eta_3^2 \},
   \eqno {\rm (125.)}$$
we suppress at first the small part $H_2$, and so form, by (88.),
these other and simpler differential equations of a motion
which we shall call {\it undisturbed\/}:
$$\left. \multieqalign{
{d\eta_1 \over dt} &= \varpi_1, &
{d\eta_2 \over dt} &= \varpi_2, &
{d\eta_3 \over dt} &= \varpi_3, \cr
{d\varpi_1 \over dt} &= 0, &
{d\varpi_2 \over dt} &= 0, &
{d\varpi_3 \over dt} &= -g. \cr}
   \right\}
   \eqno {\rm (126.)}$$

These new equations have for their rigorous integrals, of the
forms (94.) and (95.),
$$\eta_1 = e_1 + p_1 t,\quad
  \eta_2 = e_2 + p_2 t,\quad
  \eta_3 = e_3 + p_3 t - {\textstyle {1 \over 2}} g t^2,
   \eqno {\rm (127.)}$$
and
$$\varpi_1 = p_1,\quad
  \varpi_2 = p_2,\quad
  \varpi_3 = p_3 - gt;
   \eqno {\rm (128.)}$$
and the {\it principal function\/} $S_1$ of the same undisturbed
motion is, by (89.),
$$\left. \eqalign{
S_1 &= \int_0^t \left(
         {\varpi_1^2 + \varpi_2^2 + \varpi_3^2 \over 2}
         - g \eta_3 \right) dt \cr
    &= \int_0^t \left(
         {p_1^2 + p_2^2 + p_3^2 \over 2}
         - g e_3 - 2 g p_3 t + g^2 t^2 \right) dt \cr
    &= \left( {p_1^2 + p_2^2 + p_3^2 \over 2} - g e_3 \right) t
       - g p_3 t^2 + {\textstyle {1 \over 3}}g^2 t^3,\cr}
   \right\}
   \eqno {\rm (129.)}$$
or finally, by (127.),
$$S_1 = {(\eta_1 - e_1)^2 + (\eta_2 - e_2)^2 + (\eta_3 - e_3)^2
            \over 2t}
   - {\textstyle {1 \over 2}} gt (\eta_3 + e_3)
   - {\textstyle {1 \over 24}} g^2 t^3.
   \eqno {\rm (130.)}$$

This function satisfies, as it ought, the following pair of
partial differential equations,
$$\left. \eqalign{
{\delta S_1 \over \delta t} + {1 \over 2} \left\{
         \left( {\delta S_1 \over \delta \eta_1} \right)^2
       + \left( {\delta S_1 \over \delta \eta_2} \right)^2
       + \left( {\delta S_1 \over \delta \eta_3} \right)^2 \right\}
   &= - g\eta_3,\cr
{\delta S_1 \over \delta t} + {1 \over 2} \left\{
         \left( {\delta S_1 \over \delta e_1} \right)^2
       + \left( {\delta S_1 \over \delta e_2} \right)^2
       + \left( {\delta S_1 \over \delta e_3} \right)^2 \right\}
   &= - ge_3.\cr}
   \right\}
   \eqno {\rm (131.)}$$
And if, by the help of this pair, or in any other way, the form
(130.) of this {\it principal function\/} $S_1$ had been found,
the integral equations (127.) and (128.) might have been deduced
from it, by our general method, as follows:
$$\left. \eqalign{
\varpi_1 &= {\delta S_1 \over \delta \eta_1}
            = {\eta_1 - e_1 \over t},\cr
\varpi_2 &= {\delta S_1 \over \delta \eta_2}
            = {\eta_2 - e_2 \over t},\cr
\varpi_3 &= {\delta S_1 \over \delta \eta_3}
            = {\eta_3 - e_3 \over t} - {\textstyle {1 \over 2}} gt,\cr}
   \right\}
   \eqno {\rm (132.)}$$
and
$$\left. \eqalign{
p_1 &= - {\delta S_1 \over \delta e_1}
            = {\eta_1 - e_1 \over t},\cr
p_2 &= - {\delta S_1 \over \delta e_2}
            = {\eta_2 - e_2 \over t},\cr
p_3 &= - {\delta S_1 \over \delta e_3}
            = {\eta_3 - e_3 \over t} + {\textstyle {1 \over 2}} gt,\cr}
   \right\}
   \eqno {\rm (133.)}$$
the latter of these two sets coinciding with (127.), and the former
set conducting to (128.).

\bigbreak

26.
Returning now from this simpler motion to the more complex motion
first mentioned, and denoting by $S_2$ the {\it disturbing part\/}
or function which must be added to $S_1$ in order to make up the
whole principal function $S$ of that more complex motion; we have,
by applying our general method, the following rigorous expression
for this disturbing function,
$$S_2 = - \int_0^t H_2 \,dt + \int_0^t {1 \over 2} \left\{
            \left( {\delta S_2 \over \delta \eta_1} \right)^2
          + \left( {\delta S_2 \over \delta \eta_2} \right)^2
          + \left( {\delta S_2 \over \delta \eta_3} \right)^2
            \right\} dt,
   \eqno {\rm (134.)}$$
in which we may, approximately, neglect the second definite
integral, and calculate the first by the help of the equations
of undisturbed motion.  In this manner we find, approximately,
by (125.), (127.),
$$-H_2 = - {\mu^2 \over 2} \{ (e_1 + p_1 t)^2 + (e_2 + p_2 t)^2 \}
         - {\nu^2 \over 2} (e_3 + p_3 t
            - {\textstyle {1 \over 2}} g t^2)^2,
   \eqno {\rm (135.)}$$
and therefore, by integration,
$$\left. \eqalign{S_2
   &= - {\textstyle {1 \over 2}} \{
            \mu^2 (e_1^2 + e_2^2) + \nu^2 e_3^2 \} t
      - {\textstyle {1 \over 2}} \{
            \mu^2 (e_1 p_1 + e_2 p_2) + \nu^2 e_3 p_3 \} t^2 \cr
   &\mathrel{\phantom{=}} \mathord{}
      - {\textstyle {1 \over 6}} \{
            \mu^2 (p_1^2 + p_2^2) + \nu^2 (p_3^2 - g e_3) \} t^3
      + {\textstyle {1 \over 8}} \nu^2 g p_3 t^4
      - {\textstyle {1 \over 40}} \nu^2 g^2 t^5,\cr}
   \right\}
   \eqno {\rm (136.)}$$
or, by (133.),
$$\left. \eqalign{S_2
   &= - {\mu^2 t \over 6} (\eta_1^2 + e_1 \eta_1+ e_1^2
            + \eta_2^2 + e_2 \eta_2 + e_2^2) \cr
   &\mathrel{\phantom{=}} \mathord{}
      - {\nu^2 t \over 6} \{ \eta_3^2 + e_3 \eta_3+ e_3^2
            + {\textstyle {1 \over 4}} g (\eta_3 + e_3) t^2
            + {\textstyle {1 \over 40}} g^2 t^4 \}:\cr}
   \right\}
   \eqno {\rm (137.)}$$
the error being of the fourth order, with respect to the small
quantities $\mu$, $\nu$.  And neglecting this small error, we
can deduce, by our general method, approximate forms for the
integrals of the equations of disturbed motion, from the
corrected function $S_1 + S_2$, as follows:
$$\left. \eqalign{
\varpi_1 &= {\delta S_1 \over \delta \eta_1}
                + {\delta S_2 \over \delta \eta_1}
            = {\eta_1 - e_1 \over t} - {\mu^2 t \over 3} (\eta_1
                + {\textstyle {1 \over 2}} e_1),\cr
\varpi_2 &= {\delta S_1 \over \delta \eta_2}
                + {\delta S_2 \over \delta \eta_2}
            = {\eta_2 - e_2 \over t} - {\mu^2 t \over 3} (\eta_2
                + {\textstyle {1 \over 2}} e_2),\cr
\varpi_3 &= {\delta S_1 \over \delta \eta_3}
                + {\delta S_2 \over \delta \eta_3}
            = {\eta_3 - e_3 \over t}
                - {\textstyle {1 \over 2}} gt
                - {\nu^2 t \over 3} (\eta_3
                + {\textstyle {1 \over 2}} e_3
                + {\textstyle {1 \over 8}} g t^2);\cr}
   \right\}
   \eqno {\rm (138.)}$$
and
$$\left. \eqalign{
p_1 &= - {\delta S_1 \over \delta e_1}
                - {\delta S_2 \over \delta e_1}
            = {\eta_1 - e_1 \over t} + {\mu^2 t \over 3} (e_1
                + {\textstyle {1 \over 2}} \eta_1),\cr
p_2 &= - {\delta S_1 \over \delta e_2}
                - {\delta S_2 \over \delta e_2}
            = {\eta_2 - e_2 \over t} + {\mu^2 t \over 3} (e_2
                + {\textstyle {1 \over 2}} \eta_2),\cr
p_3 &= - {\delta S_1 \over \delta e_3}
                - {\delta S_2 \over \delta e_3}
            = {\eta_3 - e_3 \over t}
                + {\textstyle {1 \over 2}} gt
                + {\nu^2 t \over 3} (e_3
                + {\textstyle {1 \over 2}} \eta_3
                + {\textstyle {1 \over 8}} g t^2);\cr}
   \right\}
   \eqno {\rm (139.)}$$
or, in the same order of approximation,
$$\left. \eqalign{
\eta_1 &= e_1 + p_1 t - {\textstyle {1 \over 2}} \mu^2 t^2
            (e_1 + {\textstyle {1 \over 3}} p_1 t),\cr
\eta_2 &= e_2 + p_2 t - {\textstyle {1 \over 2}} \mu^2 t^2
            (e_2 + {\textstyle {1 \over 3}} p_2 t),\cr
\eta_3 &= e_3 + p_3 t
            - {\textstyle {1 \over 2}} g t^2
            - {\textstyle {1 \over 2}} \nu^2 t^2
            (e_3 + {\textstyle {1 \over 3}} p_3 t
            - {\textstyle {1 \over 12}} g t^2),\cr}
   \right\}
   \eqno {\rm (140.)}$$
and
$$\left. \eqalign{
\varpi_1 &= p_1 - \mu^2 t (e_1 + {\textstyle {1 \over 2}} p_1 t),\cr
\varpi_2 &= p_2 - \mu^2 t (e_2 + {\textstyle {1 \over 2}} p_2 t),\cr
\varpi_3 &= p_3  - gt - \nu^2 t (e_3
            + {\textstyle {1 \over 2}} p_3 t
            - {\textstyle {1 \over 6}} g t^2).\cr}
   \right\}
   \eqno {\rm (141.)}$$
Accordingly, if we develope the rigorous integrals of disturbed
motion, (113.) and (114.), as far as the squares (inclusive)
of the small quantities $\mu$ and $\nu$, we are conducted to
these approximate integrals; and if we develope the rigorous
expression (120.) for the principal function of such motion,
to the same degree of accuracy, we obtain the sum of the
two expressions (130.) and (137.).

\bigbreak

27.
To illustrate still further, in the present example, our general
method of successive approximation, let $S_3$ denote the small
unknown correction of the approximate expression (137.), so that
we shall now have, rigorously, for the present disturbed motion,
$$S = S_1 + S_2 + S_3,
   \eqno {\rm (142.)}$$
$S_1$ and $S_2$ being here determined rigorously by (130.)
and (137).  Then, substituting $S_1 + S_2$ for $S_1$ in the
general transformation (87.), we find, rigorously, in the
present question,
$$\left. \eqalign{S_3
   &= - \int_0^t {1 \over 2} \left\{
            \left( {\delta S_2 \over \delta \eta_1} \right)^2
          + \left( {\delta S_2 \over \delta \eta_2} \right)^2
          + \left( {\delta S_2 \over \delta \eta_3} \right)^2
            \right\} dt \cr
   &\mathrel{\phantom{=}} \mathord{}
      + \int_0^t {1 \over 2} \left\{
            \left( {\delta S_3 \over \delta \eta_1} \right)^2
          + \left( {\delta S_3 \over \delta \eta_2} \right)^2
          + \left( {\delta S_3 \over \delta \eta_3} \right)^2
            \right\} dt:\cr}
   \right\}
   \eqno {\rm (143.)}$$
and if we neglect only terms of the eighth and higher dimensions
with respect to the small quantities $\mu$, $\nu$, we may confine
ourselves to the first of these two definite integrals, and may
employ, in calculating it, the approximate expressions (140.)
for the coordinates of disturbed motion.  In this manner we
obtain the very approximate expression,
$$\left. \eqalign{S_3
   &= - {\mu^4 \over 18} \int_0^t t^2 \{
            (\eta_1 + {\textstyle {1 \over 2}} e_1)^2
          + (\eta_2 + {\textstyle {1 \over 2}} e_2)^2 \} dt \cr
   &\mathrel{\phantom{=}} \mathord{}
      - {\nu^4 \over 18} \int_0^t t^2
            (\eta_3 + {\textstyle {1 \over 2}} e_3
                    + {\textstyle {1 \over 8}} g t^2)^2 dt \cr
   &= - {\mu^4 t^3 \over 360}
            ( 4 \eta_1^2 + 7 \eta_1 e_1 + 4 e_1^2
            + 4 \eta_2^2 + 7 \eta_2 e_2 + 4 e_2^2 ) \cr
   &\mathrel{\phantom{=}} \mathord{}
      - {\nu^4 t^3 \over 360}
            ( 4 \eta_3^2 + 7 \eta_3 e_3 + 4 e_3^2 )
      - {\nu^4 g t^5 \over 240} (\eta_3 + e_3)
      - {17 \nu^4 g^2 t^7 \over 40320} \cr
   &\mathrel{\phantom{=}} \mathord{}
      - {\mu^6 t^5 \over 945}
            ( \eta_1^2 + {\textstyle {31 \over 16}} \eta_1 e_1 + e_1^2
            + \eta_2^2 + {\textstyle {31 \over 16}} \eta_2 e_2 + e_2^2 ) \cr
   &\mathrel{\phantom{=}} \mathord{}
      - {\nu^6 t^5 \over 945}
            ( \eta_3^2 + {\textstyle {31 \over 16}} \eta_3 e_3 + e_3^2 )
      - {17 \nu^6 g t^7 \over 40320} (\eta_3 + e_3)
      - {31 \nu^6 g^2 t^9 \over 725760};\cr}
   \right\}
   \eqno {\rm (144.)}$$
which is accordingly the sum of the tems of the fourth and sixth
dimensions in the development of the rigorous expression (120.),
and gives, by our general method, correspondingly approximate
expressions for the integrals of disturbed motion, under the forms
$$\left. \eqalign{
\varpi_1 &=   {\delta S_1 \over \delta \eta_1}
            + {\delta S_2 \over \delta \eta_1}
            + {\delta S_3 \over \delta \eta_1},\cr
\varpi_2 &=   {\delta S_1 \over \delta \eta_2}
            + {\delta S_2 \over \delta \eta_2}
            + {\delta S_3 \over \delta \eta_2},\cr
\varpi_3 &=   {\delta S_1 \over \delta \eta_3}
            + {\delta S_2 \over \delta \eta_3}
            + {\delta S_3 \over \delta \eta_3},\cr}
   \right\}
   \eqno {\rm (145.)}$$
and
$$\left. \eqalign{
p_1 &= - {\delta S_1 \over \delta e_1}
       - {\delta S_2 \over \delta e_1}
       - {\delta S_3 \over \delta e_1},\cr
p_2 &= - {\delta S_1 \over \delta e_2}
       - {\delta S_2 \over \delta e_2}
       - {\delta S_3 \over \delta e_2},\cr
p_3 &= - {\delta S_1 \over \delta e_3}
       - {\delta S_2 \over \delta e_3}
       - {\delta S_3 \over \delta e_3}.\cr}
   \right\}
   \eqno {\rm (146.)}$$

\medskip

28.
To illustrate by the same example the theory of gradually varying
elements, let us establish the following definitions, for the
present disturbed motion,
$$\left. \multieqalign{
\kappa_1 &= \eta_1 - \varpi_1 t, &
\kappa_2 &= \eta_2 - \varpi_2 t, &
\kappa_3 &= \eta_3 - \varpi_3 t - {\textstyle {1 \over 2}} g t^2,\cr
\lambda_1 &= \varpi_1, &
\lambda_2 &= \varpi_2, &
\lambda_3 &= \varpi_3 + gt,\cr}
   \right\}
   \eqno {\rm (147.)}$$
and let us call these six quantities
$\kappa_1$~$\kappa_2$~$\kappa_3$
$\lambda_1$~$\lambda_2$~$\lambda_3$
the {\it varying elements\/} of that motion, by analogy to the
six constant quantities $e_1$~$e_2$~$e_3$ $p_1$~$p_2$~$p_3$,
which may, for the undisturbed motion, be represented in a similar
way, namely, by (127.) and (128.),
$$\left. \multieqalign{
e_1 &= \eta_1 - \varpi_1 t, &
e_2 &= \eta_2 - \varpi_2 t, &
e_3 &= \eta_3 - \varpi_3 t - {\textstyle {1 \over 2}} g t^2,\cr
p_1 &= \varpi_1, &
p_2 &= \varpi_2, &
p_3 &= \varpi_3 + gt.\cr}
   \right\}
   \eqno {\rm (148.)}$$
We shall then have rigorously, for the six disturbed variables
$\eta_1$~$\eta_2$~$\eta_3$ $\varpi_1$~$\varpi_2$~$\varpi_3$,
expressions of the same forms as in the integrals (127.) and (128.)
of undisturbed motion, but with variable instead of constant elements,
namely, the following:
$$\left. \multieqalign{
\eta_1 &= \kappa_1 + \lambda_1 t, &
\eta_2 &= \kappa_2 + \lambda_2 t, &
\eta_3 &= \kappa_3 + \lambda_3 t - {\textstyle {1 \over 2}} g t^2,\cr
\varpi_1 &= \lambda_1, &
\varpi_2 &= \lambda_2, &
\varpi_3 &= \lambda_3 - gt;\cr}
   \right\}
   \eqno {\rm (149.)}$$
and the rigorous determination of the six varying elements
$\kappa_1$~$\kappa_2$~$\kappa_3$
$\lambda_1$~$\lambda_2$~$\lambda_3$,
as functions of the time and of their own initial values
$e_1$~$e_2$~$e_3$ $p_1$~$p_2$~$p_3$,
depends on the integration of the 6 following equations,
in ordinary differentials of the first order, of the forms (105.):
$$\left. \eqalign{
{d\kappa_1 \over dt} &= {\delta H_2 \over \delta \lambda_1}
      = \mu^2 t (\kappa_1 + \lambda_1 t),\cr
{d\kappa_2 \over dt} &= {\delta H_2 \over \delta \lambda_2}
      = \mu^2 t (\kappa_2 + \lambda_2 t),\cr
{d\kappa_3 \over dt} &= {\delta H_2 \over \delta \lambda_3}
      = \nu^2 t (\kappa_3 + \lambda_3 t
            - {\textstyle {1 \over 2}} g t^2),\cr}
   \right\}
   \eqno {\rm (150.)}$$
and
$$\left. \eqalign{
{d\lambda_1 \over dt} &= - {\delta H_2 \over \delta \kappa_1}
      = - \mu^2 (\kappa_1 + \lambda_1 t),\cr
{d\lambda_2 \over dt} &= - {\delta H_2 \over \delta \kappa_2}
      = - \mu^2 (\kappa_2 + \lambda_2 t),\cr
{d\lambda_3 \over dt} &= - {\delta H_2 \over \delta \kappa_3}
      = - \nu^2 (\kappa_3 + \lambda_3 t
            - {\textstyle {1 \over 2}} g t^2),\cr}
   \right\}
   \eqno {\rm (151.)}$$
$H_2$ being here the expression
$$H_2 = {\mu^2 \over 2} \{
           (\kappa_1 + \lambda_1 t)^2 + (\kappa_2 + \lambda_2 t)^2 \}
         + {\nu^2 \over 2}
           (\kappa_3 + \lambda_3 t - {\textstyle {1 \over 2}} g t^2)^2,
   \eqno {\rm (152.)}$$
which is obtained from (125.) by substituting for the disturbed
coordinates $\eta_1$~$\eta_2$~$\eta_3$ their values (149.),
as functions of the varying elements and of the time.  It is
not difficult to integrate rigorously this system of equations
(150.) and (151.); and we shall soon have occasion to state their
complete and accurate integrals: but we shall continue for a while
to treat these rigorous integrals as unknown, that we may take
this oppportunity to exemplify our general method of indefinite
approximation, for all such dynamical systems, founded on the
properties of the {\it functions of elements\/} $C$ and $E$.
Of these two functions either may be employed, and we shall
use here the function $C$.

\bigbreak

29.
This function, by (109.) and (152.), may rigorously be expressed
as follows:
$$C = {\mu^2 \over 2} \int_0^t
         ( \lambda_1^2 t^2 - \kappa_1^2
         + \lambda_2^2 t^2 - \kappa_2^2 ) \, dt
      + {\nu^2 \over 2} \int_0^t
         \{ ( \lambda_3 t - {\textstyle {1 \over 2}} g t^2)^2
            - \kappa_3^2 \} \, dt;
   \eqno {\rm (153.)}$$
and has therefore the following for a first approximate value,
obtained by treating the elements
$\kappa_1$~$\kappa_2$~$\kappa_3$
$\lambda_1$~$\lambda_2$~$\lambda_3$
as constant and equal to their initial values
$e_1$~$e_2$~$e_3$ $p_1$~$p_2$~$p_3$,
$$\left. \eqalign{C
   &= - {t \over 2} \{ \mu^2 (e_1^2 + e_2^2) + \nu^2 e_3^2 \}
      + {t^3 \over 6} \{ \mu^2 (p_1^2 + p_2^2) + \nu^2 p_3^2 \} \cr
   &\mathrel{\phantom{=}} \mathord{}
      - {t^4 \over 8} \nu^2 g p_3
      + {t^5 \over 40} \nu^2 g^2.\cr}
   \right\}
   \eqno {\rm (154.)}$$
In like manner we have, as first approximations, of the kind
expressed by the general formula (Z${}^1$.), the following
results deduced from the equations (151.),
$$\left. \eqalign{
\lambda_1 &= p_1 - \mu^2 (e_1 t + {\textstyle {1 \over 2}} p_1 t^2 ),\cr
\lambda_2 &= p_2 - \mu^2 (e_2 t + {\textstyle {1 \over 2}} p_2 t^2 ),\cr
\lambda_3 &= p_3 - \nu^2 (e_3 t + {\textstyle {1 \over 2}} p_3 t^2
               - {\textstyle {1 \over 6}} g t^3),\cr}
   \right\}
   \eqno {\rm (155.)}$$
and therefore, as approximations of the same kind,
$$\left. \eqalign{
e_1 &= - {\textstyle {1 \over 2}} p_1 t
         - {\lambda_1 - p_1 \over \mu^2 t},\cr
e_2 &= - {\textstyle {1 \over 2}} p_2 t
         - {\lambda_2 - p_2 \over \mu^2 t},\cr
e_3 &= - {\textstyle {1 \over 2}} p_3 t + {\textstyle {1 \over 6}} g t^2
         - {\lambda_3 - p_3 \over \nu^2 t}.\cr}
   \right\}
   \eqno {\rm (156.)}$$
Substituting these values for the initial constants $e_1$~$e_2$~$e_3$
in the approximate value (154.) for the function of elements $C$,
we obtain the following approximate expression $C_1$ for that
function, of the form supposed by our theory:
$$\left. \eqalign{C_1
   &= - {1\over 2t} \left\{
         {(\lambda_1 - p_1)^2 + (\lambda_2 - p_2)^2 \over \mu^2}
         + {(\lambda_3 - p_3)^2 \over \nu^2} \right\} \cr
   &\mathrel{\phantom{=}} \mathord{}
      - {t \over 2} \{ (\lambda_1 - p_1) p_1 + (\lambda_2 - p_2) p_2
            + (\lambda_3 - p_3) (p_3  - {\textstyle {1 \over 3}} gt) \} \cr
   &\mathrel{\phantom{=}} \mathord{}
      + {t^3 \over 24} \{ \mu^2 (p_1^2 + p_2^2) + \nu^2 p_3^2 \}
      - {t^4 \over 24} \nu^2 g p_3 + {t^5 \over 90} \nu^2 g^2.\cr}
   \right\}
   \eqno {\rm (157.)}$$
The rigorous function~$C$ must satisfy, in the present question,
by the principles of the eighteenth number, the partial differential
equation,
$${\delta C \over \delta t}
   = {\mu^2 \over 2} \left\{
         \left( {\delta C \over \delta \lambda_1} + \lambda_1 t \right)^2
       + \left( {\delta C \over \delta \lambda_2} + \lambda_2 t \right)^2
         \right\}
   + {\nu^2 \over 2}
         \left( {\delta C \over \delta \lambda_3} + \lambda_3 t
       - {\textstyle {1 \over 2}} g t^2 \right)^2;
   \eqno {\rm (158.)}$$
and if it be put under the form (U${}^1$.),
$$C = C_1 + C_2,$$
$C_1$ being a first approximation, supposed to vanish with the time,
then the correction~$C_2$ must satisfy rigorously the condition
$$\left. \eqalign{C_2
  &= \int_0^t \biggl\{ - {\delta C_1 \over \delta t}
   + {\mu^2 \over 2}
         \left( {\delta C_1 \over \delta \lambda_1} + \lambda_1 t \right)^2
       + {\mu^2 \over 2}
         \left( {\delta C_1 \over \delta \lambda_2} + \lambda_2 t \right)^2 \cr
  &\hskip 6em
   + {\nu^2 \over 2}
         \left( {\delta C_1 \over \delta \lambda_3} + \lambda_3 t
       - {\textstyle {1 \over 2}} g t^2 \right)^2 \biggr\} \, dt \cr
   &\mathrel{\phantom{=}} \mathord{}
   - {1 \over 2} \int_0^t \left\{
         \mu^2 \left( {\delta C_2 \over \delta \lambda_1} \right)^2
       + \mu^2 \left( {\delta C_2 \over \delta \lambda_2} \right)^2
       + \nu^2 \left( {\delta C_2 \over \delta \lambda_3} \right)^2
         \right\} dt.\cr}
   \right\}
   \eqno {\rm (159.)}$$

In passing to a second approximation we may neglect the second
definite integral, and may calculate the first with the help of the
approximate equations (155.); which give, in this manner,
$$\left. \eqalign{C_2
   &= - \int_0^t \{ (\lambda_1 - p_1)^2 + (\lambda_2 - p_2)^2
                  + (\lambda_3 - p_3)^2 \} \, dt \cr
   &\mathrel{\phantom{=}} \mathord{}
      + {\mu^2 \over 2} \int_0^t \{ \lambda_1 (\lambda_1 - p_1)
         + \lambda_2 (\lambda_2 - p_2) \} \, t^2 \, dt \cr
   &\mathrel{\phantom{=}} \mathord{}
      + {\nu^2 \over 2} \int_0^t
         (\lambda_3 - {\textstyle {2 \over 3}} gt)
         (\lambda_3 - p_3) \, t^2 \,dt \cr
   &= - {t \over 3} \{ (\lambda_1 - p_1)^2 + (\lambda_2 - p_2)^2
                     + (\lambda_3 - p_3)^2 \} \cr
   &\mathrel{\phantom{=}} \mathord{}
      + {t^3 \over 24} \{ \mu^2 p_1 (\lambda_1 - p_1)
                        + \mu^2 p_2 (\lambda_2 - p_2)
                        + \nu^2 p_3 (\lambda_3 - p_3) \} \cr
   &\mathrel{\phantom{=}} \mathord{}
      - {t^4 \over 45} \nu^2 g (\lambda_3 - p_3)
      + {t^5 \over 240} (\mu^4 p_1^2 + \mu^4 p_2^2 + \nu^4 p_3^2) \cr
   &\mathrel{\phantom{=}} \mathord{}
      - {t^6 \over 240} \nu^4 g p_3 + {t^7 \over 945} \nu^4 g^2.\cr}
   \right\}
   \eqno {\rm (160.)}$$
We might improve this second approximation in like manner, by
calculating a new definite integral $C_3$, with the help of the
following more approximate forms for the relations between the
varying elements $\lambda_1$~$\lambda_2$~$\lambda_3$ and the
initial constants, deduced by our general method:
$$\left. \eqalign{
e_1  =  - {\delta C_1 \over \delta p_1} - {\delta C_2 \over \delta p_1}
    &=  - {\lambda_1 - p_1 \over \mu^2 t}  \left(
            1 + {\mu^2 t^2 \over 6} + {\mu^4 t^4 \over 24} \right)
        - {t p_1 \over 2} \left(
            1 + {\mu^2 t^2 \over 12} + {\mu^4 t^4 \over 60} \right),\cr
e_2  =  - {\delta C_1 \over \delta p_2} - {\delta C_2 \over \delta p_2}
    &=  - {\lambda_2 - p_2 \over \mu^2 t}  \left(
            1 + {\mu^2 t^2 \over 6} + {\mu^4 t^4 \over 24} \right)
        - {t p_2 \over 2} \left(
            1 + {\mu^2 t^2 \over 12} + {\mu^4 t^4 \over 60} \right),\cr
e_3  =  - {\delta C_1 \over \delta p_3} - {\delta C_2 \over \delta p_3}
    &=  - {\lambda_3 - p_3 \over \nu^2 t}  \left(
            1 + {\nu^2 t^2 \over 6} + {\nu^4 t^4 \over 24} \right)
        - {t p_3 \over 2} \left(
            1 + {\nu^2 t^2 \over 12} + {\nu^4 t^4 \over 60} \right) \cr
   &\mathrel{\phantom{=}} \mathord{}
        + {g t^2 \over 6} \left(
            1 + {7 \nu^2 t^2 \over 60} + {\nu^4 t^4 \over 40} \right);\cr}
   \right\}
   \eqno {\rm (161.)}$$
in which we can only depend on the terms as far as the second order,
but which acquire a correctness of the fourth order when cleared of
the small divisors, and give then
$$\left. \eqalign{
\lambda_1 &= p_1 - \mu^2 t (e_1 + {\textstyle {1 \over 2}} p_1 t)
            + {\textstyle {1 \over 6}} \mu^4 t^3
              (e_1 + {\textstyle {1 \over 4}} p_1 t),\cr
\lambda_2 &= p_2 - \mu^2 t (e_2 + {\textstyle {1 \over 2}} p_2 t)
            + {\textstyle {1 \over 6}} \mu^4 t^3
              (e_2 + {\textstyle {1 \over 4}} p_2 t),\cr
\lambda_3 &= p_3 - \nu^2 t (e_3 + {\textstyle {1 \over 2}} p_3 t
                  - {\textstyle {1 \over 6}} g t^2)
            + {\textstyle {1 \over 6}} \nu^4 t^3
              (e_3 + {\textstyle {1 \over 4}} p_3 t
                  - {\textstyle {1 \over 20}} g t^2).\cr}
   \right\}
   \eqno {\rm (162.)}$$
But a little attention to the nature of this process shows that
all the successive corrections to which it conducts can be only
rational and integer and homogeneous functions, of the second
dimension, of the quantities
$\lambda_1$~$\lambda_2$~$\lambda_3$ $p_1$~$p_2$~$p_3$~$g$,
and that they may all be put under the following form, which is
therefore the form of their sum, or of the whole sought function~$C$;
$$\left. \eqalign{C
   &= \mathbin{\phantom{+}}
        \mu^{-2} a_\mu (\lambda_1 - p_1)^2
      + b_\mu p_1 (\lambda_1 - p_1) + \mu^2 c_\mu p_1^2 \cr
   &\mathrel{\phantom{=}} \mathord{}
      + \mu^{-2} a_\mu (\lambda_2 - p_2)^2
      + b_\mu p_2 (\lambda_2 - p_2) + \mu^2 c_\mu p_2^2 \cr
   &\mathrel{\phantom{=}} \mathord{}
      + \nu^{-2} a_\nu (\lambda_3 - p_3)^2
      + b_\nu p_3 (\lambda_3 - p_3) + \nu^2 c_\nu p_3^2 \cr
   &\mathrel{\phantom{=}} \mathord{}
      + f_\nu g (\lambda_3 - p_3) + \nu^2 h_\nu g p_3
      + \nu^2 i_\nu g^2,\cr}
   \right\}
   \eqno {\rm (163.)}$$
the cofficients $a_\mu$~$a_\nu$, \&c. being functions of the small
quantities $\mu$, $\nu$, and also of the time, of which it remains
to discover the forms.  Denoting therefore their differentials,
taken with respect to the time, as follows,
$$d a_\mu = a_\mu' \,dt,\quad d a_\nu = a_\nu' dt, \hbox{ \&c.,}
   \eqno {\rm (164.)}$$
and substituting the expression (163.) in the rigorous partial
differential equation (158.), we are conducted to the six
following equations in ordinary differentials of the first order:
$$\left. \multieqalign{
2 a_\nu' &= (2 a_\nu + \nu^2 t)^2; &
  b_\nu' &= (2 a_\nu + \nu^2 t)(b_\nu + t); &
  c_\nu' &= {\textstyle {1 \over 2}} (b_\nu + t)^2;\cr
  f_\nu' &= (2 a_\nu + \nu^2 t)(f_\nu - {\textstyle {1 \over 2}} t^2); &
  h_\nu' &= (b_\nu + t)(f_\nu - {\textstyle {1 \over 2}} t^2); &
  i_\nu' &= {\textstyle {1 \over 2}}
               (f_\nu - {\textstyle {1 \over 2}} t^2)^2;\cr}
   \right\}
   \eqno {\rm (165.)}$$
along with the 6 following conditions, to determine the 6
arbitrary constants introduced by integration,
$$a_0 = - {1 \over 2t};\quad
  b_0 = - {t \over 2};\quad
  f_0 =   {t^2 \over 6};\quad
  c_0 =   {t^3 \over 24};\quad
  h_0 = - {t^4 \over 24};\quad
  i_0 =   {t^5 \over 90}.\quad
   \eqno {\rm (166.)}$$

In this manner we find, without difficulty, observing that
$a_\mu$~$b_\mu$~$c_\mu$ may be formed from
$a_\nu$~$b_\nu$~$c_\nu$ by changing $\nu$ to $\mu$,
$$\left. \multieqalign{
a_\nu &= - {\textstyle {1 \over 2}} \nu^2 t
         - {\textstyle {1 \over 2}} \nu \mathop{\rm cotan} \nu t, &
a_\mu &= - {\textstyle {1 \over 2}} \mu^2 t
         - {\textstyle {1 \over 2}} \mu \mathop{\rm cotan} \mu t, \cr
b_\nu &= - t + {1 \over \nu} \tan {\nu t \over 2}, &
b_\mu &= - t + {1 \over \mu} \tan {\mu t \over 2}, \cr
c_\nu &= - {t \over 2 \nu^2} + {1 \over \nu^3} \tan {\nu t \over 2}, &
c_\mu &= - {t \over 2 \mu^2} + {1 \over \mu^3} \tan {\mu t \over 2}, \cr
f_\nu &= {\textstyle {1 \over 2}} t^2 - {1 \over \nu^2}
         + {t \over \nu} \mathop{\rm cotan} \nu t,\cr
h_\nu &= {t^2 \over 2 \nu^2} - {t \over \nu^3} \tan {\nu t \over 2}, \cr
i_\nu &= {t \over 2 \nu^4} - {t^3 \over 6 \nu^2}
         - {t^2 \over 2 \nu^3} \mathop{\rm cotan} \nu t.\cr}
   \right\}
   \eqno {\rm (167.)}$$

The form of the function~$C$ is therefore entirely known,
and we have for this {\it function of elements\/} the following
rigorous expression,
$$\left. \eqalign{C
   &= - {(\lambda_1 - p_1)^2 + (\lambda_2 - p_2)^2 \over 2\mu \tan \mu t}
      - {(\lambda_3 - p_3)^2 \over 2\nu \tan \nu t} \cr
   &\mathrel{\phantom{=}} \mathord{}
      - {t \over 2} \{ (\lambda_1 - p_1)^2 + (\lambda_2 - p_2)^2
                     + (\lambda_3 - p_3)^2 \} \cr
   &\mathrel{\phantom{=}} \mathord{}
      - t \{ p_1 (\lambda_1 - p_1) + p_2 (\lambda_2 - p_2)
           + p_3 (\lambda_3 - p_3) \} \cr
   &\mathrel{\phantom{=}} \mathord{}
      + {1 \over \mu} \{ p_1 (\lambda_1 - p_1)
                       + p_2 (\lambda_2 - p_2) \} \tan {\mu t \over 2}
      + {1 \over \nu}    p_3 (\lambda_3 - p_3)    \tan {\nu t \over 2} \cr
   &\mathrel{\phantom{=}} \mathord{}
      - {t \over 2} (p_1^2 + p_2 ^2 + p_3^2)
      + {1 \over \mu} (p_1^2 + p_2^2) \tan {\mu t \over 2}
      + {1 \over \nu} p_3^2 \tan {\nu t \over 2} \cr
   &\mathrel{\phantom{=}} \mathord{}
      + \left( {t^2 \over 2} - {1 \over \nu^2}
            + {t \over \nu} \mathop{\rm cotan} \nu t \right)
            g (\lambda_3 - p_3)
      + \left( {t^2 \over 2} - {t \over \nu} \tan {\nu t \over 2} \right)
            g p_3 \cr
   &\mathrel{\phantom{=}} \mathord{}
      + \left( {t \over 2\nu^2} - {t^2 \over 6}
            - {t^2 \over 2\nu} \mathop{\rm cotan} \nu t \right) g^2,\cr}
   \right\}
   \eqno {\rm (168.)}$$
which may be variously transformed, and gives by our general method
the following systems of rigorous integrals of the differential
equations of varying elements, (150.), (151.):
$$\left. \eqalign{
e_1 = - {\delta C \over \delta p_1}
   &= - {\lambda_1 - p_1 \over \mu \sin \mu t}
           - {p_1 \over \mu} \tan {\mu t \over 2},\cr
e_2 = - {\delta C \over \delta p_2}
   &= - {\lambda_2 - p_2 \over \mu \sin \mu t}
           - {p_2 \over \mu} \tan {\mu t \over 2},\cr
e_3 = - {\delta C \over \delta p_3}
   &= - {\lambda_3 - p_3 \over \nu \sin \nu t}
      - {p_3 \over \nu} \tan {\nu t \over 2}
      + {g \over \nu} \left( {t \over \sin \nu t}
         - {1 \over \nu} \right),\cr}
   \right\}
   \eqno {\rm (169.)}$$
and
$$\left. \eqalign{
\kappa_1 = {\delta C \over \delta \lambda_1}
   &= - (\lambda_1 - p_1) \left(
            t + {1 \over \mu} \mathop{\rm cotan} \mu t \right)
      + p_1 \left( -t + {1 \over \mu} \tan {\mu t \over 2} \right),\cr
\kappa_2 = {\delta C \over \delta \lambda_2}
   &= - (\lambda_2 - p_2) \left(
            t + {1 \over \mu} \mathop{\rm cotan} \mu t \right)
      + p_2 \left( -t + {1 \over \mu} \tan {\mu t \over 2} \right),\cr
\kappa_3 = {\delta C \over \delta \lambda_3}
   &= - (\lambda_3 - p_3) \left(
            t + {1 \over \nu} \mathop{\rm cotan} \nu t \right)
      + p_3 \left( -t + {1 \over \nu} \tan {\nu t \over 2} \right) \cr
   &\mathrel{\phantom{=}} \mathord{}
      + g \left( {t^2 \over 2} - {1 \over \nu^2}
            + {t \over \nu} \mathop{\rm cotan} \nu t \right);\cr}
   \right\}
   \eqno {\rm (170.)}$$
that is,
$$\left. \eqalign{
\lambda_1 &= p_1 \cos \mu t - e_1 \mu \sin \mu t,\cr
\lambda_2 &= p_2 \cos \mu t - e_2 \mu \sin \mu t,\cr
\lambda_3 &= p_3 \cos \nu t - e_3 \nu \sin \nu t
      + g \left( t - {1 \over \nu} \sin \nu t \right),\cr}
   \right\}
   \eqno {\rm (171.)}$$
and
$$\left. \eqalign{
\kappa_1 &= e_1 (\cos \mu t + \mu t \sin \mu t)
      + p_1 \left( {1 \over \mu} \sin \mu t - t \cos \mu t \right),\cr
\kappa_2 &= e_2 (\cos \mu t + \mu t \sin \mu t)
      + p_2 \left( {1 \over \mu} \sin \mu t - t \cos \mu t \right),\cr
\kappa_3 &= e_3 (\cos \nu t + \nu t \sin \nu t)
      + p_3 \left( {1 \over \nu} \sin \nu t - t \cos \nu t \right) \cr
   &\mathrel{\phantom{=}} \mathord{}
      - g \left( {\mathop{\rm vers} \nu t \over \nu^2}
         - {t \over \nu} \sin \nu t + {t^2 \over 2} \right).\cr}
   \right\}
   \eqno {\rm (172.)}$$

Accordingly, these rigorous expressions for the 6 varying elements,
in the present dynamical question, agree with the results obtained
by the ordinary methods of integration from the 6 ordinary
differential equations (150.) and (151.), and with those obtained
by elimination from the equations (113.), (114.), (147.).

\bigbreak

{\sectiontitle
Remarks on the foregoing Example.\par}

\nobreak\bigskip

30.
The example which has occupied us in the last six numbers is not
altogether ideal, but is realised to some extent by the motion
of a projectile in a void.  For if we consider the earth as a
sphere, of radius~$R$, and suppose the accelerating force of
gravity to vary inversely as the square of the distance~$r$
from the centre, and to be $= g$ at the surface, this force will
be represented generally by
$\displaystyle {g R^2 \over r^2}$;
and to adapt the differential equations (78.) to the motion
of a projectile in a void, it will be sufficient to make
$$U = g R^2 \left( {1 \over r} - {1 \over R} \right)
   \eqno {\rm (173.)}$$

If we place the origin of rectangular coordinates at the earth's
surface, and suppose the semiaxis of $+z$ to be directed vertically
upwards, we shall have
$$r = \sqrt{(R + z)^2 + x^2 + y^2},
   \eqno {\rm (174.)}$$
and
$$U = -gz + {gz^2 \over R} - {g(x^2 + y^2) \over 2R},
   \eqno {\rm (175.)}$$
neglecting only those very small terms which have the square of the
earth's radius for a divisor: neglecting therefore such terms, the
force-function $U$ in this question is of that form (110.) on which
all the reasonings of the example have been founded; the small
constants $\mu$, $\nu$, being the real and imaginary quantities
$\displaystyle \sqrt{ {g \over R} }$,
$\displaystyle \sqrt{ {-2g \over R} }$,
respectively.  We may therefore apply the results of the recent
numbers to the motions of projectiles in a void, by substituting
these values for the constants, and altering, where necessary,
trigonometric to exponential functions.  But besides the theoretical
facility and the little practical importance of researches respecting
such projectiles, the results would only be accurate as far as the
first negative power (inclusive) of the earth's radius, because the
expression (110.) for the force-function $U$ is only accurate
only so far; and therefore the rigorous and approximate investigations
of the six preceding numbers, founded on that expression, are offered
only as mathematical illustrations of a general {\it method},
extending to all problems of dynamics, at least to all those
to which the law of living forces applies.

\bigbreak

{\sectiontitle
Attracting Systems resumed: Differential Equations of internal
or Relative Motion; Integration by the Principal Function.\par}

\nobreak\bigskip

31.
Returning now from this digression on the motion of a single point,
to the more important study of an attracting or repelling system,
let us resume the differential equations (A.), which may be
thus summed up:
$$dt \, \delta H
   = \sum (d\eta \, \delta \varpi - d \varpi \delta \eta);
   \eqno {\rm (A^2.)}$$
and in order to separate the absolute motion of the whole system
in space from the motions of its points among themselves,
let us choose the following marks of position:
$$x_{\prime\prime} = {\sum \mathbin{.} mx \over \sum m},\quad
  y_{\prime\prime} = {\sum \mathbin{.} my \over \sum m},\quad
  z_{\prime\prime} = {\sum \mathbin{.} mz \over \sum m},
   \eqno {\rm (176.)}$$
and
$$\xi_i   = x_i - x_n,\quad
  \eta_i  = y_i - y_n,\quad
  \zeta_i = z_i - z_n;
   \eqno {\rm (177.)}$$
that is, the 3 rectangular coordinates of the centre of gravity
of the system, referred to an origin fixed in space, and the
$3n - 3$ rectangular coordinates of the $n - 1$ masses
$m_1, m_2,\ldots\, m_{n-1}$, referred to the $n$th mass $m_n$,
as an internal and moveable origin, but to axes parallel to
the former.  We then find, as in the former Essay,
$$\left. \eqalign{T
   &= {\textstyle {1 \over 2}}
         (x_{\prime\prime}'^2 + y_{\prime\prime}'^2 + z_{\prime\prime}'^2)
         \sum m
      + {\textstyle {1 \over 2}}
         \sum\nolimits_\prime \mathbin{.} m (\xi'^2 + \eta'^2 + \zeta'^2) \cr
   &\mathrel{\phantom{=}} \mathord{}
      - {1 \over 2 \sum m}
            \left\{
               \left( \sum\nolimits_\prime \mathbin{.} m\xi' \right)^2
             + \left( \sum\nolimits_\prime \mathbin{.} m\eta' \right)^2
             + \left( \sum\nolimits_\prime \mathbin{.} m\zeta' \right)^2
            \right\},\cr}
   \right\}
   \eqno {\rm (178.)}$$
the sign of summation $\sum\nolimits_\prime$ referring to the first
$n - 1$ masses only; and therefore,
$$\left. \eqalign{T
   &= {1 \over 2 \sum m}
         \left\{
            \left( {\delta T \over \delta x_{\prime\prime}'} \right)^2
          + \left( {\delta T \over \delta y_{\prime\prime}'} \right)^2
          + \left( {\delta T \over \delta z_{\prime\prime}'} \right)^2
         \right\} \cr
   &\mathrel{\phantom{=}} \mathord{}
      + {\textstyle {1 \over 2}} \sum\nolimits_\prime \mathbin{.} {1 \over m}
         \left\{
            \left( {\delta T \over \delta \xi'}   \right)^2
          + \left( {\delta T \over \delta \eta'}  \right)^2
          + \left( {\delta T \over \delta \zeta'} \right)^2
         \right\} \cr
   &\mathrel{\phantom{=}} \mathord{}
      + {1 \over 2 m_n} \left\{
           \left( \sum\nolimits_\prime
                  {\delta T \over \delta \xi'}   \right)^2
         + \left( \sum\nolimits_\prime
                  {\delta T \over \delta \eta'}  \right)^2
         + \left( \sum\nolimits_\prime
                  {\delta T \over \delta \zeta'} \right)^2
         \right\}.\cr}
   \right\}
   \eqno {\rm (179.)}$$

If then we put for abridgement,
$$\left. \eqalign{
x_\prime' &= {1 \over m} {\delta T \over \delta \xi'}
      = \xi' - {\sum\nolimits_\prime \mathbin{.} m \xi' \over \sum m},\cr
y_\prime' &= {1 \over m} {\delta T \over \delta \eta'}
      = \eta' - {\sum\nolimits_\prime \mathbin{.} m \eta' \over \sum m},\cr
z_\prime' &= {1 \over m} {\delta T \over \delta \zeta'}
      = \zeta' - {\sum\nolimits_\prime \mathbin{.} m \zeta' \over \sum m},\cr}
   \right\}
   \eqno {\rm (180.)}$$
we shall have the expression
$$\left. \eqalign{H
   &= {\textstyle {1 \over 2}} (x_{\prime\prime}'^2 + y_{\prime\prime}^2
            + z_{\prime\prime}'^2) \sum m
         + {\textstyle {1 \over 2}} \sum\nolimits_\prime \mathbin{.}
            m (x_\prime'^2 + y_\prime'^2 + z_\prime'^2) \cr
   &\mathrel{\phantom{=}} \mathord{}
      + {1 \over 2 m_n} \left\{
              \left( \sum\nolimits_\prime \mathbin{.} m x_\prime' \right)^2
            + \left( \sum\nolimits_\prime \mathbin{.} m y_\prime' \right)^2
            + \left( \sum\nolimits_\prime \mathbin{.} m z_\prime' \right)^2
            \right\} - U,\cr}
   \right\}
   \eqno {\rm (B^2.)}$$
of which the variation is to be compared with the following form
of (A${}^2$.),
$$\left. \eqalign{dt \, \delta H
   &= ( dx_{\prime\prime}  \, \delta x_{\prime\prime}'
      - dx_{\prime\prime}' \, \delta x_{\prime\prime}
      + dy_{\prime\prime}  \, \delta y_{\prime\prime}'
      - dy_{\prime\prime}' \, \delta y_{\prime\prime}
      + dz_{\prime\prime}  \, \delta z_{\prime\prime}'
      - dz_{\prime\prime}' \, \delta z_{\prime\prime} )
      \sum m \cr
   &\mathrel{\phantom{=}} \mathord{}
    + \sum\nolimits_\prime \mathbin{.} m
      ( d\xi       \, \delta x_\prime'
      - dx_\prime' \, \delta \xi
      + d\eta      \, \delta y_\prime'
      - dy_\prime' \, \delta \eta
      + d\zeta     \, \delta z_\prime'
      - dz_\prime' \, \delta \zeta ),\cr}
   \right\}
   \eqno {\rm (C^2.)}$$
in order to form, by our general process, $6n$ differential
equations of the first order, between the $6n$ quantities
$x_{\prime\prime}$~$y_{\prime\prime}$~$z_{\prime\prime}$
$x_{\prime\prime}'$~$y_{\prime\prime}'$~$z_{\prime\prime}'$
$\xi$~$\eta$~$\zeta$ $x_\prime'$~$y_\prime'$~$z_\prime'$
and the time~$t$.  In thus taking the variation of $H$, we are
to remember that the force-function $U$ depends only on the
$3n - 3$ internal coordinates $\xi$~$\eta$~$\zeta$, being of
the form
$$\left. \eqalign{U
   &= m_n ( m_1 f_1 + m_2 f_2 + \cdots + m_{n-1} f_{n-1} ) \cr
   &\mathrel{\phantom{=}} \mathord{}
    + m_1 m_2 f_{1,2} + m_1 m_3 f_{1,3} + \cdots
       + m_{n-2} m_{n-1} f_{n-2, n-1},\cr}
   \right\}
   \eqno {\rm (D^2.)}$$
in which $f_i$ is a function of the distance of $m_i$ from $m_n$,
and $f_{i,k}$ is a function of the distance of $m_i$ from $m_k$,
such that their derived functions or first differential coefficients,
taken with respect to the distances, express that laws of mutual
repulsion, being negative in the case of attraction; and then
we obtain, as we desired, two separate groups of equations,
for the motion of the whole system of points in space, and
for the motions of those points among themselves; namely,
first, the group
$$\left. \multieqalign{
dx_{\prime\prime}  &= x_{\prime\prime}' \, dt, &
dx_{\prime\prime}' &= 0,\cr
dy_{\prime\prime}  &= y_{\prime\prime}' \, dt, &
dy_{\prime\prime}' &= 0,\cr
dz_{\prime\prime}  &= z_{\prime\prime}' \, dt, &
dz_{\prime\prime}' &= 0,\cr}
   \right\}
   \eqno {\rm (181.)}$$
and secondly the group
$$\left. \multieqalign{
d\xi   &= \left( x_\prime'
            + {1 \over m_n} \sum\nolimits_\prime \mathbin{.} m x_\prime'
            \right) dt, &
dx_\prime' &= {1 \over m} {\delta U \over \delta \xi}   dt,\cr
d\eta  &= \left( y_\prime'
            + {1 \over m_n} \sum\nolimits_\prime \mathbin{.} m y_\prime'
            \right) dt, &
dy_\prime' &= {1 \over m} {\delta U \over \delta \eta}  dt,\cr
d\zeta &= \left( z_\prime'
            + {1 \over m_n} \sum\nolimits_\prime \mathbin{.} m z_\prime'
            \right) dt, &
dz_\prime' &= {1 \over m} {\delta U \over \delta \zeta} dt.\cr}
   \right\}
   \eqno {\rm (182.)}$$

The six differential equations of the first order (181.),
between
$x_{\prime\prime}$~$y_{\prime\prime}$~$z_{\prime\prime}$,
$x_{\prime\prime}'$~$y_{\prime\prime}'$~$z_{\prime\prime}'$
and $t$, contain the law of rectilinear and uniform motion of
the centre of gravity of the system; and the $6n - 6$ equations
of the same order, (182.), between the $6n - 6$ variables
$\xi$~$\eta$~$\zeta$ $x_\prime'$~$y_\prime'$~$z_\prime'$
and the time, are forms for the differential equations of internal
or relative motion.  We might eliminate the $3n - 3$ auxiliary
variables $x_\prime'$~$y_\prime'$~$z_\prime'$ between these
last equations, and so obtain the following other group of
$3n - 3$ equations of the second order, involving only the
relative coordinates and the time,
$$\left. \eqalign{
\xi''   &= {1 \over m} {\delta U \over \delta \xi}
            + {1 \over m_n} \sum\nolimits_\prime
               {\delta U \over \delta \xi},\cr
\eta''  &= {1 \over m} {\delta U \over \delta \eta}
            + {1 \over m_n} \sum\nolimits_\prime
               {\delta U \over \delta \eta},\cr
\zeta'' &= {1 \over m} {\delta U \over \delta \zeta}
            + {1 \over m_n} \sum\nolimits_\prime
               {\delta U \over \delta \zeta};\cr}
   \right\}
   \eqno {\rm (183.)}$$
but it is better for many purposes to retain them under the
forms (182.), omitting, however, for simplicity, the lower
accents of the auxiliary variables
$x_\prime'$~$y_\prime'$~$z_\prime'$,
because it is easy to prove that these auxiliary variables
(180.) are the components of centrobaric velocity, and because,
in investigating the properties of internal or relative motion,
we are a liberty to suppose that the centre of gravity of the
system is fixed in space, at the origin of $x$~$y$~$z$.
We may also, for simplicity, omit the lower accent of
$\sum\nolimits_\prime$, understanding that the summations
are to fall only on the first $n - 1$ masses, and denoting for
greater distinctness the $n$th mass by a separate symbol~$M$;
and then we may comprise the differential equations of relative
motion in the following simplified formula,
$$dt \, \delta H = \sum \mathbin{.} m
      ( d\xi   \, \delta x'
      - dx'    \, \delta \xi
      + d\eta  \, \delta y'
      - dy'    \, \delta \eta
      + d\zeta \, \delta z'
      - dz'    \, \delta \zeta ),
   \eqno {\rm (E^2.)}$$
in which
$$H = {\textstyle {1 \over 2}} \sum \mathbin{.}  m (x'^2 + y'^2 + z'^2)
      + {1 \over 2 M} \left\{
              \left( \sum \mathbin{.} m x' \right)^2
            + \left( \sum \mathbin{.} m y' \right)^2
            + \left( \sum \mathbin{.} m z' \right)^2
            \right\} - U,
   \eqno {\rm (F^2.)}$$

And the integrals of these equations of relative motion are
contained (by our general method) in the formula
$$\delta S = \sum \mathbin{.} m ( x' \, \delta \xi   - a' \delta \alpha
                      + y' \, \delta \eta  - b' \delta \beta
                      + z' \, \delta \zeta - c' \delta \gamma),
   \eqno {\rm (G^2.)}$$
in which $\alpha$~$\beta$~$\gamma$ $a'$~$b'$~$c'$ denote
the initial values of $\xi$~$\eta$~$\zeta$ $x'$~$y'$~$z'$,
and $S$ is the {\it principal function of relative motion\/}
of the system; that is, the former function~$S$, simplified
by the omission of the part which vanishes when the centre of
gravity is fixed, and which gives in general the laws of
motion of that centre, or the integrals of the equations (181.).

\bigbreak

{\sectiontitle
Second Example: Case of a Ternary or Multiple System with one
Predominant Mass; Equations of the undisturbed motions of the
other masses about this, in their several Binary Systems;
Differentials of all their Elements, expressed by the
coefficients of one Disturbing Function.\par}

\nobreak\bigskip

32.
Let us now suppose that the $n - 1$ masses $m$ are small in
comparision with the $n$th mass $M$; and let us separate the
expression (F${}^2$.) for $H$ into the two following parts,
$$\left. \eqalign{
H_1 &= \sum \mathbin{.} {m \over 2} \left( 1 + {m \over M} \right)
         (x'^2 + y'^2 + z'^2) - M \sum \mathbin{.} mf,\cr
H_2 &=   {m_1 m_2 \over M} (x_1' x_2' + y_1' y_2' + z_1' z_2'
              - M f_{1,2} ) + \cdots \cr
   &\mathrel{\phantom{=}} \mathord{}
       + {m_i m_k \over M} (x_i' x_k' + y_i' y_k' + z_i' z_k'
              - M f_{i,k} ) + \cdots,\cr}
   \right\}
   \eqno {\rm (H^2.)}$$
of which the latter is small in comparision with the former,
and may be neglected in a first approximation.  Suppressing
it accordingly, we are conducted to the following $6n - 6$
differential equations of the 1st order, belonging to a simpler
motion, which may be called the {\it undisturbed\/}:
$$\left. \multieqalign{
{d\xi   \over dt} &= {1 \over m} {\delta H_1 \over \delta x'}
      = \left( 1 + {m \over M} \right) x'; &
{dx' \over dt} &= - {1 \over m} {\delta H_1 \over \delta \xi}
      = M {\delta f \over \delta \xi}; \cr
{d\eta  \over dt} &= {1 \over m} {\delta H_1 \over \delta y'}
      = \left( 1 + {m \over M} \right) y'; &
{dy' \over dt} &= - {1 \over m} {\delta H_1 \over \delta \eta}
      = M {\delta f \over \delta \eta}; \cr
{d\zeta \over dt} &= {1 \over m} {\delta H_1 \over \delta z'}
      = \left( 1 + {m \over M} \right) z'; &
{dz' \over dt} &= - {1 \over m} {\delta H_1 \over \delta \zeta}
      = M {\delta f \over \delta \zeta}. \cr}
   \right\}
   \eqno {\rm (I^2.)}$$

These equations arrange themselves in $n - 1$ groups,
corresponding to the $n - 1$ binary systems $(m, M)$; and
it is easy to integrate the equations of each group separately.
We may suppose, then, these integrals found, under the forms,
$$\left. \multieqalign{
\kappa  &= \chi^{(1)} (t, \xi, \eta, \zeta, x', y', z'), &
\nu     &= \chi^{(4)} (t, \xi, \eta, \zeta, x', y', z'),\cr
\lambda &= \chi^{(2)} (t, \xi, \eta, \zeta, x', y', z'), &
\tau    &= \chi^{(5)} (t, \xi, \eta, \zeta, x', y', z'),\cr
\mu     &= \chi^{(3)} (t, \xi, \eta, \zeta, x', y', z'), &
\omega  &= \chi^{(6)} (t, \xi, \eta, \zeta, x', y', z'),\cr}
   \right\}
   \eqno {\rm (K^2.)}$$
the six quantities
$\kappa$~$\lambda$~$\mu$~$\nu$~$\tau$~$\omega$
being constant for the undisturbed motion of any one binary
system; and therefore the six functions $\chi^{(1)}$,
$\chi^{(2)}$, $\chi^{(3)}$, $\chi^{(4)}$, $\chi^{(5)}$, $\chi^{(6)}$, 
or $\kappa$, $\lambda$, $\mu$, $\nu$, $\tau$, $\omega$,
being such as to satisfy {\it identically\/} the following equation,
$$0 = m {\delta \kappa \over \delta t}
      + {\delta \kappa \over \delta \xi}   {\delta H_1 \over \delta x'}
      - {\delta \kappa \over \delta x'}    {\delta H_1 \over \delta \xi}
      + {\delta \kappa \over \delta \eta}  {\delta H_1 \over \delta y'}
      - {\delta \kappa \over \delta y'}    {\delta H_1 \over \delta \eta}
      + {\delta \kappa \over \delta \zeta} {\delta H_1 \over \delta z'}
      - {\delta \kappa \over \delta z'}    {\delta H_1 \over \delta \zeta},
   \eqno {\rm (L^2.)}$$
with five other analogous, for the five other elements
$\lambda$, $\mu$, $\nu$, $\tau$, $\omega$,
in any one binary system $(m,M)$.

\medskip

33.
Returning now to the original multiple system, we may retain
as definitions the equations (K${}^2$.), but then we can
no longer consider the elements
$\kappa_i$~$\lambda_i$~$\mu_i$~$\nu_i$~$\tau_i$~$\omega_i$
of the binary system $(m_i, M)$ as constant, because this
system is now disturbed by the other masses $m_k$; however,
the $6n - 6$ equations of disturbed relative motion, when
put under the forms
$$\left. \multieqalign{
m {d\xi   \over dt}
   &=   {\delta H_1 \over \delta x'}    + {\delta H_2 \over \delta x'},   &
m {dx'    \over dt}
   &= - {\delta H_1 \over \delta \xi}   - {\delta H_2 \over \delta \xi},  \cr
m {d\eta  \over dt}
   &=   {\delta H_1 \over \delta y'}    + {\delta H_2 \over \delta y'},   &
m {dy'    \over dt}
   &= - {\delta H_1 \over \delta \eta}  - {\delta H_2 \over \delta \eta}, \cr
m {d\zeta \over dt}
   &=   {\delta H_1 \over \delta z'}    + {\delta H_2 \over \delta z'},   &
m {dz'  \over dt}
   &= - {\delta H_1 \over \delta \zeta} - {\delta H_2 \over \delta \zeta},\cr}
   \right\}
   \eqno {\rm (M^2.)}$$
and combined with the identical equations of the kind (L${}^2$.)
give the following simple expression for the differential of the
element~$\kappa$, in its disturbed and variable state,
$$m {d\kappa \over dt}
   =    {\delta \kappa \over \delta \xi}   {\delta H_2 \over \delta x'}
      - {\delta \kappa \over \delta x'}    {\delta H_2 \over \delta \xi}
      + {\delta \kappa \over \delta \eta}  {\delta H_2 \over \delta y'}
      - {\delta \kappa \over \delta y'}    {\delta H_2 \over \delta \eta}
      + {\delta \kappa \over \delta \zeta} {\delta H_2 \over \delta z'}
      - {\delta \kappa \over \delta z'}    {\delta H_2 \over \delta \zeta},
   \eqno {\rm (N^2.)}$$
together with analogous expressions for the differentials of the
other elements.  And if we express
$\xi$~$\eta$~$\zeta$ $x'$~$y'$~$z'$,
and therefore $H_2$ itself, as depending on the time and on
these varying elements, we may transform the $6n - 6$ differential
equations of the 1st order (M${}^2$.), between
$\xi$~$\eta$~$\zeta$ $x'$~$y'$~$z'$~$t$,
into the same number of equations of the same order between the
varying elements and the time; which will be of the forms:
$$\left. \eqalign{
m {d \kappa \over dt}
   &=   \{ \kappa,  \lambda \}  {\delta H_2 \over \delta \lambda}
      + \{ \kappa,  \mu     \}  {\delta H_2 \over \delta \mu}
      + \{ \kappa,  \nu     \}  {\delta H_2 \over \delta \nu}
      + \{ \kappa,  \tau    \}  {\delta H_2 \over \delta \tau}
      + \{ \kappa,  \omega  \}  {\delta H_2 \over \delta \omega},\cr
m {d \lambda \over dt}
   &=   \{ \lambda, \kappa  \}  {\delta H_2 \over \delta \kappa}
      + \{ \lambda, \mu     \}  {\delta H_2 \over \delta \mu}
      + \{ \lambda, \nu     \}  {\delta H_2 \over \delta \nu}
      + \{ \lambda, \tau    \}  {\delta H_2 \over \delta \tau}
      + \{ \lambda, \omega  \}  {\delta H_2 \over \delta \omega},\cr
m {d \mu \over dt}
   &=   \{ \mu,     \kappa  \}  {\delta H_2 \over \delta \kappa}
      + \{ \mu,     \lambda \}  {\delta H_2 \over \delta \lambda}
      + \{ \mu,     \nu     \}  {\delta H_2 \over \delta \nu}
      + \{ \mu,     \tau    \}  {\delta H_2 \over \delta \tau}
      + \{ \mu,     \omega  \}  {\delta H_2 \over \delta \omega},\cr
m {d \nu \over dt}
   &=   \{ \nu,     \kappa  \}  {\delta H_2 \over \delta \kappa}
      + \{ \nu,     \lambda \}  {\delta H_2 \over \delta \lambda}
      + \{ \nu,     \mu     \}  {\delta H_2 \over \delta \mu}
      + \{ \nu,     \tau    \}  {\delta H_2 \over \delta \tau}
      + \{ \nu,     \omega  \}  {\delta H_2 \over \delta \omega},\cr
m {d \tau \over dt}
   &=   \{ \tau,    \kappa  \}  {\delta H_2 \over \delta \kappa}
      + \{ \tau,    \lambda \}  {\delta H_2 \over \delta \lambda}
      + \{ \tau,    \mu     \}  {\delta H_2 \over \delta \mu}
      + \{ \tau,    \nu     \}  {\delta H_2 \over \delta \nu}
      + \{ \tau,    \omega  \}  {\delta H_2 \over \delta \omega},\cr
m {d \omega \over dt}
   &=   \{ \omega,  \kappa  \}  {\delta H_2 \over \delta \kappa}
      + \{ \omega,  \lambda \}  {\delta H_2 \over \delta \lambda}
      + \{ \omega,  \mu     \}  {\delta H_2 \over \delta \mu}
      + \{ \omega,  \nu     \}  {\delta H_2 \over \delta \nu}
      + \{ \omega,  \tau    \}  {\delta H_2 \over \delta \tau},\cr}
   \right\}
   \eqno {\rm (O^2.)}$$
if we put, for abridgement,
$$\{ \kappa, \lambda \}
   =   {\delta \kappa \over \delta \xi}   {\delta \lambda \over \delta x'}
     - {\delta \kappa \over \delta x'}   {\delta \lambda \over \delta \xi}
     + {\delta \kappa \over \delta \eta}  {\delta \lambda \over \delta y'}
     - {\delta \kappa \over \delta y'}   {\delta \lambda \over \delta \eta}
     + {\delta \kappa \over \delta \zeta} {\delta \lambda \over \delta z'}
     - {\delta \kappa \over \delta z'}   {\delta \lambda \over \delta \zeta},
   \eqno {\rm (P^2.)}$$
and form the other symbols $\{ \kappa, \mu \}$,
$\{ \lambda, \kappa \}$, \&c., from this, by interchanging the
letters.  It is evident that these symbols have the properties,
$$\{ \lambda, \kappa \} = - \{ \kappa, \lambda \},\quad
  \{ \kappa, \kappa \} = 0;
   \eqno {\rm (184.)}$$
and it results from the principles of the 15th number, that
these combinations $\{ \kappa, \lambda \}$, \&c., when expressed
as functions of the elements, do not contain the time explicitly.
There are, in general, by (184.), only 15 such distinct combinations
for each of the $n - 1$ binary systems; but there would thus be,
in all, $15 n - 15$, if they admitted of no further reductions:
however it results from the principles of the 16th number, that
$12 n - 12$ of these combinations may be made to vanish
by a suitable choice of the elements.  The following is another
way of effecting as great a simplification, at least for that
extensive class of cases in which the undisturbed distance
between the two points of each binary system $(m,M)$ admits
of a minimum value.

\bigbreak

{\sectiontitle
Simplification of the Differential Expressions by a
suitable choice of the Elements.\par}

\nobreak\bigskip

34.
When the undisturbed distance~$r$ of $m$ from $M$ admists of such
a minimum~$q$, corresponding to a time~$\tau$, and satisfying
at that time the conditions
$$r' = 0,\quad r'' > 0,
   \eqno {\rm (185.)}$$
then the integrals of the group (I${}^2$.), or the known rules
of the undisturbed motion of $m$ about $M$, may be presented in
the following manner:
$$\left. \eqalign{
\kappa  &= \surd \{ (\xi y' - \eta x')^2 + (\eta z' - \zeta y')^2
                    + (\zeta x' - \xi z')^2 \};\cr
\lambda &= \kappa - \xi y' + \eta x';\cr
\mu     &= {M + m \over 2M} (x'^2 + y'^2 + z'^2) - M f(r);\cr
\nu     &= \tan^{-1} \mathbin{.}
               {\eta z' - \zeta y' \over \xi z' - \zeta x'};\cr
\tau    &= t - \int_q^r {\displaystyle
               \sqrt{ {M \over M + m} } \mathbin{.}
               {dr \over \sqrt{dr^2}} \mathbin{.} dr
            \over \displaystyle
               \sqrt{ \vphantom{\bigr(} }
               \left\{ 2 \mu + 2 M f(r) -
                  \left( 1 + {m \over M} \right)
                  {\kappa^2 \over r^2} \right\} };\cr
\omega  &= \nu + \sin^{-1}. {\kappa \zeta r^{-1}
               \over \sqrt{2 \lambda \kappa - \lambda^2}}
             - \int_q^r {\displaystyle
               \sqrt{ {M + m \over M} } \mathbin{.}
               {dr \over \sqrt{dr^2}} \mathbin{.}
                  {\kappa \over r^2} \mathbin{.} dr
            \over \displaystyle
               \sqrt{ \vphantom{\bigr(} }
               \left\{ 2 \mu + 2 M f(r) -
                  \left( 1 + {m \over M} \right)
                  {\kappa^2 \over r^2} \right\} };\cr}
   \right\}
   \eqno {\rm (Q^2.)}$$
the minimum distance~$q$ being a function of the two elements
$\kappa$, $\mu$, which must satisfy the conditions
$$2 \mu + 2 M f(q) - \left( 1 + {m \over M} \right)
                  {\kappa^2 \over q^2} = 0,\quad
  M f'(q) + \left( 1 + {m \over M} \right)
                  {\kappa^2 \over q^3} > 0;
   \eqno {\rm (186.)}$$
and $\sin^{-1} s$, $\tan^{-1} t$, being used (according to
Sir {\sc John Herschel's} notation) to express, {\it not\/} the
cosecant and cotangent, but the {\it inverse functions\/}
corresponding to sine and cosine, or the arcs which are more
commonly called
$\mathop{\rm arc} (\sin = s)$, $\mathop{\rm arc} (\tan = t)$.
It must also be observed that the factor
$\displaystyle {dr \over \sqrt{dr^2}}$, which we have
introduced under the signs of integration, is not
superfluous, but is designed to be taken as equal to
positive or negative unity, according as $dr$ is positive
or negative; that is, according as $r$ is increasing or
diminishing, so as to make the element under each integral
sign constantly positive.  In general, it appears to be
a useful rule, though not always followed by analysts,
to employ the real radical symbol $\sqrt{R}$ only for
positive quantities, unless the negative sign be expressly
prefixed; and then
$\displaystyle {r \over \sqrt{r^2}}$ will denote
positive or negative unity, according as $r$ is positive
or negative.  The arc given by its sine, in the expression
of the element $\omega$, is supposed to be so chosen as to
increase continually with the time.

\medskip

35.
After these remarks on the notation, let us apply the formula (P${}^2$.)
to calculate the values of the 15 combinations such as
$\{ \kappa, \lambda \}$, of the 6 constants or elements (Q${}^2$.).

Since
$$r = \surd \{ \xi^2 + \eta^2 + \zeta^2 \},
   \eqno {\rm (187.)}$$
it is easy to perceive that the six combinations of the
4 first elements are as follows:
$$\{ \kappa,  \lambda \} = 0,\quad
  \{ \kappa,  \mu \} = 0,\quad
  \{ \kappa,  \nu \} = 0,\quad
  \{ \lambda, \mu \} = 0,\quad
  \{ \lambda, \nu \} = 1,\quad
  \{ \mu,     \nu \} = 0.
   \eqno {\rm (188.)}$$

To form the 4 combinations of these first 4 elements with $\tau$,
we may observe, that this 5th element~$\tau$, as expressed
in (Q${}^2$.), involves explicitly (besides the time) the
distance~$r$, and the two elements $\kappa$, $\mu$; but
the combinations already determined show that these two
elements may be treated as constant in forming the
four combinations now sought; we need only attend,
therefore, to the variation of $r$, and if we interpret
by the rule (P${}^2$.) the symbols
$\{ \kappa, r \}$, $\{ \lambda, r \}$, $\{ \mu, r \}$, $\{ \nu, r \}$, 
and attend to the equations (I${}^2$.), we see that
$$\{ \kappa, r \} = 0,\quad
  \{ \lambda, r \} = 0,\quad
  \{ \mu, r \} = - {dr \over dt},\quad
  \{ \nu, r \} = 0,
   \eqno {\rm (189.)}$$
$\displaystyle {dr \over dt}$ being the total differential
coefficient of $r$ in the undisturbed motion, as determined
by the equations (I${}^2$); and, therefore, that
$$\{ \kappa, \tau \} = 0,\quad
  \{ \lambda, \tau \} = 0,\quad
  \{ \nu, \tau \} = 0,
   \eqno {\rm (190.)}$$
and
$$\{ \mu, \tau \}
   = - {\delta \tau \over \delta r}{dr \over dt}
   = + {dt \over dr}{dr \over dt} = 1:
   \eqno {\rm (191.)}$$
observing that in differentiating the expressions of the
elements (Q${}^2$), we may treat those elements as constant,
if we change the differentials of
$\xi$~$\eta$~$\zeta$ $x'$~$y'$~$z'$
to their undisturbed values.  It remains to calculate the
5 combinations of these elements with the last element~$\omega$;
which is given by (Q${}^2$.) as a function of the distance~$r$,
the coordinate $\zeta$, and the 4 elements
$\kappa$, $\lambda$, $\mu$, $\nu$; so that we may employ
this formula,
$$\{ e, \omega \}
   =   {\delta \omega \over \delta r}       \{ e, r \}
     + {\delta \omega \over \delta \zeta}   \{ e, \zeta \}
     + {\delta \omega \over \delta \kappa}  \{ e, \kappa \}
     + {\delta \omega \over \delta \lambda} \{ e, \lambda \}
     + {\delta \omega \over \delta \mu}     \{ e, \mu \}
     + {\delta \omega \over \delta \nu}     \{ e, \nu \},
   \eqno {\rm (192.)}$$
in which, if $e$ be any of the first five elements, or the
distance~$r$,
$$\{ e, r \} = - {1 \over r} \left(
        \xi   {\delta e \over \delta x'}
      + \eta  {\delta e \over \delta y'}
      + \zeta {\delta e \over \delta z'} \right),\quad
  \{ e, \zeta \} = - {\delta e \over \delta z'},\quad
  \{ e, \kappa \} = 0,
   \eqno {\rm (193.)}$$
and
$${\delta \omega \over \delta \zeta}
      = \left( {\delta \kappa \over \delta z'} \right)^{-1},\quad
  {\delta \omega \over \delta r}
      = - {d \zeta \over dr} {\delta \omega \over \delta \zeta},\quad
  {\delta \omega \over \delta \nu} = 1;
   \eqno {\rm (194.)}$$
the formula (192.) may therefore be thus written:
$$\{ e,\omega \}
   = \left\{ {\displaystyle z' \left(
           \xi   {\delta e \over \delta x'}
         + \eta  {\delta e \over \delta y'}
         + \zeta {\delta e \over \delta z'} \right) \over
         \xi x' + \eta y' + \zeta z'}
      - {\delta e \over \delta z'} \right\}
      \left( {\delta \kappa \over \delta z'} \right)^{-1}
   + \{ e, \nu \}
   + {\delta \omega \over \delta \lambda} \{ e, \lambda \}
   + {\delta \omega \over \delta \mu} \{ e, \mu \}.
   \eqno {\rm (195.)}$$
We easily find, by this formula, that
$$\{ \kappa, \omega \} = -1;\quad
  \{ \lambda, \omega \} = 0;\quad
  \{ \mu, \omega \} = 0;\quad
  \{ r, \omega \} = {dr \over dt}{\delta \omega \over \delta \mu};
   \eqno {\rm (196.)}$$
and
$$\{ \nu, \omega \}
   = - {\delta \nu \over \delta z'} {\delta \omega \over \delta \zeta}
      - {\delta \omega \over \delta \lambda} = 0.
   \eqno {\rm (197.)}$$

The formula (195.) extends to the combination $\{ \tau, \omega \}$
also; but in calculating this last combination we are to remember
that $\tau$ is given by (Q${}^2$.) as a function of
$\kappa$, $\mu$, $r$, such that
$${\delta \tau \over \delta r} = - {dt \over dr};
   \eqno {\rm (198.)}$$
and thus we see, with the help of the combinations (196.) already
determined, that
$$\{ \tau, \omega \}
   = - {\delta \tau \over \delta \kappa} - {\delta \omega \over \delta \mu}
   =   {\delta \over \delta \kappa} \int_q^r \Theta_r \,dr
     + {\delta \over \delta \mu}    \int_q^r \Omega_r \,dr,
   \eqno {\rm (199.)}$$
if we represent for abridgement by $\Theta_r$ and $\Omega_r$
the coefficients of $dr$ under the interal signs in (Q${}^2$.),
namely,
$$\left. \eqalign{
\Theta_r &= \sqrt{ {M \over M + m} } \mathbin{.} {dr \over \sqrt{dr^2}}
               \left\{ 2 \mu + 2 M f(r) -
                  {M + m \over M} \mathbin{.}
                  {\kappa^2 \over r^2} \right\}^{-{1 \over 2}},\cr
\Omega_r &= {\kappa \over r^2}
            \sqrt{ {M + m \over M} } \mathbin{.} {dr \over \sqrt{dr^2}}
               \left\{ 2 \mu + 2 M f(r) -
                  {M + m \over M} \mathbin{.}
                  {\kappa^2 \over r^2} \right\}^{-{1 \over 2}},\cr}
   \right\}
   \eqno {\rm (200.)}$$
These coefficients are evidently connected by the relation
$${\delta \Theta_r \over \delta \kappa}
  + {\delta \Omega_r \over \delta \mu} = 0,
   \eqno {\rm (201.)}$$
which gives
$${\delta \over \delta \kappa} \int_{r_\prime}^r \Theta_r \,dr
   + {\delta \over \delta \mu} \int_{r_\prime}^r \Omega_r \,dr
  = 0,
   \eqno {\rm (202.)}$$
$r_\prime$ being any quantity which does not vary with the
elements $\kappa$ and $\mu$; we might therefore at once
conclude by (199.) that the combination $\{ \tau, \omega \}$
vanishes, if a difficulty were not occasioned by the
necessity of varying the lower limit $q$, which depends on
those two elements, and by the circumstance that at this
lower limit the coefficients $\Theta_r$~$\Omega_r$ become
infinite.  However, the relation (202.) shows that we
may express the combination $\{ \tau, \omega \}$ as follows:
$$\{ \tau, \omega \}
   =   {\delta \over \delta \kappa} \int_q^{r_\prime} \Theta_r \,dr
     + {\delta \over \delta \mu}    \int_q^{r_\prime} \Omega_r \,dr,
   \eqno {\rm (203.)}$$
$r_\prime$ being an auxiliary and arbitrary quantity, which cannot
really affect the result, but may be made to facilitate the
calculation; or in other words, we may assign to the distance~$r$
any arbitrary value, not varying for infinitesimal variations
of $\kappa$, $\mu$, which may assist in calculating the value
of the expression (199.).  We may therefore suppose that the
increase of distance $r - q$ is small, and corresponds to a
small positive interval of time $t - \tau$, during which the
distance~$r$ and its differential coefficient $r'$ are
constantly increasing; and then after the first moment~$\tau$,
the quantity
$$\Theta_r = {1 \over r'}
   \eqno {\rm (204.)}$$
will be constantly finite, positive, and decreasing, during the
same interval, so that its integral must be greater than if it had
constantly its final value; that is,
$$t - \tau = \int_q^r \Theta_r \,dr > (r - q) \Theta_r.
   \eqno {\rm (205.)}$$
Hence, although $\Theta_r$ tends to infinity, yet
$(r - q) \Theta_r$ tends to zero, when by diminishing the
interval we make $r$ tend to $q$; and therefore the
following difference
$$\int_q^r \Omega_r \,dr
      - {M + m \over M} {\kappa \over q^2} \int_q^r \Theta_r \,dr
   = {M + m \over M} \int_q^r \left(
      {\kappa \over r^2} - {\kappa \over q^2}
      \right) \Theta_r \,dr
   \eqno {\rm (206.)}$$
will also tend to $0$, and so will also its partial differential
coefficient of the first order, taken with respect to $\mu$.
We find therefore the following formula for $\{ \tau, \omega \}$,
(remembering that this combination has been shown to be independent
of $r$,)
$$\{ \tau, \omega \} = \mathop{\Lambda}\limits_{r = q}
      \left\{ {\delta \over \delta \kappa} \int_q^r \Theta_r \,dr
              + {M + m \over M} {\kappa \over q^2}
              {\delta \over \delta \mu} \int_q^r \Omega_r \,dr
              \right\};
   \eqno {\rm (207.)}$$
the sign $\displaystyle \mathop{\Lambda}\limits_{r = q}$
implying that the limit is to be taken to which the expression
tends when $r$ tends to $q$.  In this last formula, as in (199.),
the integral
$\displaystyle \int_q^r \Theta_r \,dr$
may be considered as a known function of $r$, $q$, $\kappa$, $\mu$,
or simply of $r$, $q$, $\kappa$, if $\mu$ be eliminated by the
first condition (186.); and since it vanishes independently of
$\kappa$ when $r = q$, it may thus be denoted:
$$\int_q^r \Theta_r \,dr = \phi(r,q,\kappa) - \phi(q,q,\kappa),
   \eqno {\rm (208.)}$$
the form of the function~$\phi$ depending on the law of
attraction or repulsion.  This integral therefore, when
considered as depending on $\kappa$ and $\mu$, by depending
on $\kappa$ and $q$, need not be varied with respect to
$\kappa$, in calculating $\{ \tau, \omega \}$ by (207.),
because its partial differential coefficient
$\displaystyle \left(
   {\delta \over \delta \kappa} \int_q^r \Theta_r \,dr \right)$,
obtained by treating $q$ as constant, vanishes at the limit $r = q$;
nor need it be varied with respect to $q$, because, by (186.),
$${\delta q \over \delta \kappa}
   + {M + m \over M} {\kappa \over q^2} {\delta q \over \delta \mu}
      = 0;
   \eqno {\rm (209.)}$$
it may therefore be treated as constant, and we find at last
$$\{ \tau, \omega \} = 0,
   \eqno {\rm (210.)}$$
the two terms (199.) or (203.) both tending to infinity
when $r$ tends to $q$, but always destroying each other.

\medskip

36.
Collecting now our results, and presenting for greater clearness
each combination under the two forms in which it occurs when
the order of the elements is changed, we have, for each binary
system, the following thirty expressions:
$$\left. \multieqalign{
\{ \kappa,  \lambda \} &=  0, &
\{ \kappa,  \mu     \} &=  0, &
\{ \kappa,  \nu     \} &=  0, &
\{ \kappa,  \tau    \} &=  0, &
\{ \kappa,  \omega  \} &= -1, \cr
\{ \lambda, \kappa  \} &=  0, &
\{ \lambda, \mu     \} &=  0, &
\{ \lambda, \nu     \} &=  1, &
\{ \lambda, \tau    \} &=  0, &
\{ \lambda, \omega  \} &=  0, \cr
\{ \mu,     \kappa  \} &=  0, &
\{ \mu,     \lambda \} &=  0, &
\{ \mu,     \nu     \} &=  0, &
\{ \mu,     \tau    \} &=  1, &
\{ \mu,     \omega  \} &=  0, \cr
\{ \nu,     \kappa  \} &=  0, &
\{ \nu,     \lambda \} &= -1, &
\{ \nu,     \mu     \} &=  0, &
\{ \nu,     \tau    \} &=  0, &
\{ \nu,     \omega  \} &=  0, \cr
\{ \tau,    \kappa  \} &=  0, &
\{ \tau,    \lambda \} &=  0, &
\{ \tau,    \mu     \} &= -1, &
\{ \tau,    \nu     \} &=  0, &
\{ \tau,    \omega  \} &=  0, \cr
\{ \omega,  \kappa  \} &=  1, &
\{ \omega,  \lambda \} &=  0, &
\{ \omega,  \mu     \} &=  0, &
\{ \omega,  \nu     \} &=  0, &
\{ \omega,  \tau    \} &=  0; \cr}
   \right\}
   \eqno {\rm (R^2.)}$$
so that the three combinations
$$\{ \mu, \tau \},\quad \{ \omega, \kappa \},\quad \{ \lambda, \nu \}$$
are each equal to positive unity; the three inverse combinations
$$\{ \tau, \mu \},\quad \{ \kappa, \omega \},\quad \{ \nu, \lambda \}$$
are each equal to negative unity; and all the others vanish.  The six
differential equations of the first order, for the 6 varying elements
of any one binary system $(m,M)$, are therefore, by (O${}^2$.),
$$\left. \multieqalign{
m {d\mu \over dt}     &=   {\delta H_2 \over \delta \tau}, &
m {d\tau \over dt}    &= - {\delta H_2 \over \delta \mu},\cr
m {d\omega \over dt}  &=   {\delta H_2 \over \delta \kappa}, &
m {d\kappa \over dt}  &= - {\delta H_2 \over \delta \omega},\cr
m {d\lambda \over dt} &=   {\delta H_2 \over \delta \nu}, &
m {d\nu \over dt}     &= - {\delta H_2 \over \delta \lambda},\cr}
   \right\}
   \eqno {\rm (S^2.)}$$
and if we still omit the variation of $t$, they may all be
summed up in this form for the variation of $H_2$,
$$\delta H_2
   = \sum \mathbin{.} m (\mu' \delta \tau - \tau' \delta \mu
      + \omega' \delta \kappa - \kappa' \delta \omega
      + \lambda' \delta \nu - \nu' \delta \lambda),
   \eqno {\rm (T^2.)}$$
which single formula enables us to derive all the $6n - 6$
differential equations of the first order, for all the
varying elements of all the binary systems, from the variation
or from the partial differential coefficients of a single
quantity $H_2$, expressed as a function of those elements.

If we choose to introduce into the expression (T${}^2$.),
for $\delta H_2$, the variation of the time~$t$, we have only
to change $\delta \tau$ to $\delta \tau - \delta t$, because,
by (Q${}^2$.), $\delta t$ enters only so accompanied; that is,
$t$ enters only under the form $t - \tau_i$, in the expressions
of $\xi_i$~$\eta_i$~$\zeta_i$ $x'_i$~$y'_i$~$z'_i$
as functions of the time and of the elements; we have,
therefore
$${\delta H_2 \over \delta t}
   = - \sum {\delta H_2 \over \delta \tau}
   = - \sum \mathbin{.} m \mu';
   \eqno {\rm (211.)}$$
and since, by (H${}^2$.), (Q${}^2$),
$$H_1 = \sum \mathbin{.} m \mu,
   \eqno {\rm (212.)}$$
we find, finally,
$${d H_1 \over dt} = - {\delta H_2 \over \delta t}.
   \eqno {\rm (U^2.)}$$

This remarkable form for the differential of $H_1$, considered
as a varying element, is general for all problems of dynamics.
It may be deduced by the general method from the formul{\ae}
of the 13th and 14th numbers, which give
$$\left. \eqalign{
{d H_1 \over dt}
   &=   {\delta H_2 \over \delta \kappa_1}    \sum \left(
            {\delta H_1 \over \delta \eta}
            {\delta \kappa_1 \over \delta \varpi}
          - {\delta H_1 \over \delta \varpi}
            {\delta \kappa_1 \over \delta \eta} \right)
      + \cdots
      + {\delta H_2 \over \delta \kappa_{6n}} \sum \left(
            {\delta H_1 \over \delta \eta}
            {\delta \kappa_{6n} \over \delta \varpi}
          - {\delta H_1 \over \delta \varpi}
            {\delta \kappa_{6n} \over \delta \eta} \right) \cr
   &=    {\delta H_2 \over \delta \kappa_1}
            {\delta \kappa_1 \over \delta t}
       + {\delta H_2 \over \delta \kappa_2}
            {\delta \kappa_2 \over \delta t}
       + \cdots
       + {\delta H_2 \over \delta \kappa_{6n}}
            {\delta \kappa_{6n} \over \delta t}
    =  - {\delta H_2 \over \delta t},\cr}
   \right\}
   \eqno {\rm (213.)}$$
$\kappa_1 \, \kappa_2 \,\ldots\, \kappa_{6n}$ being any $6n$
elements of a system expressed as functions of the time and of
the quantities $\eta$~$\omega$; or concisely by this special
consideration, that $H_1 + H_2$ is constant in the disturbed
motion, and that in taking the first total differential
coefficient of $H_2$ with respect to the time, the elements
may by (F${}^1$.) be treated as constant.  It is also a
remarkable corollary of the general principles just referred
to, but not one difficult to verify, that the first
partial differential coefficient
$\displaystyle {\delta \kappa_s \over \delta t}$, of any
element $\kappa_s$, taken with respect to the time, may
be expressed as a function of the elements alone, not
involving the time explicitly.

\bigbreak

{\sectiontitle
On the essential distinction between the Systems of Varying Elements
considered in this Essay and those hitherto employed by mathematicians.\par}

\nobreak\bigskip

37.
When we shall have integrated the differential equations of
varying elements (S${}^2$.), we can then calculate the varying
relative coordinates $\xi$~$\eta$~$\zeta$, for any binary
system $(m,M)$, by the rules of undisturbed motion, as expressed
by the equations (I${}^2$.), (Q${}^2$.), or by the following
connected formul{\ae}:
$$\left. \eqalign{
\xi   &= r \left( \cos \theta
                  + {\lambda \over \kappa} \sin (\theta - \nu) \sin \nu
                  \right),\cr
\eta  &= r \left( \sin \theta
                  - {\lambda \over \kappa} \sin (\theta - \nu) \cos \nu
                  \right),\cr
\zeta &= {r \over \kappa}
            \sqrt{2\lambda \kappa - \lambda^2}
            \sin (\theta - \nu);\cr}
   \right\}
   \eqno {\rm (V^2.)}$$
in which the distance~$r$ is determined as a function of the
time~$t$ and of the elements $\tau$, $\kappa$, $\mu$, by
the 5th equation (Q${}^2$.), and in which
$$\theta = \omega + \int_q^r {\displaystyle
               \sqrt{ {M + m \over M} } \mathbin{.}
               {dr \over \sqrt{dr^2}} \mathbin{.}
                  {\kappa \over r^2} dr
            \over \displaystyle
               \sqrt{ \vphantom{\bigr(} }
               \left\{ 2 \mu + 2 M f(r) -
                  \left( 1 + {m \over M} \right)
                  {\kappa^2 \over r^2} \right\} },
   \eqno {\rm (W^2.)}$$
$q$ being still the minimum of $r$, when the orbit is treated
as constant, and being still connected with the elements
$\kappa$, $\mu$, by the first equation of condition (186.).
In astronomical language, $M$ is the sun, $m$ a planet,
$\xi$~$\eta$~$\zeta$ are the heliocentric rectangular
coordinates, $r$ is the radius vector, $\theta$ the longitude
in the orbit, $\omega$ the longitude of the perihelion,
$\nu$ of the node, $\theta - \omega$ is the true anomaly,
$\theta - \nu$ the argument of latitude, $\mu$ the constant
part of the half square of undisturbed heliocentric velocity,
diminished in the ratio of the sun's mass ($M$) to the
sum ($M + m$) of masses of sun and planet, $\kappa$ is the double
of the areal velocity diminished in the same ratio,
$\displaystyle {\lambda \over \kappa}$
is the versed sine of the inclination of the orbit, $q$ the
perihelion distance, and $\tau$ the time of perihelion
passage.  The law of attraction or repulsion is here left
undetermined; for {\sc Newton's} law, $\mu$ is the sun's mass
divided by the axis major of the orbit taken negatively,
and $\kappa$ is the square root of the semiparameter,
multiplied by the sun's mass, and divided by the square
root of the sum of the masses of sun and planet.
But the varying ellipse or other orbit, which the foregoing
formul{\ae} require, differs essentially (though little)
from that hitherto employed by astronomers: because it
gives correctly the heliocentric coordinates, but {\it not\/}
the heliocentric components of velocity, without differentiating
the elements in the calculation; and therefore does {\it not
touch\/} but {\it cuts\/}, (though under a very small angle,)
the actual heliocentric orbit, described under the influence
of all the disturbing forces.

\medskip

38.
For it results from the foregoing theory, that if we differentiate
the expressions (V${}^2$.) for the heliocentric coordinates,
without differentiating the elements, and then assign to those
new varying elements their values as functions of the time,
obtained from the equations (S${}^2$), and deduce the centrobaric
components of velocity by the formul{\ae} (I${}^2$.), or by
the following:
$$x' = {M \xi'   \over M + m},\quad
  y' = {M \eta'  \over M + m},\quad
  z' = {M \zeta' \over M + m};
   \eqno {\rm (214.)}$$
then these centrobaric components will be the same functions
of the time and of the new varying elements which might
be otherwise deduced by elimination from the integrals (Q${}^2$.),
and will represent rigorously (by the extension given in the
theory to those last-mentioned integrals) the components of
velocity of the disturbed planet~$m$, relatively to the centre
of gravity of the whole solar system.  We chose, as more
suitable to the general course of our method, that these
centrobaric components of velocity should be the auxiliary
variables to be combined with the heliocentric coordinates,
and to have their disturbed values rigorously expressed
by the formul{\ae} of undisturbed motion; but in making
this choice it became necessary to modify these latter
formul{\ae}, and to determine a varying orbit essentially
distinct in theory (though little differing in practice)
from that conceived so beautifully by {\sc Lagrange}.  The orbit
which he imagined was more simply connected with the
heliocentric motion of a {\it single planet}, since it
gave, for such heliocentric motion, the velocity as
well as the position; the orbit which we have chosen
is perhaps more closely combined with the conception of a
{\it multiple system}, moving about its common centre of
gravity, and influenced in every part by the actions of
all the rest.  Whichever orbit shall be hereafter adopted
by astronomers, they will remember that both are equally
fit to represent the celestial appearances, if the numeric
elements of either set be suitable determined by
observation, and the elements of the other set of orbits
be deduced from these by calculation.  Meantime mathematicians
will judge, whether in sacrificing a part of the simplicity
of that geometrical conception on which the theories of
{\sc Lagrange} and {\sc Poisson} are founded, a simplicity of
another kind has not been introduced, which was wanting in
those admirable theories; by our having succeeded in expressing
rigorously the differentials of {\it all\/} our own new
varying elements through the coefficients of a {\it single\/}
function: whereas it has seemed necessary hitherto to
employ one function for the Earth disturbed by Venus, and
another function for Venus disturbed by the Earth.

\bigbreak

{\sectiontitle
Integration of the Simplified Equations, which determine the
new varying Elements.\par}

\nobreak\bigskip

39.
The simplified differential equations of varying elements,
(S${}^2$.), are of the same form as the equations (A.),
and may be integrated in a similar manner.  If we put,
for abridgement,
$$(\tau, \kappa, \nu)
   = \int_0^t \left\{ \sum \left(
           \tau {\delta H_2 \over \delta \tau}
         + \kappa {\delta H_2 \over \delta \kappa}
         + \nu {\delta H_2 \over \delta \nu}
         \right) - H_2 \right\} dt,
   \eqno {\rm (X^2.)}$$
and interpret similarly the symbols $(\mu, \omega, \lambda)$, \&c.,
we can easily assign the variations of the following 8 combinations,
$(\tau, \kappa, \nu)$, $(\mu, \omega, \lambda)$, $(\mu, \kappa, \nu)$,
$(\tau, \omega, \lambda)$, $(\tau, \omega, \nu)$, $(\mu, \kappa, \lambda)$,
$(\tau, \kappa, \lambda)$, $(\mu, \omega, \nu)$; namely
$$\left. \eqalign{
\delta (\tau, \kappa, \nu) &= \sum \mathbin{.} m
         ( \tau \, \delta \mu - \tau_0 \, \delta \mu_0
         + \kappa \, \delta \omega - \kappa_0 \, \delta \omega_0
         + \nu \, \delta \lambda - \nu_0 \, \delta \lambda_0 )
         - H_2 \, \delta t,\cr
\delta (\mu, \omega, \lambda) &= \sum \mathbin{.} m
         ( \mu_0 \, \delta \tau_0 - \mu \, \delta \tau
         + \omega_0 \, \delta \kappa_0 - \omega \, \delta \kappa
         + \lambda_0 \, \delta \nu_0 - \lambda \, \delta \nu )
         - H_2 \, \delta t,\cr
\delta (\mu, \kappa, \nu) &= \sum \mathbin{.} m
         ( \mu_0 \, \delta \tau_0 - \mu \, \delta \tau
         + \kappa \, \delta \omega - \kappa_0 \, \delta \omega_0
         + \nu \, \delta \lambda - \nu_0 \, \delta \lambda_0 )
         - H_2 \, \delta t,\cr
\delta (\tau, \omega, \lambda) &= \sum \mathbin{.} m
         ( \tau \, \delta \mu - \tau_0 \, \delta \mu_0
         + \omega_0 \, \delta \kappa_0 - \omega \, \delta \kappa
         + \lambda_0 \, \delta \nu_0 - \lambda \, \delta \nu )
         - H_2 \, \delta t,\cr
\delta (\tau, \omega, \nu) &= \sum \mathbin{.} m
         ( \tau \, \delta \mu - \tau_0 \, \delta \mu_0
         + \omega_0 \, \delta \kappa_0 - \omega \, \delta \kappa
         + \nu \, \delta \lambda - \nu_0 \, \delta \lambda_0 )
         - H_2 \, \delta t,\cr
\delta (\mu, \kappa, \lambda) &= \sum \mathbin{.} m
         ( \mu_0 \, \delta \tau_0 - \mu \, \delta \tau
         + \kappa \, \delta \omega - \kappa_0 \, \delta \omega_0
         + \lambda_0 \, \delta \nu_0 - \lambda \, \delta \nu )
         - H_2 \, \delta t,\cr
\delta (\tau, \kappa, \lambda) &= \sum \mathbin{.} m
         ( \tau \, \delta \mu - \tau_0 \, \delta \mu_0
         + \kappa \, \delta \omega - \kappa_0 \, \delta \omega_0
         + \lambda_0 \, \delta \nu_0 - \lambda \, \delta \nu )
         - H_2 \, \delta t,\cr
\delta (\mu, \omega, \nu) &= \sum \mathbin{.} m
         ( \mu_0 \, \delta \tau_0 - \mu \, \delta \tau
         + \omega_0 \, \delta \kappa_0 - \omega \, \delta \kappa
         + \nu \, \delta \lambda - \nu_0 \, \delta \lambda_0 )
         - H_2 \, \delta t,\cr}
   \right\}
   \eqno {\rm (Y^2.)}$$
$\kappa_0$~$\lambda_0$~$\mu_0$~$\nu_0$~$\tau_0$~$\omega_0$
being the initial values of the varying elements
$\kappa$~$\lambda$~$\mu$~$\nu$~$\tau$~$\omega$.
If, then, we consider, for example, the first of these 8
combinations $(\tau, \kappa, \nu)$, as a function of all of
the $3n - 3$ elements $\mu_i$~$\omega_i$~$\lambda_i$,
and of their initial values
$\mu_{0,i}$~$\omega_{0,i}$~$\lambda_{0,i}$,
involving also in general the time explicitly, we shall have
the following forms for the $6n - 6$ rigorous integrals
of the $6n - 6$ equations (S${}^2$.):
$$\left. \multieqalign{
m_i \tau_i
   &=   {\delta \over \delta \mu_i} (\tau, \kappa, \nu); &
m_i \tau_{0,i}
   &= - {\delta \over \delta \mu_{0,i}} (\tau, \kappa, \nu); \cr
m_i \kappa_i
   &=   {\delta \over \delta \omega_i} (\tau, \kappa, \nu); &
m_i \kappa_{0,i}
   &= - {\delta \over \delta \omega_{0,i}} (\tau, \kappa, \nu); \cr
m_i \nu_i
   &=   {\delta \over \delta \lambda_i} (\tau, \kappa, \nu); &
m_i \nu_{0,i}
   &= - {\delta \over \delta \lambda_{0,i}} (\tau, \kappa, \nu); \cr}
   \right\}
   \eqno {\rm (Z^2.)}$$
and in like manner we can deduce forms for the same rigorous
integrals, from any one of the eight combinations (Y${}^2$.).
The determination of all the varying elements would therefore
be fully accomplished, if we could find the complete expression
for any one of these 8 combinations.

\bigbreak

40.
A first approximate expression for any one of them can be found
from the form under which we have supposed $H_2$ to be put,
namely, as a function of the elements and of the time,
which may be thus denoted:
$$H_2 = H_2(t,
   \kappa_1, \lambda_1, \mu_1, \nu_1, \tau_1, \omega_1,\cdots\,
   \kappa_{n-1}, \lambda_{n-1}, \mu_{n-1}, \nu_{n-1}, \tau_{n-1},
   \omega_{n-1} );
   \eqno {\rm (A^3.)}$$
by changing in this function the varying elements to their
initial values, and employing the following approximate integrals
of the equations (S${}^2$.),
$$\left. \multieqalign{
\mu &= \mu_0 + {1 \over m} \int_0^t
         {\delta H_2 \over \delta \tau_0} \, dt, &
\tau &= \tau_0 - {1 \over m} \int_0^t
         {\delta H_2 \over \delta \mu_0} \, dt, \cr
\omega &= \omega_0 + {1 \over m} \int_0^t
         {\delta H_2 \over \delta \kappa_0} \, dt, &
\kappa &= \kappa_0 - {1 \over m} \int_0^t
         {\delta H_2 \over \delta \omega_0} \, dt, \cr
\lambda &= \lambda_0 + {1 \over m} \int_0^t
         {\delta H_2 \over \delta \nu_0} \, dt, &
\nu &= \nu_0 - {1 \over m} \int_0^t
         {\delta H_2 \over \delta \lambda_0} \, dt. \cr}
   \right\}
   \eqno {\rm (B^3.)}$$
For if we denote, for example, the first of the 8 combinations
(Y${}^2$.) by $G$, so that
$$G = (\tau, \kappa, \nu),
   \eqno {\rm (C^3.)}$$
we shall have, as a first approximate value,
$$G_1 = \int_0^t \left\{ \sum \left(
           \tau_0 {\delta H_2 \over \delta \tau_0}
         + \kappa_0 {\delta H_2 \over \delta \kappa_0}
         + \nu_0 {\delta H_2 \over \delta \nu_0}
         \right) - H_2 \right\} dt;
   \eqno {\rm (D^3.)}$$
and after thus expressing $G_1$ as a function of the time,
and of the initial elements, we can eliminate the initial
quantities of the forms $\tau_0$~$\kappa_0$~$\nu_0$,
and introduce in their stead the final quantities
$\mu$~$\omega$~$\lambda$, so as to obtain an
expression for $G_1$ of the kind supposed in (Z${}^2$.),
namely, a function of the time~$t$, the varying elements
$\mu$~$\omega$~$\lambda$, and their initial values
$\mu_0$~$\omega_0$~$\lambda_0$.  An approximate
expression thus found may be corrected by a process of
that kind, which has often been employed in this Essay
for other similar purposes.  For the function $G$, or the
combination $(\tau, \kappa, \nu)$, must satisfy rigorously,
by (Y${}^2$.), (A${}^3$.), the following partial differential
equation:
$$0 = {\delta G \over \delta t} + H_2 \left( t,
         {1 \over m_1} {\delta G \over \delta \omega_1},
         \lambda_1, \mu_1,
         {1 \over m_1} {\delta G \over \delta \lambda_1},
         {1 \over m_1} {\delta G \over \delta \mu_1},
         \omega_1,
         {1 \over m_2} {\delta G \over \delta \omega_2},
         \cdots \, \omega_{n-1} \right);
   \eqno {\rm (E^3.)}$$
and each of the other analogous functions or combinations (Y${}^2$.)
must satisfy an analogous equation: if then we change $G$ to
$G_1 + G_2$, and neglect the squares and products of the
coefficients of the small correction $G_2$, $G_1$ being a
first approximation such as that already found, we are conducted,
as a second approximation on principles already explained, to
the following expression for this correction $G_2$:
$$G_2 = - \int_0^t \left\{ {\delta G_1 \over \delta t} + H_2 \left(
         {1 \over m_1} {\delta G_1 \over \delta \omega_1},
         \lambda_1, \mu_1,
         {1 \over m_1} {\delta G_1 \over \delta \lambda_1},
         {1 \over m_1} {\delta G_1 \over \delta \mu_1},
         \omega_1,\cdots \,
         \right) \right\} dt:
   \eqno {\rm (F^3.)}$$
which may be continually and indefinitely improved by a repetition
of the same process of correction.  We may therefore, theoretically,
consider the problem as solved; but it must remain for future
consideration, and perhaps for actual trial, to determine which
of all these various processes of successive and indefinite
approximation, deduced in the present Essay and in the former,
as corollaries of one general Method, and as consequences of one
central Idea, is best adapted for numeric application, and for the
mathematical study of phenomena.

\bye
